% !TEX root = ../Elements.tex
%%%%%%
% BOOK 13
%%%%%%
\pdfbookmark[0]{Book 13}{book13}
\pagestyle{empty}
\begin{center}
{\Huge ELEMENTS BOOK 13}\\
\spa\spa\spa
{\huge\it The Platonic Solids\symbolfootnote[2]{The five regular solids---the cube, tetrahedron ({\rm i.e.}, pyramid), octahedron, icosahedron, and dodecahedron---were problably discovered by the school of Pythagoras. They are generally termed ``Platonic'' solids because they feature prominently in Plato's famous dialogue {\em Timaeus}. Many of the theorems contained in this book---particularly those which pertain to the last two solids---are ascribed to Theaetetus of Athens.}}
\end{center}
\vspace{2cm}
\epsfysize=5in
\centerline{\epsffile{Book13/front.eps}}
\newpage

%%%%
%13.1
%%%%
\pdfbookmark[1]{Proposition 13.1}{pdf13.1}
\pagestyle{fancy}
\cfoot{{\greekfont τ̀ηεπαγε}}
\lhead{\large{\greekfont ΣΤΟΙΘΕΙΩΝ γ̀γν{13}.}}
\rhead{\large ELEMENTS BOOK 13}
\begin{Parallel}{}{}
\ParallelLText{
\begin{center}
{\large \ggn{1}.}
\end{center}\vspace*{-7pt}

{\greekfont Ἐὰν εὐθεῖα γραμμὴ ἄκρον καὶ μέσον λόγον τμ ηθῇ, τὸ
μεῖζον τμῆμα προσλαβὸν τὴν ἡμίσειαν τῆς ὅλης πενταπλάσιον
δύναται τοῦ ἀπὸ τῆς ἡμισείας τετραγώνου.}

\epsfysize=2.5in
\centerline{\epsffile{Book13/fig01g.eps}}

{\greekfont Εὐθεῖα γὰρ γραμμὴ ἡ ΑΒ ἄκρον καὶ μέσον λόγον τετμήσθω
κατὰ τὸ Γ σημεῖον, καὶ ἔστω μεῖζον τμῆμα τὸ ΑΓ, καὶ
ἐκβεβλήσθω ἐπ εὐθείας τῇ ΓΑ εὐθεῖα ἡ ΑΔ, καὶ κείσθω
τῆς ΑΒ ἡμίσεια ἡ ΑΔ; λέγω, ὅτι πενταπλάσιόν ἐστι τὸ ἀπὸ
τῆς ΓΔ τοῦ ἀπὸ τῆς ΔΑ.}

{\greekfont Ἀναγεγράφθωσαν γὰρ ἀπὸ τῶν ΑΒ, ΔΓ τετράγωνα τὰ ΑΕ,
ΔΖ, καὶ καταγεγράφθω ἐν τῷ ΔΖ τὸ σχῆμα, καὶ διήθθω ἡ ΖΓ
ἐπὶ τὸ Η. καὶ ἐπεὶ ἡ ΑΒ ἄκρον καὶ μέσον λόγον τέτμ ηται
κατὰ τὸ Γ, τὸ ἄρα ὑπὸ τῶν ΑΒΓ ἴσον ἐστὶ τῷ ἀπὸ τῆς
ΑΓ. καί ἐστι τὸ μὲν ὑπὸ τῶν ΑΒΓ τὸ ΓΕ, τὸ δὲ ἀπὸ τῆς
ΑΓ τὸ ΖΘ; ἴσον ἄρα τὸ ΓΕ τῷ ΖΘ. καὶ ἐπεὶ διπλῆ ἐστιν
ἡ ΒΑ τῆς ΑΔ, ἴση δὲ ἡ μὲν ΒΑ τῇ ΚΑ, ἡ δὲ ΑΔ τῇ
ΑΘ, διπλῆ ἄρα καὶ ἡ ΚΑ τῆς ΑΘ. ὡς δὲ ἡ ΚΑ πρὸς τὴν
ΑΘ, οὕτως τὸ ΓΚ πρὸς τὸ ΓΘ;
διπλάσιον ἄρα τὸ ΓΚ τοῦ ΓΘ. εἰσὶ δὲ καὶ τὰ ΛΘ, ΘΓ διπλάσια
τοῦ ΓΘ. ἴσον ἄρα τὸ ΚΓ τοῖς ΛΘ, ΘΓ. ἐδείθθη δὲ καὶ τὸ ΓΕ
τῷ ΘΖ ἴσον; ὅλον ἄρα τὸ ΑΕ τετράγωνον ἴσον ἐστὶ τῷ
ΜΝΧ γνώμονι. καὶ ἐπεὶ διπλῆ ἐστιν ἡ ΒΑ τῆς ΑΔ,
τετραπλάσιόν ἐστι τὸ ἀπὸ τῆς ΒΑ τοῦ ἀπὸ τῆς ΑΔ, τουτέστι
τὸ ΑΕ τοῦ ΔΘ. ἴσον δὲ τὸ ΑΕ τῷ ΜΝΧ γνώμονι; καὶ ὁ
ΜΝΧ ἄρα γνώμων τετραπλάσιός ἐστι τοῦ ΑΟ; ὅλον ἄρα
τὸ ΔΖ πενταπλάσιόν ἐστι τοῦ ΑΟ. καί ἐστι τὸ μὲν ΔΖ τὸ
ἀπὸ τῆς ΔΓ, τὸ δὲ ΑΟ τὸ ἀπὸ τῆς ΔΑ; τὸ ἄρα ἀπὸ
τῆς ΓΔ πενταπλάσιόν ἐστι τοῦ ἀπὸ τῆς ΔΑ.}

{\greekfont Ἐὰν ἄρα εὐθεῖα ἄκρον καὶ μέσον λόγον τμ ηθῇ,
τὸ μεῖζον τμῆμα προσλαβὸν τὴν ἡμίσειαν τῆς ὅλης
πενταπλάσιον δύναται τοῦ ἀπὸ τῆς ἡμισείας τετραγώνου;
ὅπερ ἔδει δεῖχαι.}}

\ParallelRText{
\begin{center}
{\large Proposition 1}
\end{center}

If a straight-line is cut in extreme and mean ratio then the square on the
greater piece, added to half of the whole, is  five times the square
on the half.

\epsfysize=2.5in
\centerline{\epsffile{Book13/fig01e.eps}}

For let the straight-line $AB$ be cut in extreme and mean ratio at
point $C$, and let $AC$ be the greater piece. And let the straight-line
$AD$ be produced in a straight-line with $CA$. And let
$AD$ be made (equal to) half of $AB$. I say that the (square) on
$CD$ is five times the (square) on $DA$. 

For let the squares $AE$ and $DF$ be described on $AB$ and
$DC$ (respectively). And let the figure in $DF$ be drawn.  
And let $FC$ be drawn across to $G$. And since $AB$ has been
cut in extreme and mean ratio at $C$, the (rectangle
contained) by $ABC$ is thus equal to the (square) on $AC$ [Def.~6.3, Prop.~6.17].
And  $CE$ is the (rectangle contained) by $ABC$, and $FH$ the
(square) on $AC$. Thus, $CE$ (is) equal to $FH$. 
And since $BA$ is double $AD$, and $BA$ (is) equal to $KA$,
and $AD$ to $AH$, $KA$ (is) thus also double $AH$. 
And as $KA$ (is) to $AH$, so $CK$ (is) to $CH$ [Prop.~6.1].
Thus, $CK$ (is) double $CH$. And $LH$ plus $HC$ is also double $CH$
[Prop.~1.43].  Thus, $KC$ (is) equal to $LH$ plus $HC$. And
$CE$ was also shown (to be) equal to $HF$. Thus, the whole square
$AE$ is equal to the gnomon $MNO$. And since $BA$ is double
$AD$, the (square) on $BA$ is four times the (square) on $AD$---that
is to say, $AE$ (is four times) $DH$. And $AE$ (is) equal to
gnomon $MNO$. And, thus, gnomon $MNO$ is also four times $AP$. 
Thus, the whole of $DF$ is five times $AP$. And $DF$ is the (square)
on $DC$, and $AP$ the (square) on $DA$. Thus, the
(square) on $CD$ is five times the (square) on $DA$. 

Thus, if a straight-line is cut in extreme and mean ratio then the square on the
greater piece, added to half of the whole, is five times the square
on the half. (Which is) the very thing it was required to show.\\}
\end{Parallel}

%%%%
%13.2
%%%%
\pdfbookmark[1]{Proposition 13.2}{pdf13.2}
\begin{Parallel}{}{}
\ParallelLText{
\begin{center}
{\large \ggn{2}.}
\end{center}\vspace*{-7pt}

{\greekfont Ἐὰν εὐθεῖα γραμμὴ τμήματος ἑαυτῆς πενταπλάσιον δύνηται,
τῆς διπλασίας τοῦ εἰρημένου τμήματος ἄκρον καὶ μέσον
λόγον τεμνομένης τὸ μεῖζον τμῆμα τὸ λοιπὸν μέρος ἐστὶ
τῆς ἐχ ἀρθῆς εὐθείας.}

\epsfysize=2.5in
\centerline{\epsffile{Book13/fig02g.eps}}

{\greekfont Εὐθεῖα γὰρ γραμμὴ ἡ ΑΒ τμήματος ἑαυτῆς τοῦ ΑΓ
πενταπλάσιον δυνάσθω, τῆς δὲ ΑΓ διπλῆ ἔστω ἡ ΓΔ. λέγω,
ὅτι τῆς ΓΔ ἄκρον καὶ μέσον λόγον τεμνομένος τὸ μεῖζον
τμῆμά ἐστιν ἡ ΓΒ.}

{\greekfont Ἀναγεγράφθω γὰρ ἀφ ἑκατέρας τῶν ΑΒ, ΓΔ τετράγωνα
τὰ ΑΖ, ΓΗ, καὶ καταγεγράφθω ἐν τῷ ΑΖ τὸ σχῆμα, καὶ διήθθω
ἡ ΒΕ. καὶ ἐπεὶ πενταπλάσιόν ἐστι τὸ ἀπό τῆς ΒΑ τοῦ ἀπὸ
τῆς ΑΓ, πενταπλάσιόν ἐστι τὸ ΑΖ τοῦ ΑΘ. τετραπλάσιος
ἄρα ὁ ΜΝΧ γνώμων τοῦ ΑΘ. καὶ ἐπεὶ διπλῆ ἐστιν ἡ ΔΓ
τῆς ΓΑ, τετραπλάσιον ἄρα ἐστὶ τὸ ἀπὸ ΔΓ τοῦ ἀπὸ ΓΑ,
τουτέστι τὸ ΓΗ τοῦ ΑΘ. ἐδείθθη δὲ καὶ ὁ ΜΝΧ
γνώμων τετραπλάσιος τοῦ ΑΘ; ἴσος ἄρα ὁ ΜΝΧ γνώμων
τῷ ΓΗ. καὶ ἐπεὶ διπλῆ ἐστιν ἡ ΔΓ τῆς ΓΑ, ἴση δὲ ἡ
μὲν ΔΓ τῇ ΓΚ, ἡ δὲ ΑΓ τῇ ΓΘ, [διπλῆ ἄρα καὶ ἡ ΚΓ τῆς ΓΘ], διπλάσιον ἄρα καὶ τὸ ΚΒ τοῦ ΒΘ. εἰσὶ δὲ καὶ
τὰ ΛΘ, ΘΒ τοῦ ΘΒ διπλάσια; ἴσον ἄρα τὸ ΚΒ τοῖς ΛΘ, ΘΒ.
ἐδείθθη δὲ καὶ ὅλος ὁ ΜΝΧ γνώμων ὅλῳ τῷ ΓΗ ἴσος;
καὶ λοιπὸν ἄρα τὸ ΘΖ τῷ ΒΗ ἐστιν ἴσον. καί ἐστι τὸ μὲν
ΒΗ τὸ ὑπὸ τῶν ΓΔΒ; ἴση γὰρ ἡ ΓΔ τῇ ΔΗ; τὸ δὲ ΘΖ
τὸ ἀπὸ τῆς ΓΒ; τὸ ἄρα ὑπὸ τῶν ΓΔΒ ἴσον ἐστὶ τῷ
ἀπὸ τῆς ΓΒ. ἔστιν ἄρα ὡς ἡ ΔΓ πρὸς τὴν ΓΒ, οὕτως
ἡ ΓΒ πρὸς τὴν ΒΔ. μείζων δὲ ἡ ΔΓ τῆς ΓΒ; μείζων ἄρα καὶ
ἡ ΓΒ τῆς ΒΔ. τῆς ΓΔ ἄρα εὐθείας ἄκρον καὶ μέσον λόγον
τεμνομένης τὸ μεῖζον τμῆμά ἐστιν ἡ ΓΒ.}

{\greekfont Ἐὰν ἄρα εὐθεῖα γραμμὴ τμήματος ἑαυτῆς πενταπλάσιον
δύνηται, τῆς διπλασίας τοῦ εἰρημένου τμήμ-ατος ἄκρον καὶ
μέσον λόγον τεμνομένης τὸ μεῖζον τμῆμα τὸ λοιπὸν μέρος
ἐστὶ τῆς ἐχ ἀρθῆς εὐθείας; ὅπερ ἔδει δεῖχαι.}}

\ParallelRText{
\begin{center}
{\large Proposition 2}
\end{center}

If the square on a straight-line is five times the (square) on a piece of
it, and  double the aforementioned piece is cut in extreme and
mean ratio, then the greater piece  is  the remaining part of the
original straight-line. 

\epsfysize=2.5in
\centerline{\epsffile{Book13/fig02e.eps}}

For let the square on the straight-line $AB$ be five times the (square)
on the piece of it, $AC$. And let $CD$ be double $AC$. I say that
if $CD$ is cut in extreme and mean ratio then the greater piece
is $CB$. 

For let the squares $AF$ and $CG$ be described on each of
$AB$ and $CD$ (respectively). And let the figure  in $AF$ be drawn.
And let $BE$ be drawn across. And since the (square) on $BA$
is five times the (square) on $AC$, $AF$ is five times $AH$. 
Thus, gnomon $MNO$ (is) four times $AH$. 
And since
$DC$ is double $CA$, the (square) on $DC$ is thus four times the (square)
on $CA$---that is to say, $CG$ (is four times) $AH$. And the gnomon
$MNO$ was also shown (to be) four times $AH$. Thus, gnomon $MNO$
(is) equal to $CG$. And since $DC$ is double $CA$, and $DC$
(is) equal to $CK$, and $AC$ to $CH$, [$KC$ (is) thus also double $CH$], (and) $KB$ (is) also double $BH$
[Prop.~6.1]. And $LH$ plus $HB$ is also double $HB$ [Prop.~1.43].
Thus, $KB$ (is) equal to $LH$ plus $HB$. And the whole gnomon
$MNO$ was also shown (to be) equal to the whole of $CG$. Thus, the
remainder $HF$ is also equal to (the remainder) $BG$.  And $BG$
is the (rectangle contained) by $CDB$. For  $CD$ (is) equal to $DG$.
And $HF$ (is) the square on $CB$. Thus, the (rectangle contained)
by $CDB$ is equal to the (square) on $CB$. Thus, as $DC$ is to $CB$,
so $CB$ (is) to $BD$ [Prop.~6.17]. And $DC$ (is) greater than $CB$ (see lemma). Thus, $CB$ (is) also greater than $BD$ [Prop.~5.14]. Thus, if the
straight-line $CD$ is cut in extreme and mean ratio then the greater piece is $CB$.

Thus, if the square on a straight-line is five times the (square) on a piece of
itself, and  double the aforementioned piece is cut in extreme and
mean ratio, then the greater piece is the remaining part of the
original straight-line.  (Which is) the very thing it was required to show.}
\end{Parallel}

\begin{Parallel}{}{}
\ParallelLText{
\begin{center}
{\large {\greekfont Λῆμμα}.}
\end{center}\vspace*{-7pt}

{\greekfont Ὅτι δὲ ἡ διπλῆ τῆς ΑΓ μείζων ἐστὶ τῆς ΒΓ, οὕτως
δεικτέον.}

{\greekfont Εἰ γὰρ μή, ἔστω, εἰ δυνατόν, ἡ ΒΓ διπλῆ τῆς ΓΑ.
τετραπλάσιον ἄρα τὸ ἀπὸ τῆς ΒΓ τοῦ ἀπὸ τῆς ΓΑ;
πενταπλάσια ἄρα τὰ ἀπὸ τῶν ΒΓ, ΓΑ τοῦ ἀπὸ
τῆς ΓΑ. ὑπόκειται δὲ καὶ τὸ ἀπὸ τῆς ΒΑ πενταπλάσιον
τοῦ ἀπὸ τῆς ΓΑ; τὸ ἄρα ἀπὸ τῆς ΒΑ ἴσον ἐστὶ
τοῖς ἀπὸ τῶν ΒΓ, ΓΑ; ὅπερ ἀδύνατον. οὐκ ἄρα
ἡ ΓΒ διπλασία ἐστὶ τῆς ΑΓ. ὁμοίως δὴ δείχομεν,
ὅτι οὐδὲ ἡ ἐλάττων τῆς ΓΒ διπλασίων ἐστὶ τῆς ΓΑ;
πολλῷ γὰρ [μεῖζον] τὸ ἄτοπον.}

{\greekfont Ἡ ἄρα τῆς ΑΓ διπλῆ μείζων ἐστὶ τῆς ΓΒ; ὅπερ
ἔδει δεῖχαι.}}

\ParallelRText{
\begin{center}
{\large Lemma}
\end{center}

And it can be shown that double $AC$ ({\em i.e.}, $DC$) is greater than $BC$, as follows.

For if  (double $AC$ is) not (greater than $BC$), if possible, let $BC$ be  double $CA$. Thus, the (square) on $BC$ (is) four times the (square)  on $CA$. Thus, the (sum of) the (squares)
on $BC$ and $CA$ (is) five times the (square)  on $CA$. And the
(square) on $BA$ was assumed (to be) five times the (square) on $CA$. 
Thus, the (square) on $BA$ is equal to the (sum of) the (squares) on $BC$
and $CA$. The very thing (is) impossible [Prop.~2.4]. Thus,
$CB$ is not  double $AC$. So, similarly, we can show that 
a (straight-line) less than $CB$ is not double $AC$ either. 
For (in this case) the absurdity is much [greater].

Thus, double $AC$ is greater than $CB$. (Which is) the very thing it
was required to show.}
\end{Parallel}

%%%%
%13.3
%%%%
\pdfbookmark[1]{Proposition 13.3}{pdf13.3}
\begin{Parallel}{}{}
\ParallelLText{
\begin{center}
{\large \ggn{3}.}
\end{center}\vspace*{-7pt}

{\greekfont Ἐὰν εὐθεῖα γραμμὴ ἄκρον καὶ μέσον 
λόγον τμ ηθῇ, τὸ ἔλασσον τμῆμα προσλαβὸν τὴν ἡμίσειαν
τοῦ μείζονος τμήματος πενταπλάσιον δύναται
τοῦ ἀπὸ τῆς ἡμισείας τοῦ μείζονος τμήματος
τετραγώνου.}

\epsfysize=2.5in
\centerline{\epsffile{Book13/fig03g.eps}}

{\greekfont Εὐθεῖα γάρ τις ἡ ΑΒ ἄκρον καὶ μέσον λόγον τετμήσθω
κατὰ τὸ Γ σημεῖον, καὶ ἔστω μεῖζον τμῆμα
τὸ ΑΓ, καὶ τετμήσθω ἡ ΑΓ δίθα κατὰ τὸ Δ; λέγω,
ὅτι πενταπλάσιόν ἐστι τὸ ἀπὸ τῆς ΒΔ τοῦ ἀπὸ
τῆς ΔΓ.}

{\greekfont Ἀναγεγράφθω γὰρ ἀπὸ τῆς ΑΒ τετράγωνον τὸ ΑΕ, καὶ
καταγεγράφθω διπλοῦν τὸ σχῆμα. ἐπεὶ διπλῆ
ἐστιν ἡ ΑΓ τῆς ΔΓ, τετραπλάσιον ἄρα τὸ ἀπὸ τῆς ΑΓ
τοῦ ἀπὸ τῆς ΔΓ, τουτέστι τὸ ΡΣ τοῦ ΖΗ. καὶ ἐπεὶ
τὸ ὑπὸ τῶν ΑΒΓ ἴσον ἐστὶ τῷ ἀπὸ τῆς ΑΓ,
καί ἐστι τὸ ὑπὸ τῶν ΑΒΓ τὸ ΓΕ, τὸ ἄρα ΓΕ ἴσον
ἐστὶ τῷ ΡΣ. τετραπλάσιον δὲ τὸ ΡΣ τοῦ ΖΗ; τετραπλάσιον
ἄρα καὶ τὸ ΓΕ τοῦ ΖΗ. πάλιν ἐπεὶ ἴση ἐστὶν ἡ ΑΔ
τῇ ΔΓ, ἴση ἐστὶ καὶ ἡ ΘΚ τῇ ΚΖ. ὥστε καὶ τὸ ΗΖ
τετράγωνον ἴσον ἐστὶ  τῷ ΘΛ τετραγώνῳ.
ἴση ἄρα ἡ ΗΚ τῇ ΚΛ, τουτέστιν ἡ ΜΝ τῇ ΝΕ;
ὥστε καὶ τὸ  ΜΖ τῷ ΖΕ ἐστιν ἴσον. ἀλλὰ
τὸ ΜΖ τῷ ΓΗ ἐστιν ἴσον; καὶ τὸ ΓΗ ἄρα τῷ ΖΕ ἐστιν
ἴσον. κοινὸν προσκείσθω τὸ ΓΝ; ὁ ἄρα ΧΟΠ γνώμων
ἴσος ἐστὶ τῷ ΓΕ. ἀλλὰ τὸ ΓΕ τετραπλάσιον ἐδείθθῃ τοῦ
ΗΖ; καὶ ὁ ΧΟΠ ἄρα γνώμων τετραπλάσιός ἐστι τοῦ ΖΗ
τετραγώνου. ὁ ΧΟΠ ἄρα γνώμων καὶ τὸ ΖΗ τετράγωνον
πενταπλάσιός ἐστι τοῦ ΖΗ. ἀλλὰ ὁ ΧΟΠ γνώμων καὶ τὸ
ΖΗ τετράγωνόν ἐστι τὸ ΔΝ. καί
ἐστι τὸ μὲν ΔΝ τὸ ἀπὸ τῆς ΔΒ, τὸ δὲ ΗΖ τὸ ἀπὸ
τῆς ΔΓ. τὸ ἄρα ἀπὸ τῆς ΔΒ πενταπλάσιόν ἐστι τοῦ
ἀπὸ τῆς ΔΓ; ὅπερ ἔδει δεῖχαι.}}

\ParallelRText{
\begin{center}
{\large Proposition 3}
\end{center}

If a straight-line is cut in extreme and mean ratio then the square
on the lesser piece added to half of the greater piece is five
times the square on half of the greater piece.\\

\epsfysize=2.5in
\centerline{\epsffile{Book13/fig03e.eps}}

For let some straight-line $AB$ be cut in extreme and mean ratio
at point $C$. And let $AC$ be the greater piece. And let $AC$ be
cut in half at $D$. I say that the (square) on $BD$ is five times the
(square) on $DC$. 

For let the square $AE$ be described on $AB$. And let the
figure be drawn double. Since $AC$ is double $DC$, the
(square) on $AC$ (is) thus four times the (square) on $DC$---that
is to say, $RS$ (is four times) $FG$. And since the (rectangle contained)
by $ABC$ is equal to the (square) on $AC$ [Def.~6.3, Prop.~6.17],
and $CE$ is the (rectangle contained) by $ABC$,  $CE$ is thus equal
to $RS$. And $RS$ (is) four times $FG$. Thus, $CE$ (is) also four times
$FG$. Again, since $AD$ is equal to $DC$, $HK$ is also equal to
$KF$. Hence, square $GF$ is also equal to square $HL$. Thus,
$GK$ (is) equal to $KL$---that is to say, $MN$ to $NE$. 
Hence, $MF$ is also equal to $FE$. But, $MF$ is equal to $CG$.
Thus, $CG$ is also equal to $FE$. Let $CN$ be added to both.
Thus, gnomon $OPQ$ is equal to $CE$. But, $CE$ was shown (to be)
equal to four times $GF$. Thus, gnomon $OPQ$  is also four times
square $FG$. Thus, gnomon $OPQ$ plus square $FG$ is five times
$FG$. But, gnomon $OPQ$ plus square $FG$ is (square) $DN$. 
And  $DN$ is the (square) on $DB$, and $GF$ the (square)
on $DC$. Thus, the (square) on $DB$ is five times the (square)
on $DC$. (Which is) the very thing it was required to show.}
\end{Parallel}

%%%%
%13.4
%%%%
\pdfbookmark[1]{Proposition 13.4}{pdf13.4}
\begin{Parallel}{}{}
\ParallelLText{
\begin{center}
{\large \ggn{4}.}
\end{center}\vspace*{-7pt}

{\greekfont Ἐὰν εὐθεῖα γραμμὴ ἄκρον καὶ μέσον λόγον τμ ηθῇ,
τὸ ἀπὸ τῆς ὅλης καὶ τοῦ ἐλάσσονος τμήματος, τὰ συναμφότερα
τετράγωνα, τριπλάσιά ἐστι τοῦ ἀπὸ τοῦ μείζονος τμήματος
τετραγώνου.}

\epsfysize=2.5in
\centerline{\epsffile{Book13/fig04g.eps}}

{\greekfont Ἔστω εὐθεῖα ἡ ΑΒ, καὶ τετμήσθω ἄκρον καὶ μέσον
λόγον κατὰ τὸ Γ, καὶ ἔστω μεῖζον τμῆμα τὸ ΑΓ; λέγω,
ὅτι τὰ ἀπὸ τῶν ΑΒ, ΒΓ τριπλάσιά ἐστι τοῦ ἀπὸ τῆς
ΓΑ.}

{\greekfont Ἀναγεγράφθω γὰρ ἀπὸ τῆς ΑΒ τετράγωνον τὸ ΑΔΕΒ, καὶ
καταγεγράφθω τὸ σχῆμα. ἐπεὶ οὖν ἡ ΑΒ ἄκρον καὶ
μέσον λόγον τέτμ ηται κατὰ τὸ Γ, καὶ τὸ μεῖζον τμῆμά ἐστιν
ἡ ΑΓ, τὸ ἄρα ὑπὸ τῶν ΑΒΓ ἴσον ἐστὶ τῷ ἀπὸ τῆς
ΑΓ. καί ἐστι τὸ μὲν ὑπὸ τῶν ΑΒΓ τὸ ΑΚ, τὸ δὲ ἀπὸ
τῆς ΑΓ τὸ ΘΗ; ἴσον ἄρα ἐστὶ τὸ ΑΚ τῷ ΘΗ. καὶ ἐπεὶ
ἴσον ἐστὶ τὸ ΑΖ τῷ ΖΕ, κοινὸν προσκείσθω τὸ ΓΚ; ὅλον
ἄρα τὸ ΑΚ ὅλῳ τῷ ΓΕ ἐστιν ἴσον; τὰ ἄρα ΑΚ, ΓΕ
τοῦ ΑΚ ἐστι διπλάσια. ἀλλὰ τὰ ΑΚ, ΓΕ ὁ ΛΜΝ γνώμων
ἐστὶ καὶ τὸ ΓΚ τετράγωνον; ὁ ἄρα ΛΜΝ γνώμων
καὶ τὸ ΓΚ τετράγωνον  διπλάσιά ἐστι τοῦ ΑΚ. ἀλλὰ μὴν καὶ τὸ
ΑΚ τῷ ΘΗ ἐδείθθη ἴσον; ὁ ἄρα ΛΜΝ γνώμων καὶ
[τὸ ΓΚ τετράγωνον διπλάσιά ἐστι τοῦ ΘΗ; ὥστε ὁ ΛΜΝ
γνώμων καὶ] τὰ ΓΚ, ΘΗ τετράγωνα τριπλάσιά ἐστι
τοῦ ΘΗ τετραγώνου. καί ἐστιν ὁ [μὲν] ΛΜΝ γνώμων
καὶ τὰ ΓΚ, ΘΗ τετράγωνα ὅλον τὸ ΑΕ καὶ τὸ ΓΚ, ἅπερ
ἐστὶ τὰ ἀπὸ τῶν ΑΒ, ΒΓ τετράγωνα, τὸ δὲ ΗΘ τὸ ἀπὸ
τῆς ΑΓ τετράγωνον. τὰ ἄρα ἀπὸ τῶν ΑΒ, ΒΓ τετράγωνα
τριπλάσιά ἐστι τοῦ ἀπὸ τῆς ΑΓ τετραγώνου; ὅπερ ἔδει
δεῖχαι.}}

\ParallelRText{
\begin{center}
{\large Proposition 4}
\end{center}

If a straight-line is cut in extreme and mean ratio then the sum of the
squares on the whole and the lesser piece is three times the square on
the greater piece.\\

\epsfysize=2.5in
\centerline{\epsffile{Book13/fig04e.eps}}

Let $AB$ be a straight-line, and let it be cut in extreme and mean
ratio at $C$, and let $AC$ be the greater piece. I say that the (sum of the
squares) on $AB$ and $BC$ is three times the (square) on $CA$.

For let the square $ADEB$ be described on $AB$, and
let the (remainder of the) figure be drawn. Therefore, since
$AB$ has been cut in extreme and mean ratio at $C$,  and $AC$ is the greater piece, the (rectangle contained) by $ABC$ is thus equal to the
(square) on $AC$ [Def.~6.3, Prop.~6.17]. And $AK$ is the (rectangle
contained) by $ABC$, and $HG$ the (square) on $AC$. Thus,
$AK$ is equal to $HG$.  And since $AF$ is equal to $FE$ [Prop.~1.43],
let $CK$ be added to both. Thus, the whole of $AK$  is equal to
the whole of $CE$. Thus, $AK$ plus $CE$ is double $AK$. But, $AK$
plus $CE$ is the gnomon $LMN$ plus the square $CK$. Thus, 
gnomon $LMN$ plus square $CK$ is double $AK$. But, indeed,
$AK$ was also shown (to be) equal to $HG$. Thus, gnomon $LMN$
plus [square $CK$ is double $HG$. Hence, gnomon $LMN$ plus] the
 squares $CK$ and $HG$ is three times the square $HG$.
And  gnomon $LMN$ plus the  squares  $CK$ and
$HG$ is the whole of $AE$ plus $CK$---which are the 
squares on $AB$ and $BC$ (respectively)---and $GH$ (is) the square on $AC$. 
Thus, the (sum of the) squares on $AB$ and $BC$ is three times the
square on $AC$. (Which is) the very thing it was required to show.}
\end{Parallel}

%%%%
%13.5
%%%%
\pdfbookmark[1]{Proposition 13.5}{pdf13.5}
\begin{Parallel}{}{}
\ParallelLText{
\begin{center}
{\large \ggn{5}.}
\end{center}\vspace*{-7pt}

{\greekfont Ἐὰν εὐθεῖα γραμμὴ ἄκρον καὶ μέσον λόγον τμ ηθῇ, καὶ
προστεθῇ α ὐτῇ ἴση τῷ μείζονι τμήματι, ἡ ὅλη εὐθεῖα
ἄκρον καὶ μέσον λόγον τέτμ ηται, καὶ τὸ μεῖζον τμῆμά ἐστιν
ἡ ἐχ ἀρθῆς εὐθεῖα.}

\epsfysize=2.in
\centerline{\epsffile{Book13/fig05g.eps}}

{\greekfont Εὐθεῖα γὰρ γραμμὴ ἡ ΑΒ ἄκρον καὶ μέσον λόγον τετμήσθω
κατὰ τὸ Γ σημεῖον, καὶ ἔστω μεῖζον τμῆμα ἡ ΑΓ, καὶ τῇ
ΑΓ ἴση [κείσθω] ἡ ΑΔ. λέγω, ὅτι ἡ ΔΒ εὐθεῖα ἄκρον
καὶ μέσον λόγον τέτμ ηται κατὰ τὸ Α, καὶ τὸ μεῖζον τμῆμά ἐστιν
ἡ ἐχ ἀρθῆς εὐθεῖα ἡ ΑΒ.}

{\greekfont Ἀναγεγράφθω γὰρ ἀπὸ τῆς ΑΒ τετράγωνον τὸ ΑΕ, καὶ
καταγεγράφθω τὸ σχῆμα. ἐπεὶ ἡ ΑΒ ἄκρον καὶ μέσον
λόγον τέτμ ηται κατὰ τὸ Γ, τὸ ἄρα ὑπὸ ΑΒΓ ἴσον ἐστὶ
τῷ ἀπὸ ΑΓ. καί ἐστι τὸ μὲν ὑπὸ ΑΒΓ τὸ ΓΕ, τὸ δὲ ἀπὸ
τῆς ΑΓ τὸ ΓΘ; ἴσον ἄρα τὸ ΓΕ τῷ ΘΓ. ἀλλὰ τῷ μὲν ΓΕ ἴσον
ἐστὶ τὸ ΘΕ, τῷ δὲ ΘΓ ἴσον τὸ ΔΘ; καὶ τὸ ΔΘ ἄρα ἴσον ἐστὶ
τῷ ΘΕ [κοινὸν προσκείσθω τὸ ΘΒ]. ὅλον ἄρα τὸ ΔΚ ὅλῳ τῷ
ΑΕ ἐστιν ἴσον. καί ἐστι τὸ μὲν ΔΚ τὸ ὑπὸ τῶν ΒΔ, ΔΑ; ἴση
γὰρ ἡ ΑΔ τῇ ΔΛ; τὸ δὲ ΑΕ τὸ ἀπὸ τῆς ΑΒ; τὸ ἄρα ὑπὸ
τῶν ΒΔΑ ἴσον ἐστὶ τῷ ἀπὸ τῆς ΑΒ. ἔστιν ἄρα ὡς ἡ ΔΒ πρὸς
τὴν ΒΑ, οὕτως ἡ ΒΑ πρὸς τὴν ΑΔ. μείζων δὲ ἡ ΔΒ τῆς ΒΑ;
μείζων ἄρα καὶ ἡ ΒΑ τῆς ΑΔ.}

{\greekfont Ἡ ἄρα ΔΒ ἄκρον καὶ μέσον λόγον τέτμ ηται κατὰ τὸ Α,
καὶ τὸ μεῖζον τμῆμά ἐστιν ἡ ΑΒ; ὅπερ ἔδει δεῖχαι.}}

\ParallelRText{
\begin{center}
{\large Proposition 5}
\end{center}

If a straight-line is cut in extreme and mean ratio, and a
(straight-line) equal to the greater piece is added to it, then the
whole straight-line has been cut in extreme and mean ratio, and
 the original straight-line is the greater piece.

\epsfysize=2.in
\centerline{\epsffile{Book13/fig05e.eps}}

For let the straight-line $AB$ be cut in extreme and mean ratio
at point $C$. And let $AC$ be the greater piece. And let $AD$
be [made] equal to $AC$. I say that the straight-line $DB$ has been cut in extreme and
mean ratio at $A$, and that the original straight-line $AB$ is
the greater piece.

For let the square $AE$ be described on $AB$, and
let the (remainder of the) figure be drawn. And since
$AB$ has been cut in extreme and mean ratio at $C$, the (rectangle
contained) by $ABC$ is thus equal to the (square) on $AC$
[Def.~6.3, Prop.~6.17].  And $CE$ is the (rectangle contained)
by  $ABC$,  and $CH$ the (square) on $AC$. But, $HE$ is equal
to $CE$ [Prop.~1.43], and $DH$ equal to $HC$. Thus, $DH$ is also
equal to $HE$. [Let $HB$ be added to both.] Thus, the whole
of $DK$ is equal to the whole of $AE$. And $DK$ is the
(rectangle contained) by $BD$ and $DA$. For $AD$ (is) equal to
$DL$. And $AE$ (is) the (square) on $AB$.  Thus, the (rectangle
contained) by $BDA$ is equal to the (square) on $AB$. Thus, as
$DB$ (is) to $BA$, so $BA$ (is) to $AD$ [Prop.~6.17]. And
$DB$ (is) greater than $BA$. Thus, $BA$ (is) also greater
than $AD$ [Prop.~5.14].

Thus, $DB$ has been cut in extreme and mean ratio at $A$, and the
greater piece is $AB$. (Which is) the very thing it was required to show.}
\end{Parallel}

%%%%
%13.6
%%%%
\pdfbookmark[1]{Proposition 13.6}{pdf13.6}
\begin{Parallel}{}{}
\ParallelLText{
\begin{center}
{\large \ggn{6}.}
\end{center}\vspace*{-7pt}

{\greekfont Ἐὰν εὐθεῖα ῥητη ἄκρον καὶ μέσον λόγον τμ ηθῇ, ἑκάτερον
τῶν τμ ημάτων ἄλογός ἐστιν ἡ καλουμένη ἀποτομή.}\\

\epsfysize=0.3in
\centerline{\epsffile{Book13/fig06g.eps}}

{\greekfont Ἔστω εὐθεῖα ῥητὴ ἡ ΑΒ καὶ τετμήσθω ἄκρον καὶ μέσον
λόγον κατὰ τὸ Γ, καὶ ἔστω μεῖζον τμῆμα ἡ ΑΓ; λέγω, ὅτι
ἑκατέρα τῶν ΑΓ, ΓΒ ἄλογός ἐστιν ἡ καλουμένη ἀποτομή.}

{\greekfont Ἐκβεβλήσθω γὰρ ἡ ΒΑ, καὶ κείσθω τῆς ΒΑ ἡμίσεια ἡ ΑΔ. ἐπεὶ
οὖν εὐθεῖα ἡ ΑΒ τέτμ ηται ἄκρον
καὶ μέσον λόγον κατὰ τὸ Γ, καὶ τῷ μείζονι τμήματι τῷ ΑΓ
πρόσκειται ἡ ΑΔ ἡμίσεια οὖσα τῆς ΑΒ, τὸ ἄρα ἀπὸ ΓΔ
τοῦ ἀπὸ ΔΑ πενταπλάσιόν ἐστιν. τὸ ἄρα ἀπὸ ΓΔ πρὸς τὸ
ἀπὸ ΔΑ λόγον ἔχει, ὃν ἀριθμὸς πρὸς ἀριθμόν; σύμμετρον
ἄρα τὸ ἀπὸ ΓΔ τῷ ἀπὸ
ΔΑ. ῥητὸν δὲ τὸ ἀπὸ ΔΑ; ῥητὴ γάρ [ἐστιν]
ἡ ΔΑ ἡμίσεια οὖσα τῆς ΑΒ ῥητῆς οὔσης; ῥητὸν
ἄρα καὶ τὸ ἀπὸ ΓΔ; ῥητὴ ἄρα ἐστὶ καὶ ἡ ΓΔ. καὶ ἐπεὶ
τὸ ἀπὸ ΓΔ πρὸς τὸ ἀπὸ ΔΑ λόγον οὐκ ἔχει, ὃν τετράγωνος
ἀριθμὸς πρὸς τετράγωνον ἀριθμόν, ἀσύμμετρος ἄρα μήκει
ἡ ΓΔ τῇ ΔΑ; αἱ ΓΔ, ΔΑ ἄρα ῥηταί εἰσι δυνάμει μόνον
σύμμετροι; ἀποτομὴ ἄρα ἐστὶν ἡ ΑΓ.  πάλιν, ἐπεὶ ἡ ΑΒ
ἄκρον καὶ μέσον λόγον τέτμ ηται, καὶ τὸ μεῖζον  τμῆμά
ἐστιν ἡ ΑΓ, τὸ ἄρα ὑπὸ ΑΒ, ΒΓ τῷ ἀπὸ ΑΓ ἴσον
ἐστίν. τὸ ἄρα ἀπὸ τῆς ΑΓ ἀποτομῆς παρὰ τὴν ΑΒ ῥητὴν
παραβληθὲν πλάτος ποιεῖ τὴν ΒΓ. τὸ δὲ ἀπὸ ἀποτομῆς
παρὰ ῥητὴν παραβαλλόμενον πλάτος ποιεῖ ἀποτομὴν πρώτην; ἀποτομὴ
ἄρα πρώτη ἐστὶν ἡ ΓΒ. ἐδείθθη δὲ καὶ ἡ ΓΑ ἀποτομή.}

{\greekfont Ἐὰν ἄρα εὐθεῖα ῥητὴ ἄκρον καὶ μέσον λόγον τμ ηθῇ,
ἑκάτερον τῶν τμ ημάτων ἄλογός ἐστιν ἡ καλουμένη
ἀποτομή; ὅπερ ἔδει δεῖχαι.}}

\ParallelRText{
\begin{center}
{\large Proposition 6}
\end{center}

If a rational straight-line is cut in extreme and mean ratio then each of
the pieces is that irrational (straight-line) called an apotome.

\epsfysize=0.3in
\centerline{\epsffile{Book13/fig06e.eps}}

Let  $AB$ be a rational straight-line cut in extreme and mean ratio
at $C$, and let $AC$ be the greater piece. I say that  $AC$ and
$CB$ is each that irrational (straight-line) called an apotome.

For let $BA$ be produced, and let $AD$ be made (equal) to
half of $BA$. Therefore, since the straight-line $AB$ has been cut in
extreme and mean ratio at $C$, and $AD$, which is half of $AB$, has been added to the greater piece $AC$, the (square) on $CD$ is thus 
five times the (square) on $DA$ [Prop.~13.1]. Thus, the (square) on $CD$
has to the (square) on $DA$ the ratio which a number (has) to a number.
The (square) on $CD$ (is) thus commensurable with the (square) on
$DA$ [Prop.~10.6]. And the (square) on $DA$ (is) rational.
For $DA$ [is] rational, being half of $AB$, which is rational.
Thus, the (square) on $CD$ (is) also rational [Def.~10.4]. 
Thus, $CD$ is also rational. And since the (square) on $CD$ does not have to
the (square) on $DA$ the ratio which a square number (has) to a square
number, $CD$ (is) thus incommensurable in length with $DA$ [Prop.~10.9].
Thus, $CD$  and $DA$ are rational (straight-lines which are) commensurable in square only. Thus, $AC$ is an apotome [Prop.~10.73].  Again, since
$AB$ has been cut in extreme and mean ratio, and $AC$ is the greater
piece, the (rectangle contained) by $AB$ and $BC$ is thus equal
to the (square) on $AC$ [Def.~6.3, Prop.~6.17]. Thus, the (square) on
the apotome $AC$, applied to the rational (straight-line) $AB$, makes $BC$
as width. And the (square) on an apotome, applied to a rational
(straight-line), makes a first apotome as width [Prop.~10.97].
Thus, $CB$ is a first apotome. And $CA$ was also
shown (to be) an apotome.

Thus, if a rational straight-line is cut in extreme and mean ratio then each of
the pieces is that irrational (straight-line) called an apotome.}
\end{Parallel}

%%%%
%13.7
%%%%
\pdfbookmark[1]{Proposition 13.7}{pdf13.7}
\begin{Parallel}{}{}
\ParallelLText{
\begin{center}
{\large \ggn{7}.}
\end{center}\vspace*{-7pt}

{\greekfont Ἐὰν πενταγώνου ἰσοπλεύρου αἱ τρεῖς γωνίαι ἤτοι αἱ κατὰ
τὸ ἑχῆς ἢ αἱ μὴ κατὰ τὸ ἑχῆς ἴσαι ὦσιν, ἰσογώνιον ἔσται
τὸ πεντάγωνον.}

{\greekfont Πενταγώνου γὰρ ἰσοπλεύρον τοῦ ΑΒΓΔΕ αἱ τρεῖς γωνίαι
πρότερον αἱ κατὰ τὸ ἑχῆς αἱ πρὸς τοῖς Α, Β, Γ ἴσαι ἀλλήλαις
ἔστωσαν; λέγω, ὅτι ἰσογώνιόν ἐστι τὸ ΑΒΓΔΕ πεντάγωνον.}

{\greekfont Ἐπεζεύθθωσαν γὰρ αἱ ΑΓ, ΒΕ, ΖΔ. καὶ ἐπεὶ δύο αἱ ΓΒ, ΒΑ
δυσὶ ταῖς ΒΑ, ΑΕ ἴσαι ἐισὶν ἑκατέρα ἑκατέρᾳ, καὶ γωνία ἡ ὑπὸ ΓΒΑ γωνίᾳ
τῇ ὑπὸ ΒΑΕ ἐστιν ἴση, βάσις ἄρα ἡ ΑΓ βάσει τῇ ΒΕ ἐστιν
ἴση, καὶ τὸ ΑΒΓ τρίγωνον τῷ ΑΒΕ τριγώνῳ ἴσον, καὶ αἱ
λοιπαὶ γωνίαι ταῖς λοιπαῖς γωνίαις ἴσαι ἔσονται, ὑφ ἃς αἱ
ἴσαι πλευραὶ ὑποτείνουσιν, ἡ μὲν ὑπὸ ΒΓΑ τῇ ὑπὸ ΒΕΑ,
ἡ δὲ ὑπὸ ΑΒΕ τῇ ὑπὸ ΓΑΒ; ὥστε καὶ πλευρὰ ἡ ΑΖ πλευρᾷ
τῇ ΒΖ ἐστιν ἴση. ἐδείθθη δὲ καὶ ὅλη ἡ ΑΓ ὅλῃ τῇ
ΒΕ ἴση; καὶ λοιπὴ ἄρα ἡ ΖΓ λοιπῇ τῇ ΖΕ ἐστιν ἴση. ἔστι
δὲ καὶ ἡ ΓΔ τῇ ΔΕ ἴση. δύο δὴ αἱ ΖΓ, ΓΔ δυσὶ ταῖς
ΖΕ, ΕΔ ἴσαι εἰσίν; καὶ βάσις α ὐτῶν κοινὴ ἡ ΖΔ; γωνία
ἄρα ἡ ὑπὸ ΖΓΔ γωνίᾳ τῇ ὑπὸ ΖΕΔ ἐστιν ἴση. ἐδείθθη
δὲ καὶ ἡ ὑπὸ ΒΓΑ τῇ ὑπὸ ΑΕΒ ἴση; καὶ ὅλη ἄρα ἡ ὑπὸ
ΒΓΔ ὅλῃ τῇ ὑπὸ ΑΕΔ ἴση. ἀλλ ἡ ὑπὸ ΒΓΔ ἴση ὑπόκειται
ταῖς πρὸς τοῖς Α, Β γωνίαις; καὶ ἡ ὑπὸ ΑΕΔ ἄρα ταῖς πρὸς
τοῖς Α, Β γωνίαις ἴση ἐστίν. ὁμοίως δὴ δείχομεν, ὅτι
καὶ ἡ ὑπὸ ΓΔΕ γωνία ἴση ἐστὶ ταῖς πρὸς τοῖς Α, Β, Γ γωνίαις;
ἰσογώνιον ἄρα ἐστὶ τὸ ΑΒΓΔΕ πεντάγωνον.}\\~\\

\epsfysize=2.in
\centerline{\epsffile{Book13/fig07g.eps}}

{\greekfont Ἀλλὰ δὴ μὴ ἔστωσαν ἴσαι αἱ κατὰ τὸ ἑχῆς γωνίαι,
ἀλλ ἔστωσαν ἴσαι αἱ πρὸς τοῖς Α, Γ, Δ σημείοις;
λέγω, ὅτι καὶ οὕτως ἰσογώνιόν ἐστι τὸ ΑΒΓΔΕ πεντάγωνον.}

{\greekfont Ἐπεζεύθθω γὰρ ἡ ΒΔ. καὶ ἐπεὶ δύο αἱ ΒΑ, ΑΕ δυσὶ ταῖς
ΒΓ, ΓΔ ἴσαι εἰσὶ καὶ γωνίας ἴσας περιέθουσιν, βάσις ἄρα ἡ ΒΕ βάσει τῇ ΒΔ ἴση ἐστίν, καὶ τὸ ΑΒΕ τρίγωνον τῷ ΒΓΔ τριγώνῳ
ἴσον ἐστίν, καὶ αἱ λοιπαὶ γωνίαι ταῖς λοιπαῖς γωνίαις
ἴσαι ἔσονται, ὑφ  ἃς αἱ ἴσαι πλευραὶ ὑποτείνουσιν; ἴση
ἄρα ἐστὶν ἡ ὑπὸ ΑΕΒ γωνία τῇ ὑπὸ ΓΔΒ. ἔστι δὲ καὶ ἡ
ὑπὸ ΒΕΔ γωνία τῇ ὑπὸ ΒΔΕ ἴση, ἐπεὶ καὶ πλευρὰ
ἡ ΒΕ πλευρᾷ τῇ ΒΔ ἐστιν ἴση.
 καὶ ὅλη ἄρα ἡ ὑπὸ
ΑΕΔ γωνία ὅλῃ τῇ ὑπὸ ΓΔΕ ἐστιν  ἴση. ἀλλὰ ἡ ὑπὸ
ΓΔΕ ταῖς πρὸς τοῖς Α, Γ γωνίαις ὑπόκειται ἴση;  καὶ ἡ ὑπὸ
ΑΕΔ ἄρα γωνία ταῖς πρὸς τοῖς Α, Γ ἴση ἐστίν. διὰ τὰ α ὐτὰ
δὴ καὶ ἡ ὑπὸ ΑΒΓ ἴση ἐστὶ ταῖς πρὸς τοῖς Α, Γ, Δ γωνίαις.
ἰσογώνιον ἄρα ἐστὶ τὸ ΑΒΓΔΕ πεντάγωνον; ὅπερ ἔδει
δεῖχαι.}}

\ParallelRText{
\begin{center}
{\large Proposition 7}
\end{center}

If  three angles, either  consecutive or not consecutive, of an equilateral pentagon
are equal then the pentagon will be equiangular.

For let three angles of the equilateral pentagon $ABCDE$---first of
all, the consecutive (angles) at $A$, $B$, and $C$----be equal to
one another. I say that pentagon $ABCDE$ is equiangular.

For let $AC$, $BE$, and $FD$ be joined. And since the two (straight-lines) $CB$ and $BA$ are equal to the two (straight-lines) $BA$ and
$AE$, respectively, and angle $CBA$ is equal to angle $BAE$,  base
$AC$ is thus equal to base $BE$, and triangle $ABC$  equal to triangle
$ABE$, and the remaining angles will be equal to the remaining angles
which the equal sides subtend [Prop.~1.4],  (that is), $BCA$ (equal) to $BEA$, and $ABE$ to $CAB$. And hence side $AF$ is also equal to side $BF$ [Prop.~1.6]. And the whole  of $AC$ was also shown (to be) equal to
the whole of $BE$. Thus, the remainder $FC$ is also equal to the remainder $FE$.
And $CD$ is also equal to $DE$. So, the two (straight-lines) $FC$ and
$CD$ are equal to the two (straight-lines) $FE$ and $ED$ (respectively). 
And $FD$ is their common base. Thus, angle $FCD$ is equal to angle
$FED$ [Prop.~1.8]. And $BCA$ was also shown (to be) equal to 
$AEB$. And thus the whole of  $BCD$ (is) equal to the whole of $AED$.
But, (angle) $BCD$ was assumed (to be) equal to the angles at $A$ and
$B$. Thus, (angle) $AED$ is also equal to the angles at $A$ and $B$. So, 
similarly, we can show that angle $CDE$ is also equal to the angles
at $A$, $B$, $C$. Thus, pentagon $ABCDE$ is equiangular.

\epsfysize=2.in
\centerline{\epsffile{Book13/fig07e.eps}}

And so let consecutive  angles not be equal, but let
the (angles) at points $A$, $C$, and $D$ be equal. I say that pentagon $ABCDE$ is also
equiangular in this case.

For let $BD$ be joined. And since the two (straight-lines) $BA$
and $AE$ are equal to the (straight-lines) $BC$ and $CD$, and they contain
equal angles,   base $BE$ is thus equal to base $BD$, and triangle $ABE$
is equal to triangle $BCD$, and the remaining angles will be equal
to the remaining angles which the equal sides subtend [Prop.~1.4]. 
Thus, angle $AEB$ is equal to (angle) $CDB$. And angle $BED$
is also equal to (angle) $BDE$,  since side $BE$ is also equal to side
$BD$ [Prop.~1.5]. Thus, the whole angle $AED$ is also equal to the
whole (angle) $CDE$. But, (angle) $CDE$ was assumed (to be) equal
to the angles at $A$ and $C$.  Thus, angle $AED$ is also equal to the (angles)
at $A$ and $C$. So, for the same (reasons), (angle) $ABC$ is also equal
to the angles at  $A$, $C$, and $D$. Thus, pentagon $ABCDE$ is equiangular.
(Which is) the very thing it was required to show.}
\end{Parallel}

%%%%
%13.8
%%%%
\pdfbookmark[1]{Proposition 13.8}{pdf13.8}
\begin{Parallel}{}{}
\ParallelLText{
\begin{center}
{\large \ggn{8}.}
\end{center}\vspace*{-7pt}

{\greekfont Ἐὰν πενταγώνου ἰσοπλεύρου καὶ ἰσογωνίου τὰς κατὰ
τὸ ἑχῆς δύο γωνίας ὑποτείνωσιν εὐθεῖαι, ἄκρον καὶ
μέσον λόγον τέμνουσιν ἀλλήλας, καὶ τὰ μείζονα α ὐτῶν
τμήματα ἴσα ἐστὶ τῇ τοῦ πενταγώνου πλευρᾷ.}

{\greekfont Πενταγώνου γὰρ ἰσοπλεύρον καὶ ἰσογωνίου τοῦ ΑΒΓΔΕ δύο
γωνίας τὰς κατὰ τὸ ἑχῆς τὰς πρὸς τοῖς Α, Β ὑποτεινέτωσαν εὐθεῖαι
αἱ ΑΓ, ΒΕ τέμνουσαι ἀλλήλας κατὰ τὸ Θ σημεὶον; λέγω, ὅτι
ἑκατέρα α ὐτῶν ἄκρον καὶ μέσον λόγον τέτμ ηται κατὰ τὸ Θ
σημεῖον, καὶ τὰ μείζονα α ὐτῶν τμήματα ἴσα ἐστὶ τῇ τοῦ
πενταγώνου πλευρᾷ.}

\epsfysize=2.in
\centerline{\epsffile{Book13/fig08g.eps}}

{\greekfont Περιγεγράφθω γὰρ περὶ τὸ ΑΒΓΔΕ πεντάγωνον κύκλος ὁ
ΑΒΓΔΕ. καὶ ἐπεὶ δύο εὐθεῖαι αἱ ΕΑ, ΑΒ δυσὶ ταῖς ΑΒ, ΒΓ
ἴσαι εἰσὶ καὶ γωνίας ἴσας περιέθουσιν, βάσις ἄρα ἡ ΒΕ βάσει
τῇ ΑΓ ἴση ἐστίν, καὶ τὸ ΑΒΕ τρίγωνον τῷ ΑΒΓ τριγώνῳ
ἴσον ἐστίν, καὶ αἱ λοιπαὶ γωνίαι ταῖς λοιπαῖς γωνίαις
ἴσαι ἔσονται ἑκατέρα ἑκατέρᾳ, ὑφ ἃς αἱ ἴσαι πλευραὶ
ὑποτείνουσιν. ἴση ἄρα ἐστὶν ἡ ὑπὸ ΒΑΓ γωνία τῇ
ὑπὸ ΑΒΕ; διπλῆ ἄρα ἡ ὑπὸ ΑΘΕ τῆς ὑπὸ ΒΑΘ. ἔστι
δὲ καὶ ἡ ὑπὸ ΕΑΓ τῆς ὑπὸ ΒΑΓ διπλῆ, ἐπειδήπερ καὶ
περιφέρεια ἡ ΕΔΓ περιφερείας τῆς ΓΒ ἐστι διπλῆ; ἴση ἄρα
ἡ ὑπὸ ΘΑΕ γωνία τῇ ὑπὸ ΑΘΕ; ὥστε καὶ ἡ ΘΕ εὐθεῖα
τῇ ΕΑ, τουτέστι τῇ ΑΒ ἐστιν ἴση. καὶ ἐπεὶ ἴση
ἐστὶν ἡ ΒΑ εὐθεῖα τῇ ΑΕ, ἴση ἐστὶ καὶ γωνία ἡ ὑπὸ
ΑΒΕ τῇ ὑπὸ ΑΕΒ. ἀλλὰ ἡ ὑπὸ ΑΒΕ τῇ ὑπὸ ΒΑΘ
ἐδείθθη ἴση; καὶ ἡ ὑπὸ ΒΕΑ ἄρα τῇ ὑπὸ ΒΑΘ
ἐστιν ἴση. καὶ κοινὴ τῶν δύο τριγώνων τοῦ τε ΑΒΕ καὶ
τοῦ ΑΒΘ ἐστιν ἡ ὑπὸ ΑΒΕ; λοιπὴ ἄρα ἡ ὑπὸ ΒΑΕ
γωνία λοιπῇ τῇ ὑπὸ ΑΘΒ ἐστιν ἴση; ἰσογώνιον ἄρα
ἐστὶ τὸ ΑΒΕ τρίγωνον τῷ ΑΒΘ τριγώνῳ; ἀνάλογον ἄρα
ἐστὶν ὡς ἡ ΕΒ πρὸς τὴν ΒΑ, οὕτως ἡ ΑΒ πρὸς
τὴν ΒΘ. ἴση δὲ ἡ ΒΑ τῇ ΕΘ; ὡς ἄρα ἡ ΒΕ πρὸς τὴν
ΕΘ, οὕτως ἡ ΕΘ πρὸς τὴν ΘΒ. μείζων δὲ ἡ ΒΕ
τῆς ΕΘ; μείζων ἄρα καὶ ἡ ΕΘ τῆς ΘΒ. ἡ ΒΕ
ἅρα ἄκρον καὶ μέσον λόγον τέτμ ηται κατὰ τὸ Θ, καὶ τὸ
μεῖζον τμῆμα τὸ ΘΕ ἴσον ἐστὶ τῇ τοῦ πενταγώνου
πλευρᾷ. ὁμοίως δὴ  δείχομεν, ὅτι καὶ ἡ ΑΓ
ἄκρον καὶ μέσον λόγον τέτμ ηται κατὰ τὸ Θ, καὶ τὸ μεῖζον α ὐτῆς
τμῆμα ἡ ΓΘ ἴσον ἐστὶ τῇ τοῦ πενταγώνου πλευρᾷ; ὅπερ
ἔδει δεῖχαι.}}

\ParallelRText{
\begin{center}
{\large Proposition 8}
\end{center}

If straight-lines subtend two consecutive angles of an equilateral and equiangular pentagon then they cut one another in extreme and
mean ratio, and their greater pieces are equal to  the sides of the pentagon.

For let the two straight-lines, $AC$ and $BE$, cutting one another
at point $H$, have subtended  two consecutive angles, at $A$ and $B$ (respectively), of
the equilateral and equiangular pentagon $ABCDE$. I say that
each of them has been cut in extreme and mean ratio at point $H$, and
that their greater pieces are equal to the sides of the pentagon.

\epsfysize=2.in
\centerline{\epsffile{Book13/fig08e.eps}}

For let the circle $ABCDE$ be circumscribed about pentagon
$ABCDE$ [Prop.~4.14]. And since the two straight-lines $EA$ and
$AB$ are equal to the two (straight-lines) $AB$ and $BC$ (respectively),
and they contain equal angles, the base $BE$ is thus equal to the base
$AC$, and triangle $ABE$ is equal to triangle $ABC$, and the remaining
angles will be equal to the remaining angles, respectively,
which the equal sides subtend [Prop.~1.4]. Thus, angle $BAC$ is equal
to (angle) $ABE$. Thus, (angle) $AHE$ (is) double (angle) $BAH$ [Prop.~1.32]. And $EAC$ is also double $BAC$, inasmuch as
circumference $EDC$ is also double circumference $CB$ [Props.~3.28, 6.33].
Thus, angle $HAE$ (is) equal to (angle) $AHE$. Hence, straight-line
$HE$ is also equal to (straight-line) $EA$---that is to say, to (straight-line)
$AB$ [Prop.~1.6]. And since straight-line $BA$ is equal to $AE$,
angle $ABE$ is also equal to $AEB$ [Prop.~1.5]. But, $ABE$
was shown (to be) equal to $BAH$. Thus, $BEA$ is also equal to
$BAH$. And (angle) $ABE$ is common to the two triangles
$ABE$ and $ABH$. Thus, the remaining angle $BAE$ is equal to
the remaining (angle) $AHB$ [Prop.~1.32]. Thus, triangle $ABE$
is equiangular to triangle $ABH$. Thus, proportionally, as $EB$ is
to $BA$, so $AB$ (is) to $BH$ [Prop.~6.4].  And $BA$ (is) equal to
$EH$. Thus, as $BE$ (is) to $EH$, so $EH$ (is) to $HB$. And
$BE$ (is) greater than $EH$.  $EH$ (is) thus also greater than $HB$
[Prop.~5.14]. Thus, $BE$ has been cut in extreme and mean ratio at
$H$, and the greater piece $HE$ is equal to the side of the pentagon. 
So, similarly, we can show that $AC$ has also been cut in extreme and
mean ratio at $H$, and that  its greater piece $CH$ is equal to the side
of the pentagon. (Which is) the very thing it was required to show.}
\end{Parallel}

%%%%
%13.9
%%%%
\pdfbookmark[1]{Proposition 13.9}{pdf13.9}
\begin{Parallel}{}{}
\ParallelLText{
\begin{center}
{\large \ggn{9}.}
\end{center}\vspace*{-7pt}

{\greekfont Ἐὰν ἡ τοῦ ἑχαγώνου πλευρὰ καὶ ἡ τοῦ δεκαγώνου τῶν
εἰς τὸν α ὐτὸν κύκλον ἐγγραφομένων συντεθῶσιν, ἡ
ὅλη εὐθεῖα ἄκρον καὶ μέσον λόγον τέτμ ηται, καὶ τὸ
μεῖχον α ὐτῆς τμῆμά ἐστιν ἡ τοῦ ἑχαγώνου πλευρά.}\\

{\greekfont Ἔστω κύκλος ὁ ΑΒΓ, καὶ τῶν εἰς τὸν ΑΒΓ κύκλον ἐγγραφομένων
σθημάτων, δεκαγώνου μὲν ἔστω πλευρὰ ἡ ΒΓ, ἑχαγώνου δὲ
ἡ ΓΔ, καὶ ἔστωσαν ἐπ
εὐθείας; λέγω, ὅτι ἡ ὅλη εὐθεῖα ἡ ΒΔ ἄκρον καὶ μέσον
λόγον τέτμ ηται, καὶ τὸ μεῖζον α ὐτῆς τμῆμά
ἐστιν ἡ ΓΔ.}

\epsfysize=2.25in
\centerline{\epsffile{Book13/fig09g.eps}}

{\greekfont Εἰλήφθω γὰρ τὸ κέντρον τοῦ κύκλου τὸ Ε σημεῖον,
καὶ ἐπεζεύθθωσαν αἱ ΕΒ, ΕΓ, ΕΔ, καὶ διήθθω ἡ ΒΕ ἐπὶ
τὸ Α. ἐπεὶ δεκαγώνου ἰσοπλεύρον πλευρά ἐστιν ἡ ΒΓ,
πενταπλασίων ἄρα ἡ ΑΓΒ περιφέρεια
τῆς ΒΓ περιφερείας; τετραπλασίων ἄρα ἡ ΑΓ περιφέρεια τῆς
ΓΒ. ὡς δὲ ἡ ΑΓ περιφέρεια πρὸς τὴν
ΓΒ, οὕτως ἡ ὑπὸ ΑΕΓ γωνία πρὸς τὴν ὑπὸ ΓΕΒ;
τετραπλασίων ἄρα ἡ ὑπὸ ΑΕΓ τῆς ὑπὸ ΓΕΒ. καὶ
ἐπεὶ ἴση ἡ ὑπὸ ΕΒΓ γωνία τῇ ὑπὸ ΕΓΒ, ἡ
ἄρα ὑπὸ ΑΕΓ γωνία διπλασία ἐστὶ τῆς ὑπὸ ΕΓΒ. καὶ
ἐπεὶ ἴση ἐστὶν ἡ ΕΓ εὐθεῖα τῇ ΓΔ; ἑκατέρα γὰρ
α ὐτῶν ἴση ἐστὶ τῇ τοῦ ἑχαγώνου πλευρᾷ τοῦ εἰς
τὸν ΑΒΓ κύκλον [ἐγγραφομένου]; ἴση ἐστὶ καὶ ἡ
ὑπὸ ΓΕΔ γωνία τῇ ὑπὸ ΓΔΕ γωνίᾳ; διπλασία ἄρα ἡ
ὑπὸ ΕΓΒ γωνία τῆς ὑπὸ ΕΔΓ. ἀλλὰ τῆς ὑπὸ ΕΓΒ
διπλασία ἐδείθθη ἡ ὑπὸ ΑΕΓ; τετραπλασία ἄρα ἡ ὑπὸ
ΑΕΓ τῆς ὑπὸ ΕΔΓ. ἐδείθθη δὲ καὶ τῆς ὑπὸ ΒΕΓ τετραπλασία
ἡ ὑπὸ ΑΕΓ; ἴση ἄρα ἡ ὑπὸ ΕΔΓ τῇ ὑπὸ ΒΕΓ. κοινὴ
δὲ τῶν δύο τριγώνων, τοῦ τε
ΒΕΓ καὶ τοῦ ΒΕΔ, ἡ ὑπὸ ΕΒΔ γωνία;
καὶ λοιπὴ ἄρα ἡ ὑπὸ ΒΕΔ τῇ ὑπὸ ΕΓΒ ἐστιν ἴση;
ἰσογώνιον ἄρα ἐστὶ τὸ ΕΒΔ τρίγωνον τῷ ΕΒΓ τριγώνῳ.
ἀνάλογον ἄρα ἐστὶν ὡς ἡ ΔΒ πρὸς τὴν ΒΕ, οὕτως
ἡ ΕΒ πρὸς τὴν ΒΓ. ἴση δὲ ἡ ΕΒ τῇ ΓΔ. ἔστιν
ἄρα ὡς ἡ ΒΔ πρὸς τὴν ΔΓ, οὕτως ἡ ΔΓ πρὸς τὴν ΓΒ.
μείζων δὲ ἡ ΒΔ τῆς ΔΓ; μείζων ἄρα καὶ ἡ ΔΓ
τῆς ΓΒ. ἡ ΒΔ ἄρα εὐθεῖα ἄκρον καὶ μέσον λόγον
τέτμ ηται [κατὰ τὸ Γ], καὶ τὸ μεῖζον τμῆμα α ὐτῆς
ἐστιν ἡ ΔΓ; ὅπερ ἔδει δεῖχαι.}}

\ParallelRText{
\begin{center}
{\large Proposition 9}
\end{center}

If the side of a hexagon and  of a decagon inscribed in the same circle
are added together then the whole straight-line has been cut in extreme and
mean ratio (at the junction point), and its greater piece is the side of the hexagon.$^\dag$

Let $ABC$ be a circle. And of the figures inscribed in circle $ABC$, 
let $BC$ be the side of a decagon, and $CD$ (the side) of a hexagon.
And let them be (laid down) straight-on (to one another). I say that
the whole straight-line $BD$ has been cut in extreme and mean
ratio (at $C$), and that   $CD$ is its greater piece. \\

\epsfysize=2.25in
\centerline{\epsffile{Book13/fig09e.eps}}

For let the center of the circle, point $E$, be found [Prop.~3.1],
and let $EB$, $EC$, and $ED$ be joined, and let $BE$ be
drawn across to $A$. Since $BC$ is a side on an equilateral decagon,
circumference $ACB$ (is) thus five times circumference $BC$. Thus,
circumference $AC$ (is) four times $CB$. And as circumference $AC$
(is) to $CB$, so angle $AEC$ (is) to $CEB$ [Prop.~6.33]. Thus,
(angle) $AEC$ (is) four times $CEB$. And since angle $EBC$
(is) equal to $ECB$ [Prop.~1.5], angle $AEC$ is thus double $ECB$ [Prop.~1.32]. 
And since straight-line $EC$ is equal to $CD$---for each of them is equal
to the side of the hexagon [inscribed] in circle $ABC$ [Prop.~4.15~corr.]---
angle $CED$ is also equal to angle $CDE$ [Prop.~1.5]. Thus,
angle $ECB$ (is) double $EDC$ [Prop.~1.32]. But,
$AEC$ was shown (to be) double $ECB$. Thus, $AEC$ (is) four times
$EDC$. And $AEC$ was also shown (to be) four times $BEC$. 
Thus, $EDC$ (is) equal to $BEC$. And angle $EBD$ (is) common
to the two triangles $BEC$ and $BED$. Thus, the remaining (angle)
$BED$ is equal to the (remaining angle) $ECB$ [Prop.~1.32]. 
Thus, triangle $EBD$ is equiangular to triangle $EBC$. Thus,
proportionally, as $DB$ is to $BE$, so $EB$ (is) to $BC$ [Prop.~6.4].
And $EB$ (is) equal to $CD$. Thus, as $BD$ is to $DC$, so $DC$
(is) to $CB$. And $BD$ (is) greater than $DC$. Thus,
$DC$ (is) also greater than $CB$ [Prop.~5.14]. Thus, the
straight-line $BD$ has been cut in extreme and mean ratio [at $C$],
and $DC$  is its greater piece. (Which is), the very thing it was required to show.}
\end{Parallel}{\footnotesize\noindent$^\dag$ If the circle is of unit radius
then the side of the hexagon is 1, whereas the side of the decagon is $(1/2)\,(\sqrt{5}-1)$.}

%%%%
%13.10
%%%%
\pdfbookmark[1]{Proposition 13.10}{pdf13.10}
\begin{Parallel}{}{}
\ParallelLText{
\begin{center}
{\large \ggn{10}.}
\end{center}\vspace*{-7pt}

{\greekfont Ἐὰν εἰς κύκλον πεντάγωνον ἰσόπλευρον ἐγγραφῇ, ἡ τοῦ
πενταγώνου πλευρὰ δύναται τήν τε τοῦ ἑχαγώνου καὶ τὴν τοῦ
δεκαγώνου τῶν εἰς τὸν α ὐτὸν κύκλον ἐγγραφομένων.}\\

\epsfysize=2.5in
\centerline{\epsffile{Book13/fig10g.eps}}

{\greekfont Ἔστω κύκλος ὁ ΑΒΓΔΕ, καὶ εἰς τὸ ΑΒΓΔΕ κύκλον
πεντάγωνον ἰσόπλευρον ἐγγεγράφθω τὸ ΑΒΓΔΕ. λέγω,
ὅτι ἡ τοῦ ΑΒΓΔΕ πενταγώνου πλευρὰ δύναται τήν τε τοῦ
ἑχαγώνου καὶ τὴν τοῦ δεκαγώνου πλευρὰν τῶν εἰς
τὸν ΑΒΓΔΕ κύκλον ἐγγραφομένων.}

{\greekfont Εἰλήφθω γὰρ τὸ κέντρον τοῦ κύκλου τὸ Ζ σημεὶον, καὶ
ἐπιζευθθεῖσα ἡ ΑΖ διήθθω ἐπὶ τὸ Η σημεῖον, καὶ ἐπεζεύθθω
ἡ ΖΒ, καὶ ἀπὸ τοῦ Ζ ἐπὶ τὴν ΑΒ κάθετος ἤθθω ἡ ΖΘ,
καὶ διήθθω ἐπὶ τὸ Κ, καὶ ἐπεζεύθθωσαν αἱ ΑΚ, ΚΒ, 
καὶ πάλιν ἀπὸ τοῦ Ζ ἐπὶ τὴν ΑΚ κάθετος ἤθθω ἡ ΖΛ, καὶ
διήθθω ἐπὶ τὸ Μ, καὶ ἐπεζεύθθω ἡ ΚΝ.}

{\greekfont Ἐπεὶ ἴση ἐστὶν
ἡ ΑΒΓΗ περιφέρεια τῇ ΑΕΔΗ περιφερείᾳ, ὧν ἡ ΑΒΓ τῇ
ΑΕΔ ἐστιν ἴση, λοιπὴ ἄρα ἡ ΓΗ περιφέρεια λοιπῇ τῇ ΗΔ ἐστιν
ἴση. πενταγώνου δὲ ἡ ΓΔ; δεκαγώνου ἄρα ἡ ΓΗ. καὶ ἐπεὶ
ἴση ἐστὶν ἡ ΖΑ τῇ ΖΒ, καὶ κάθετος ἡ ΖΘ, ἴση ἄρα καὶ ἡ
ὑπὸ ΑΖΚ γωνία τῇ ὑπὸ ΚΖΒ. ὥστε καὶ περιφέρεια ἡ ΑΚ τῇ
ΚΒ ἐστιν ἴση; διπλῆ ἄρα ἡ ΑΒ περιφέρεια τῆς ΒΚ περιφερείας;
δεκαγώνου ἄρα πλευρά ἐστιν ἡ ΑΚ εὐθεῖα. διὰ τὰ α ὐτὰ
δὴ καὶ ἡ ΑΚ τῆς ΚΜ ἐστι διπλῆ. καὶ ἐπεὶ διπλῆ ἐστιν
ἡ ΑΒ περιφέρεια τῆς ΒΚ περιφερείας, ἴση δὲ ἡ ΓΔ περιφέρεια
τῇ ΑΒ περιφερείᾳ, διπλῆ ἄρα καὶ ἡ ΓΔ περιφέρεια
τῆς ΒΚ περιφερείας. ἔστι δὲ ἡ ΓΔ περιφέρεια καὶ τῆς ΓΗ διπλῆ; ἴση
ἄρα ἡ ΓΗ περιφέρεια τῇ ΒΚ περιφερείᾳ. ἀλλὰ ἡ ΒΚ τῆς ΚΜ ἐστι
διπλῆ, ἐπεὶ καὶ ἡ ΚΑ; καὶ ἡ ΓΗ ἄρα τῆς ΚΜ ἐστι διπλῆ.
ἀλλὰ μὴν καὶ ἡ ΓΒ περιφέρεια τῆς ΒΚ περιφερείας ἐστὶ διπλῆ;
ἴση γὰρ ἡ ΓΒ περιφέρεια τῇ ΒΑ. καὶ ὅλη ἄρα ἡ ΗΒ περιφέρεια
τῆς ΒΜ ἐστι διπλῆ; ὥστε καὶ γωνία ἡ ὑπὸ ΗΖΒ γωνίας τῆς
ὑπὸ ΒΖΜ [ἐστι] διπλῆ. ἔστι δὲ ἡ ὑπὸ ΗΖΒ καὶ τῆς ὑπὸ
ΖΑΒ διπλῆ; ἴση γὰρ ἡ ὑπὸ ΖΑΒ τῇ ὑπὸ ΑΒΖ. καὶ
ἡ ὑπὸ ΒΖΝ ἄρα τῇ ὑπὸ ΖΑΒ ἐστιν ἴση. κοινὴ δὲ τῶν
δύο τριγώνων, τοῦ τε ΑΒΖ καὶ τοῦ ΒΖΝ, ἡ ὑπὸ ΑΒΖ γωνία;
λοιπὴ ἄρα ἡ ὑπὸ ΑΖΒ λοιπῇ τῇ ὑπὸ ΒΝΖ ἐστιν ἴση;
ἴσογώνιον ἄρα ἐστὶ τὸ ΑΒΖ τρίγωνον τῷ ΒΖΝ τριγώνῳ. ἀνάλογον
ἄρα ἐστὶν ὡς ἡ ΑΒ εὐθεῖα πρὸς τὴν ΒΖ, οὕτως
ἡ ΖΒ πρὸς τὴν ΒΝ; τὸ ἄρα ὑπὸ τῶν ΑΒΝ ἴσον ἐστὶ τῷ
ἀπὸ ΒΖ. πάλιν ἐπεὶ ἴση ἐστὶν ἡ ΑΛ τῇ ΛΚ, κοινὴ δὲ καὶ
πρὸς ὀρθὰς ἡ ΛΝ, βάσις ἄρα ἡ ΚΝ βάσει τῇ ΑΝ ἐστιν
ἴση; καὶ γωνία ἄρα ἡ ὑπὸ ΛΚΝ γωνίᾳ τῇ ὑπὸ ΛΑΝ
ἐστιν ἴση. ἀλλὰ ἡ ὑπὸ ΛΑΝ τῇ ὑπὸ ΚΒΝ ἐστιν ἴση;
καὶ ἡ ὑπὸ ΛΚΝ ἄρα τῇ ὑπὸ ΚΒΝ ἐστιν ἴση. καὶ κοινὴ
τῶν δύο τριγώνων τοῦ τε ΑΚΒ καὶ τοῦ ΑΚΝ ἡ πρὸς τῷ Α.
λοιπὴ ἄρα ἡ ὑπὸ ΑΚΒ λοιπῇ τῇ ὑπὸ ΚΝΑ ἐστιν ἴση;
ἰσογώνιον ἄρα ἐστὶ τὸ ΚΒΑ τρίγωνον τῷ ΚΝΑ τριγώνῳ.
ἀνάλογον ἄρα ἐστὶν ὡς ἡ ΒΑ εὐθεῖα πρὸς τὴν ΑΚ, οὕτως
ἡ ΚΑ πρὸς τὴν ΑΝ; τὸ ἄρα ὑπὸ τῶν ΒΑΝ ἴσον ἐστὶ τῷ
ἀπὸ τῆς ΑΚ. ἐδείθθη δὲ καὶ τὸ ὑπὸ τῶν ΑΒΝ ἴσον τῷ ἀπὸ
τῆς ΒΖ; τὸ ἄρα ὑπὸ τῶν ΑΒΝ μετὰ τοῦ  ὑπὸ ΒΑΝ, ὅπερ
ἐστὶ τὸ ἀπὸ τῆς ΒΑ, ἴσον ἐστὶ τῷ ἀπὸ τῆς ΒΖ μετὰ τοῦ
ἀπὸ τῆς ΑΚ. καί ἐστιν ἡ μὲν ΒΑ πενταγώνου πλευρά, ἡ δὲ
ΒΖ ἑχαγώνου, ἡ δὲ ΑΚ δεκαγώνου.}

{\greekfont Ἡ ἄρα τοῦ πενταγώνου πλευρὰ δύναται τήν τε τοῦ ἑχαγώνου
καὶ τὴν τοῦ δεκαγώνου τῶν εἰς τὸν α ὐτὸν κύκλον ἐγγραφομένων;
ὅπερ ἔδει δεῖχαι.}}

\ParallelRText{
\begin{center}
{\large Proposition 10}
\end{center}

If an equilateral pentagon is inscribed in a circle then the square on the
side of the pentagon is (equal to) the (sum of the squares) on the (sides) of
the hexagon and of the decagon inscribed in the same circle.$^\dag$

\epsfysize=2.5in
\centerline{\epsffile{Book13/fig10e.eps}}

Let $ABCDE$ be a circle. And let the equilateral pentagon $ABCDE$
be inscribed in circle $ABCDE$. I say that the square on the
side of pentagon $ABCDE$ is the (sum of the squares) on the sides
of the hexagon and of the decagon inscribed in circle $ABCDE$. 

For let the center of the circle, point $F$, be found [Prop.~3.1]. And,
$AF$ being joined, let it be drawn across to point $G$. And
let $FB$ be joined. And let $FH$ be drawn from $F$
perpendicular to $AB$. And let it be drawn across to $K$. 
And let $AK$ and $KB$ be joined.  And, again, let $FL$ have
been drawn from $F$ perpendicular to $AK$. And let it be drawn
across to $M$. And let $KN$ be joined.

Since circumference $ABCG$ is equal to circumference $AEDG$,
of which $ABC$ is equal to $AED$, the remaining circumference
$CG$ is thus equal to the remaining (circumference) $GD$. And $CD$ (is the side) of the pentagon. 
$CG$ (is) thus (the side) of the decagon. And since  $FA$ is equal to $FB$,
and $FH$ is perpendicular (to $AB$), angle $AFK$ (is) thus also equal
to $KFB$ [Props.~1.5, 1.26]. Hence, circumference $AK$
is also equal to $KB$ [Prop.~3.26]. Thus, circumference $AB$ (is)
double circumference $BK$. Thus, straight-line $AK$ is the side of the decagon. 
So, for the same (reasons, circumference)  $AK$ is also double $KM$. And since
circumference $AB$ is double circumference $BK$, and
circumference $CD$ (is) equal to circumference $AB$, circumference
$CD$ (is) thus also double circumference $BK$. And
circumference $CD$ is also double $CG$. Thus, circumference
$CG$ (is) equal to circumference $BK$. But, $BK$
is double $KM$, since $KA$ (is) also (double $KM$). Thus,
(circumference) $CG$ is also double $KM$. But, indeed, circumference $CB$ is also
double circumference $BK$. For circumference $CB$ (is) equal to $BA$. 
Thus, the whole circumference $GB$ is also double $BM$. 
Hence, angle $GFB$ [is] also double  angle $BFM$ [Prop.~6.33].
And $GFB$ (is) also double $FAB$. For $FAB$ (is) equal to $ABF$. 
Thus, $BFN$ is also equal to $FAB$. And angle $ABF$ (is) common
to the two triangles $ABF$ and $BFN$. Thus, the
remaining (angle) $AFB$ is equal to the remaining (angle) $BNF$
[Prop.~1.32]. Thus, triangle $ABF$ is equiangular to triangle
$BFN$. Thus, proportionally, as straight-line $AB$ (is) to
$BF$, so $FB$ (is) to $BN$ [Prop.~6.4]. Thus, the (rectangle contained)
by $ABN$ is equal to the (square) on $BF$ [Prop.~6.17]. Again,
since $AL$ is equal to $LK$, and $LN$ is common and at right-angles (to $KA$), 
base $KN$ is thus equal to base $AN$  [Prop.~1.4]. And, thus, angle
$LKN$ is equal to angle $LAN$. But, $LAN$ is equal to
$KBN$ [Props.~3.29, 1.5]. Thus, $LKN$ is also equal to $KBN$. 
And the (angle) at $A$ (is) common to the two triangles $AKB$ and
$AKN$. Thus, the remaining (angle) $AKB$ is equal to the
remaining (angle) $KNA$ [Prop.~1.32]. Thus, triangle $KBA$ is
equiangular to triangle $KNA$. Thus, proportionally, as straight-line
$BA$ is to $AK$, so $KA$ (is) to $AN$ [Prop.~6.4]. Thus, the
(rectangle contained) by $BAN$ is equal to the (square) on $AK$ [Prop.~6.17]. 
And  the (rectangle contained) by $ABN$ was also shown (to be)
equal to the (square) on $BF$. Thus, the (rectangle contained) by $ABN$
plus the (rectangle contained) by $BAN$, which is the (square) on $BA$
[Prop.~2.2],
is equal to the (square) on $BF$ plus the (square) on $AK$. 
And $BA$ is the side of the pentagon, and $BF$ (the side) of the hexagon
[Prop.~4.15~corr.], and $AK$ (the side) of the decagon.

Thus, the square on the side of the pentagon  (inscribed in a circle) is (equal to) the (sum of the squares)
on the (sides) of the hexagon and of  the decagon inscribed in the
same circle.}
\end{Parallel}
{\footnotesize\noindent$^\dag$ If the circle is of unit radius then the side
of the pentagon is $(1/2)\,\sqrt{10-2\,\sqrt{5}}$.}

%%%%
%13.11
%%%%
\pdfbookmark[1]{Proposition 13.11}{pdf13.11}
\begin{Parallel}{}{}
\ParallelLText{
\begin{center}
{\large \ggn{11}.}
\end{center}\vspace*{-7pt}

{\greekfont Ἐὰν εἰς κύκλον ῥητὴν ἔθοντα τὴν διάμετρον πεντάγω-νον
ἰσόπλευρον ἐγγραφῇ, ἡ τοῦ πενταγώνου πλευρὰ ἄλογός
ἐστιν ἡ καλουμένη ἐλάσσων.}

\epsfysize=2.5in
\centerline{\epsffile{Book13/fig11g.eps}}

{\greekfont Εἰς γὰρ κύκλον τὸν ΑΒΓΔΕ ῥητὴν ἔθοντα τὴν δίαμετρον
πεντάγωνον ἰσόπλευρον  ἐγγεγράφθω τὸ ΑΒΓΔΕ; λέγω,
ὅτι ἡ τοῦ [ΑΒΓΔΕ] πενταγώνου πλευρὰ ἄλογός ἐστιν
ἡ καλουμένη ἐλάσσων.}

{\greekfont Εἰλήφθω γὰρ τὸ κέντρον τοῦ κύκλου τὸ Ζ σημεῖον, καὶ
ἐπεζεύθθωσαν αἱ ΑΖ, ΖΒ καὶ διήθθωσαν ἐπὶ τὰ Η, Θ σημεῖα,
καὶ ἐπεζεύθθω ἡ ΑΓ, καὶ κείσθω τῆς ΑΖ τέταρτον μέρος
ἡ ΖΚ. ῥητὴ δὲ ἡ ΑΖ; ῥητὴ ἄρα καὶ ἡ ΖΚ. ἔστι
δὲ καὶ ἡ ΒΖ ῥητή; ὅλη ἄρα ἡ ΒΚ ῥητή ἐστιν. καὶ ἐπεὶ
ἴση ἐστὶν ἡ ΑΓΗ περιφέρεια τῇ ΑΔΗ περιφερείᾳ, ὧν ἡ ΑΒΓ τῇ
ΑΕΔ ἐστιν ἴση, λοιπὴ ἄρα ἡ ΓΗ λοιπῇ τῇ ΗΔ ἐστιν ἴση.
καὶ ἐὰν ἐπιζεύχωμεν τὴν ΑΔ, συνάγονται ὀρθαὶ αἱ πρὸς τῷ
Λ γωνίαι, καὶ διπλῆ ἡ ΓΔ τῆς ΓΛ. διὰ τὰ α ὐτὰ δὴ καὶ αἱ
πρὸς τῷ Μ ὀρθαί εἰσιν, καὶ διπλῆ ἡ ΑΓ τῆς ΓΜ. ἐπεὶ
οὖν ἴση ἐστὶν ἡ ὑπὸ ΑΛΓ γωνία τῇ ὑπὸ ΑΜΖ, κοινὴ δὲ
τῶν δύο τριγώνων τοῦ τε ΑΓΛ καὶ τοῦ ΑΜΖ ἡ ὑπὸ ΛΑΓ,
λοιπὴ ἄρα ἡ ὑπὸ ΑΓΛ λοιπῇ τῇ ὑπὸ ΜΖΑ ἐστιν ἴση;
ἰσογώνιον ἄρα ἐστὶ τὸ ΑΓΛ τρίγωνον τῷ ΑΜΖ τριγώνῳ;
ἀνάλογον ἄρα ἐστὶν ὡς ἡ ΛΓ πρὸς ΓΑ, οὕτως ἡ 
ΜΖ πρὸς ΖΑ; καὶ τῶν ἡγουμένων τὰ διπλάσια; ὡς ἄρα
ἡ τῆς ΛΓ διπλῆ πρὸς τὴν ΓΑ, οὕτως ἡ τῆς ΜΖ
διπλῆ πρὸς τὴν ΖΑ. ὡς δὲ ἡ τῆς ΜΖ διπλῆ πρὸς τὴν ΖΑ,
οὕτως ἡ ΜΖ πρὸς τὴν ἡμίσειαν τῆς ΖΑ; καὶ ὡς ἄρα
ἡ τῆς ΛΓ διπλῆ πρὸς τὴν ΓΑ, οὕτως ἡ ΜΖ πρὸς τὴν
ἡμίσειαν τῆς ΖΑ; καὶ 
τῶν ἑπομένων τὰ ἡμίσεα; ὡς ἄρα ἡ τῆς ΛΓ διπλῆ πρὸς τὴν
ἡμίσειαν τῆς ΓΑ,  οὕτως ἡ ΜΖ πρὸς τὸ τέτατρον τῆς ΖΑ. καί
ἐστι τῆς μὲν ΛΓ διπλῆ ἡ ΔΓ, τῆς δὲ ΓΑ ἡμίσεια ἡ ΓΜ,
τῆς δὲ ΖΑ τέτατρον μέρος ἡ ΖΚ; ἔστιν ἄρα ὡς ἡ ΔΓ πρὸς
τὴν ΓΜ, οὕτως ἡ ΜΖ πρὸς τὴν ΖΚ. συνθέντι καὶ ὡς συναμφότερος
ἡ ΔΓΜ πρὸς τὴν ΓΜ, οὕτως ἡ ΜΚ πρὸς ΚΖ; καὶ ὡς ἄρα
τὸ ἀπὸ συναμφοτέρου τῆς ΔΓΜ πρὸς τὸ ἀπὸ ΓΜ, οὕτως τὸ
ἀπὸ ΜΚ πρὸς τὸ ἀπὸ ΚΖ. καὶ ἐπεὶ τῆς ὑπὸ δύο
πλευρὰς τοῦ πενταγώνου ὑποτεινούσης, οἷον τῆς ΑΓ, ἄκρον
καὶ μέσον λόγον τεμνομένης τὸ μεῖζον τμῆμα ἴσον ἐστὶ
τῇ τοῦ πενταγώνου πλευρᾷ, τουτέστι τῇ ΔΓ, τὸ δὲ μεῖζον
τμῆμα προσλαβὸν τὴν ἡμίσειαν τῆς ὅλῆς πενταπλάσιον
δύναται τοῦ ἀπὸ τῆς ἡμισείας τῆς ὅλης, καί ἐστιν ὅλης τῆς
ΑΓ ἡμίσεια ἡ ΓΜ, τὸ ἄρα ἀπὸ τῆς ΔΓΜ ὡς μιᾶς
πενταπλάσιόν ἐστι τοῦ ἀπὸ τῆς ΓΜ. ὡς δὲ τὸ ἀπὸ τῆς
ΔΓΜ ὡς μιᾶς πρὸς τὸ ἀπὸ τῆς ΓΜ, οὕτως ἐδείθθη τὸ
ἀπὸ τῆς ΜΚ πρὸς τὸ ἀπὸ τῆς ΚΖ; πενταπλάσιον ἄρα τὸ ἀπὸ
τῆς ΜΚ τοῦ ἀπὸ τῆς ΚΖ. ῥητὸν δὲ τὸ ἀπὸ τῆς ΚΖ;
ῥητὴ γὰρ ἡ  διάμετρος; ῥητὸν ἄρα καὶ τὸ ἀπὸ τῆς ΜΚ;
ῥητὴ ἄρα  ἐστὶν ἡ ΜΚ [δυνάμει μόνον]. καὶ ἐπεὶ
τετραπλασία ἐστὶν ἡ ΒΖ τῆς ΖΚ, πενταπλασία ἄρα ἐστὶν
ἡ ΒΚ τῆς  ΚΖ;
εἰκοσιπενταπλάσιον ἄρα τὸ ἀπὸ τῆς ΒΚ τοῦ ἀπὸ τῆς ΚΖ.
πενταπλάσιον δὲ τὸ ἀπὸ τῆς ΜΚ τοῦ ἀπὸ τῆς ΚΖ; πενταπλάσιον
ἄρα τὸ ἀπὸ τῆς ΒΚ τοῦ ἀπὸ τῆς ΚΜ; τὸ ἄρα ἀπὸ τῆς
ΒΚ πρὸς τὸ ἀπὸ ΚΜ λόγον οὐκ ἔχει, ὃν τετράγωνος ἀριθμὸς
πρὸς τετράγωνον ἀριθμόν; ἀσύμμετρος ἄρα ἐστὶν ἡ ΒΚ τῇ ΚΜ
μήκει. καί ἐστι ῥητὴ ἑκατέρα α ὐτῶν. αἱ ΒΚ, ΚΜ ἄρα
ῥηταί εἰσι δυνάμει μόνον σύμμετροι. ἐὰν δὲ ἀπὸ ῥητῆς
ῥητὴ ἀφαιρεθῇ δυνάμει μόνον σύμμετρος οὖσα τῇ
ὅλῃ, ἡ λοιπὴ ἄλογός ἐστιν ἀποτομή; ἀποτομὴ ἄρα
ἐστὶν ἡ ΜΒ, προσαρμόζουσα δὲ α ὐτῇ ἡ ΜΚ. λέγω δή,
ὅτι καὶ τετάρτη. ᾧ δὴ μεῖζόν ἐστι τὸ ἀπὸ τῆς ΒΚ
τοῦ ἀπὸ τῆς ΚΜ, ἐκείνῳ ἴσον ἔστω τὸ ἀπὸ τῆς
Ν; ἡ ΒΚ ἄρα τῆς ΚΜ
μεῖζον δύναται τῇ Ν. καὶ ἐπεὶ
σύμμετρός ἐστιν ἡ ΚΖ τῇ ΖΒ, καὶ συνθέντι σύμμετρός ἐστι ἡ ΚΒ τῇ ΖΒ. ἀλλὰ ἡ ΒΖ τῇ ΒΘ σύμμετρός ἐστιν; καὶ ἡ ΒΚ ἄρα τῇ
ΒΘ σύμμετρός ἐστιν. καὶ ἐπεὶ πενταπλάσιόν ἐστι τὸ ἀπὸ τῆς
ΒΚ τοῦ ἀπὸ τῆς ΚΜ, τὸ ἄρα ἀπὸ τῆς ΒΚ πρὸς τὸ ἀπὸ
τῆς ΚΜ λόγον ἔχει, ὃνὸv{ε} πρὸς ἕν. ἀναστρέψαντι ἄρα
τὸ ἀπὸ τῆς ΒΚ πρὸς τὸ ἀπὸ τῆς Ν λόγον ἔχει, ὃν ὸv{ε}
πρὸς ὸv{δ}, οὐθ ὃν τετράγωνος πρὸς τετράγωνον; ἀσύμμετρος
ἄρα ἐστὶν ἡ ΒΚ τῇ Ν; ἡ ΒΚ ἄρα τῆς ΚΜ μεῖζον
δύναται τῷ ἀπὸ ἀσυμμέτρου ἑαυτῇ. ἐπεὶ οὖν ὅλη
ἡ ΒΚ τῆς προσαρμοζούσης τῆς ΚΜ μεῖζον δύναται τῷ ἀπὸ
ἀσυμμέτρου ἑαυτῇ, καὶ ὅλη ἡ ΒΚ σύμμετρός ἐστι τῇ ἐκκειμένῃ
ῥητῇ τῇ ΒΘ, ἀποτομ ὴ ἄρα τετάρτη ἐστὶν ἡ ΜΒ. τὸ δὲ ὑπὸ ῥητῆς
καὶ ἀποτομ ῆς τετάρτης περιεχόμενον ὀρθογώνιον ἄλογόν ἐστιν,
καὶ ἡ δυναμένη α ὐτὸ ἄλογός ἐστιν, καλεῖται δὲ ἐλάττων. δύναται
δὲ τὸ ὑπὸ τῶν ΘΒΜ ἡ ΑΒ διὰ τὸ ἐπιζευγνυμένης
τῆς ΑΘ ἰσογώνιον γίνεσθαι τὸ ΑΒΘ τρίγωνον τῷ ΑΒΜ τριγώνῳ
καὶ εἶναι ὡς τὴν ΘΒ πρὸς τὴν ΒΑ, οὕτως τὴν ΑΒ πρὸς
τὴν ΒΜ.}

{\greekfont Ἡ ἄρα ΑΒ τοῦ πενταγώνου πλευρὰ ἄλογός ἐστιν ἡ καλουμένη
ἐλάττων; ὅπερ ἔδει δεῖχαι.}}

\ParallelRText{
\begin{center}
{\large Proposition 11}
\end{center}

If an equilateral pentagon is inscribed in a circle which has a rational  diameter  then the side of the pentagon is that irrational
(straight-line) called minor.

\epsfysize=2.5in
\centerline{\epsffile{Book13/fig11e.eps}}

For let the equilateral pentagon $ABCDE$ be inscribed  in the circle $ABCDE$ which has a rational diameter. I say that the side of pentagon
[$ABCDE$] is that irrational (straight-line) called minor.

For let the center of the circle, point $F$, be found [Prop.~3.1].
And let $AF$ and $FB$ be joined. And let them be drawn
across to points $G$ and $H$ (respectively). And let $AC$ be joined.
And let $FK$ made (equal) to the fourth part of $AF$. And $AF$ (is) rational. $FK$ (is) thus also rational. And $BF$ is also rational. Thus,
the whole of $BK$ is rational. And since circumference $ACG$ is equal
to circumference $ADG$, of which $ABC$ is equal to $AED$, the
remainder $CG$ is thus equal to the remainder $GD$.
And if we join $AD$ then the angles at $L$ are inferred (to be) right-angles, 
and $CD$ (is inferred to be) double $CL$ [Prop.~1.4].  So, for the same (reasons), 
the (angles) at $M$ are also right-angles, and $AC$ (is) double $CM$. Therefore,
since angle $ALC$ (is) equal to $AMF$, and (angle) $LAC$ (is) common
to the two triangles $ACL$ and $AMF$, the remaining (angle) $ACL$ is thus
equal to the remaining (angle) $MFA$ [Prop.~1.32]. Thus, triangle 
$ACL$  is equiangular to triangle $AMF$. Thus, proportionally, 
as $LC$ (is) to $CA$, so $MF$ (is) to $FA$ [Prop.~6.4]. And (we can take) the
doubles of the leading (magnitudes). Thus, as double $LC$ (is) to $CA$,
so double $MF$ (is) to $FA$. And as double $MF$ (is) to $FA$, so
$MF$ (is) to half of $FA$. And, thus, as double $LC$ (is) to $CA$,
so $MF$ (is) to half of $FA$. And (we can take) the halves of the following
(magnitudes). Thus, as double $LC$ (is) to half of $CA$, so $MF$
(is) to the fourth of $FA$. And $DC$ is double $LC$, and  $CM$ half
of $CA$, and $FK$ the fourth part of $FA$.  Thus, as $DC$
is to $CM$, so $MF$ (is) to $FK$. Via composition, as the
sum of $DCM$ ({\em i.e.}, $DC$ and $CM$)
(is) to $CM$, so $MK$ (is) to $KF$ [Prop.~5.18]. And, thus, as the (square) on the
sum of $DCM$ (is) to the (square) on $CM$, so the (square) on $MK$
(is) to the (square) on $KF$. And since the
greater piece of a (straight-line) subtending two sides of a
pentagon, such as $AC$, (which is) cut in extreme and mean ratio is equal to the side of the pentagon [Prop.~13.8]---that is to say,
to $DC$----and the square on the greater piece added to half of the
whole is five times the (square) on half of the whole [Prop.~13.1], 
and $CM$ (is) half of the whole, $AC$, thus the (square) on
$DCM$, (taken) as one, is five times the (square) on $CM$.
And the (square) on $DCM$, (taken) as one,  (is)  to the (square) on $CM$, so the (square) on $MK$ was shown (to be) to the (square) on 
$KF$.  Thus, the (square) on $MK$ (is) five times the (square) on 
$KF$. And the square on $KF$ (is) rational. For the diameter
(is) rational. Thus, the (square) on $MK$ (is) also rational. 
Thus, $MK$ is rational [in square only]. And since $BF$ is four times
$FK$, $BK$ is thus five times $KF$. Thus, the (square) on $BK$
(is) twenty-five times the (square) on $KF$.  And the (square) on $MK$
(is) five times the square on $KF$. Thus, the (square) on $BK$ (is) five
times the (square) on $KM$.  Thus, the (square) on $BK$ does not
have to the (square) on $KM$ the ratio which a square number (has) to
a square number. Thus, $BK$ is incommensurable in length with $KM$ [Prop.~10.9].
And each of them is a rational (straight-line).  Thus, $BK$ and $KM$ are rational
(straight-lines which are) commensurable in square only. And if from a rational (straight-line)
a rational (straight-line) is subtracted, which is commensurable in square
only with the whole, then the remainder is that irrational (straight-line
called) an apotome [Prop.~10.73]. Thus, $MB$ is an apotome, and $MK$ its attachment. So, I say that (it is) also a fourth (apotome). So, let the
(square) on $N$ be (made) equal to that (magnitude) by which the (square)
on $BK$ is greater than the (square) on $KM$. Thus, the
square on $BK$ is greater than the (square) on $KM$ by the (square) on $N$.
And since $KF$ is commensurable (in length) with $FB$ then, via composition, 
$KB$ is also commensurable (in length) with $FB$ [Prop.~10.15]. But, $BF$ is commensurable (in length)
with $BH$. Thus, $BK$ is also commensurable (in length) with $BH$ [Prop.~10.12]. And since the
(square) on $BK$ is five times the (square) on $KM$, the (square) on $BK$
thus has to the (square) on $KM$ the ratio which 5 (has) to one. Thus, via
conversion, the (square) on $BK$ has to the (square) on $N$ the ratio
which 5 (has) to 4 [Prop.~5.19~corr.], which is not (that) of a
square (number) to a square (number). $BK$ is thus incommensurable (in length)
with $N$ [Prop.~10.9]. Thus, the square on $BK$ is greater than
the (square) on $KM$ by the (square) on (some straight-line which is) incommensurable (in length) with ($BK$). Therefore, since the square on the whole,
$BK$, is greater than the (square) on the attachment, $KM$, by the
(square) on (some straight-line which is) incommensurable (in length) with ($BK$),
and the whole, $BK$, is commensurable (in length) with the (previously) laid down
rational (straight-line) $BH$, $MB$ is thus a fourth apotome
[Def.~10.14]. And the rectangle contained by a rational (straight-line)
and a fourth apotome is irrational, and its square-root is that irrational
(straight-line) called minor [Prop.~10.94]. And the square on $AB$
is the rectangle contained by $HBM$, on account of joining $AH$, (so that)
triangle $ABH$ becomes equiangular  with triangle $ABM$ [Prop.~6.8], and (proportionally) as
$HB$ is to $BA$, so $AB$ (is) to $BM$.

Thus, the side $AB$ of the pentagon is that irrational (straight-line)
called minor.$^\dag$ (Which is) the very thing it was required to show.}
\end{Parallel}
{\footnotesize\noindent$^\dag$ If the circle has unit radius then the side of the pentagon is $(1/2)\,\sqrt{10-2\,\sqrt{5}}$. However, this length can
be written in the ``minor'' form (see Prop.~10.94) $(\rho/\sqrt{2})\,\sqrt{1+k/\sqrt{1+k^2}} - (\rho/\sqrt{2})\,\sqrt{1-k/\sqrt{1+k^2}}$, with $\rho=\sqrt{5/2}$ and $k=2$.}

%%%%
%13.12
%%%%
\pdfbookmark[1]{Proposition 13.12}{pdf13.12}
\begin{Parallel}{}{}
\ParallelLText{
\begin{center}
{\large \ggn{12}.}
\end{center}\vspace*{-7pt}

{\greekfont Ἐὰν εἰς κύκλον τρίγωνον ἰσόπλευρον ἐγγραφῇ, ἡ τοῦ
τριγώνου πλευρὰ δυνάμει τριπλασίων ἐστὶ τῆς ἐκ τοῦ
κέντρου τοῦ κύκλου.}

{\greekfont Ἐστω κύκλος ὁ ΑΒΓ, καὶ εἰς α ὐτὸν τρίγωνον ἰσόπλευρ-ον
ἐγγεγράφθω τὸ ΑΒΓ; λέγω, ὅτι τοῦ ΑΒΓ τριγώνου μία πλευρὰ
δυνάμει τριπλασίων ἐστὶ τῆς ἐκ τοῦ κέντρου τοῦ ΑΒΓ κύκλου.}

\epsfysize=2.5in
\centerline{\epsffile{Book13/fig12g.eps}}

{\greekfont Εἰλήφθω γὰρ τὸ κέντρον τοῦ ΑΒΓ κύκλου τὸ Δ, καὶ
ἐπιζευθθεῖσα ἡ ΑΔ διήθθω ἐπὶ τὸ Ε, καὶ ἐπεζεύθθω ἡ ΒΕ.}

{\greekfont Καὶ ἐπεὶ ἰσόπλευρόν
ἐστι τὸ ΑΒΓ τρίγωνον, ἡ ΒΕΓ ἄρα περιφέρεια τρίτον
μέρος ἐστὶ τῆς τοῦ ΑΒΓ κύκλου περιφερείας. ἡ ἄρα ΒΕ
περιφέρεια ἕκτον ἐστὶ μέρος τῆς τοῦ κύκλου περιφερείας;
ἑχαγώνου ἄρα ἐστὶν ἡ ΒΕ εὐθεῖα; ἴση ἄρα ἐστὶ τῇ
ἐκ τοῦ κέντρου τῇ ΔΕ. καὶ ἐπεὶ διπλῆ ἐστιν ἡ ΑΕ
τῆς ΔΕ, τετραπλάσιον ἐστι τὸ ἀπὸ τῆς ΑΕ τοῦ ἀπὸ
τῆς ΕΔ, τουτέστι τοῦ ἀπὸ τῆς ΒΕ. ἴσον δὲ τὸ ἀπὸ
τῆς ΑΕ τοῖς ἀπὸ τῶν ΑΒ, ΒΕ; τὰ ἄρα ἀπὸ τῶν ΑΒ, ΒΕ
τετραπλάσιά ἐστι τοῦ ἀπὸ τῆς ΒΕ. διελόντι ἄρα τὸ ἀπὸ
τῆς ΑΒ τριπλάσιόν ἐστι τοῦ ἀπὸ ΒΕ. ἴση δὲ ἡ ΒΕ τῇ ΔΕ;
τὸ ἄρα ἀπὸ τῆς ΑΒ τριπλάσιόν ἐστι τοῦ ἀπὸ τῆς ΔΕ.}

{\greekfont Ἡ ἄρα τοῦ τριγώνου πλευρὰ δυνάμει τριπλασία
ἐστὶ τῆς ἐκ τοῦ κέντρου [τοῦ κύκλου]; ὅπερ ἔδει
δεῖχαι.}}

\ParallelRText{
\begin{center}
{\large Proposition 12}
\end{center}

If an equilateral triangle is inscribed in a circle then the
square on the side of the triangle is three times the (square) on
the radius of the circle.

Let there be a circle $ABC$, and let the equilateral triangle $ABC$ have
been inscribed in it  [Prop.~4.2]. I say that the square on one side of triangle
$ABC$ is three times the (square) on the radius of circle $ABC$. 

\epsfysize=2.5in
\centerline{\epsffile{Book13/fig12e.eps}}

For let the center, $D$, of circle $ABC$ be found [Prop.~3.1].
And  $AD$ (being) joined, let it be drawn across to $E$. 
And let $BE$ be joined.

And since triangle $ABC$ is equilateral, circumference $BEC$ is thus the
third part of the circumference of circle $ABC$. Thus, 
circumference $BE$ is the sixth part of the circumference of the circle. 
Thus, straight-line $BE$ is (the side) of a hexagon. Thus, it is equal to the radius
$DE$ [Prop.~4.15~corr.]. And since $AE$ is double $DE$, the (square)
on $AE$ is four times the (square) on $ED$---that is to say, of the
(square) on $BE$.  And the (square) on $AE$ (is) equal to
the (sum of the squares) on $AB$ and $BE$ [Props.~3.31, 1.47]. 
Thus, the (sum of the squares) on $AB$ and $BE$ is four times the
(square) on $BE$. Thus, via separation, the (square) on $AB$ is three
times the (square) on $BE$. And $BE$ (is) equal to $DE$. Thus, the
(square) on $AB$ is three times the (square) on $DE$.

Thus, the square on the side of the triangle is three times the (square)
on the radius [of the circle]. (Which is) the very thing it was required to show.}
\end{Parallel}

%%%%
%13.13
%%%%
\pdfbookmark[1]{Proposition 13.13}{pdf13.13}
\begin{Parallel}{}{}
\ParallelLText{
\begin{center}
{\large \ggn{13}.}
\end{center}\vspace*{-7pt}

{\greekfont Πυραμίδα συστήσασθαι καὶ σφαίρᾳ περιλαβεῖν τῇ δοθείσῃ καὶ
δεῖχαι, ὅτι ἡ τῆς σφαίρας διάμετρος δυνάμει ἡμιολία ἐστὶ
τῆς πλευρᾶς τῆς πυραμίδος.}

{\greekfont Ἐκκείσθω ἡ τῆς δοθείσης σφαίρας δίαμετρος ἡ ΑΒ, καὶ
τετμήσθω κατὰ τὸ Γ σημεῖον, ὥστε διπλασίαν εἶναι τὴν
ΑΓ τῆς ΓΒ; καὶ γεγράφθω ἐπὶ τῆς ΑΒ ἡμικύκλιον τὸ
ΑΔΒ, καὶ ἤθθω ἀπὸ τοῦ Γ σημείου τῇ ΑΒ πρὸς ὀρθὰς
ἡ ΓΔ, καὶ ἐπεζεύθθω ἡ ΔΑ; καὶ ἐκκείσθω κύκλος
ὁ ΕΖΗ ἴσην ἔθων τὴν ἐκ τοῦ κέντρου τῇ ΔΓ, καὶ
ἐγγεγράφθω εἰς τὸν ΕΖΗ κύκλον τρίγωνον ἰσόπλευρον
τὸ ΕΖΗ; καὶ εἰλήφθω τὸ κέντρον τοῦ κύκλου τὸ Θ σημεῖον,
καὶ ἐπεζεύθθωσαν αἱ ΕΘ, ΘΖ, ΘΗ; καὶ ἀνεστάτω ἀπὸ τοῦ
Θ σημείου τῷ τοῦ ΕΖΗ κύκλου ἐπιπέδῳ πρὸς ὀρθὰς
ἡ ΘΚ, καὶ ἀφῃρήσθω ἀπὸ τῆς ΘΚ τῇ ΑΓ εὐθείᾳ
ἴση ἡ ΘΚ, καὶ ἐπεζεύθθωσαν αἱ ΚΕ, ΚΖ, ΚΗ. καὶ
ἐπεὶ ἡ ΚΘ ὀρθή ἐστι πρὸς τὸ τοῦ ΕΖΗ κύκλου ἐπίπεδον,
καὶ πρὸς πάσας ἄρα τὰς ἁπτομένας α ὐτῆς εὐθείας
καὶ οὔσας ἐν τῷ τοῦ ΕΖΗ κύκλου ἐπιπέδῳ
ὀρθὰς ποιήσει γωνίας. ἅπτεται δὲ α ὐτῆς ἑκάστη
τῶν ΘΕ, ΘΖ, ΘΗ; ἡ ΘΚ ἄρα πρὸς ἑκάστη τῶν ΘΕ, ΘΖ,
ΘΗ
 ὀρθή ἐστιν. καὶ ἐπεὶ ἴση ἐστὶν ἡ μὲν
ΑΓ τῇ ΘΚ, ἡ δὲ ΓΔ τῇ ΘΕ, καὶ ὀρθὰς γωνίας περιέθουσιν,
βάσις ἄρα ἡ ΔΑ βάσει τῇ ΚΕ ἐστιν ἴση. διὰ τὰ α ὐτὰ
δὴ καὶ ἑκατέρα τῶν ΚΖ, ΚΗ τῇ ΔΑ ἐστιν ἴση;
αἱ τρεῖς ἄρα αἱ ΚΕ, ΚΖ, ΚΗ ἴσαι ἀλλήλαις εἰσίν.
καὶ ἐπεὶ διπλῆ ἐστιν ἡ ΑΓ τῆς ΓΒ, τριπλῆ ἄρα ἡ
ΑΒ τῆς ΒΓ. ὡς δὲ ἡ ΑΒ πρὸς τὴν ΒΓ, οὕτως τὸ ἀπὸ
τῆς ΑΔ πρὸς τὸ ἀπὸ τῆς ΔΓ, ὡς ἑχῆς δειθθήσεται.
τριπλάσιον ἄρα τὸ ἀπὸ τῆς ΑΔ τοῦ ἀπὸ τῆς
ΔΓ. ἔστι δὲ καὶ τὸ ἀπὸ τῆς ΖΕ τοῦ ἀπὸ τῆς ΕΘ τριπλάσιον,
καί ἐστιν ἴση ἡ ΔΓ τῇ ΕΘ; ἴση ἄρα καὶ ἡ ΔΑ τῇ ΕΖ.
ἀλλὰ ἡ ΔΑ ἑκάστῃ τῶν ΚΕ, ΚΖ, ΚΗ ἐδείθθη ἴση;
καὶ ἑκάστη ἄρα τῶν ΕΖ, ΖΗ, ΗΕ ἑκάστῃ τῶν ΚΕ, ΚΖ, ΚΗ
ἐστιν ἴση; ἰσόπλευρα ἄρα ἐστὶ τὰ τέσσαρα τρίγωνα τὰ ΕΖΗ, ΚΕΖ,
ΚΖΗ, ΚΕΗ. πυραμὶς ἄρα συνέσταται ἐκ τεσσάρων τριγώνων
ἰσοπλέυρων, ἧς βάσις μέν ἐστι τὸ ΕΖΗ τρίγωνον,
κορυφὴ δὲ τὸ Κ σημεῖον.}\\~\\~\\

\epsfysize=3.5in
\centerline{\epsffile{Book13/fig13g.eps}}

{\greekfont Δεῖ δὴ α ὐτὴν καὶ σφαίρᾳ περιλαβεῖν τῇ δοθείσῃ καὶ δεῖχαι,
ὅτι ἡ τῆς σφαίρας διάμετρος ἡμιολία ἐστὶ δυνάμει τῆς
πλευρᾶς τῆς πυραμίδος.}

{\greekfont Ἐκβεβλήσθω γὰρ ἐπ εὐθείας τῇ ΚΘ εὐθεῖα ἡ ΘΛ, καὶ
κείσθω τῇ ΓΒ ἴση ἡ ΘΛ. καὶ ἐπεί ἐστιν ὡς ἡ ΑΓ πρὸς
τὴν ΓΔ, οὕτως ἡ ΓΔ πρὸς τὴν ΓΒ, ἴση δὲ ἡ μὲν ΑΓ τῇ
ΚΘ, ἡ δὲ ΓΔ τῇ ΘΕ, ἡ δὲ ΓΒ τῇ ΘΛ, ἔστιν ἄρα ὡς ἡ ΚΘ πρὸς
τὴν ΘΕ, οὕτως ἡ ΕΘ πρὸς τὴν ΘΛ; τὸ ἄρα ὑπὸ τῶν
ΚΘ, ΘΛ ἴσον ἐστὶ τῷ ἀπὸ τῆς ΕΘ. καί ἐστιν ὀρθὴ ἑκατέρα
τῶν ὑπὸ ΚΘΕ, ΕΘΛ γωνιῶν; τὸ ἄρα ἐπὶ τῆς ΚΛ γραφόμενον
ἡμικύκλιον ἥχει καὶ διὰ τοῦ Ε [ἐπειδήπερ ἐὰν ἐπιζεύχωμεν
τὴν ΕΛ, ὀρθὴ γίνεται ἡ ὑπὸ ΛΕΚ γωνία διὰ τὸ ἰσογώνιον
γίνεσθαι τὸ ΕΛΚ τρίγωνον ἑκατέρῳ τῶν ΕΛΘ, ΕΘΚ τριγώνων].
ἐὰν δὴ μενούσης τῆς ΚΛ περιενεθθὲν τὸ ἡμικύκλιον
εἰς τὸ α ὐτὸ πάλιν ἀποκατασταθῇ, ὅθεν ἤρχατο φέρεσθαι,
ἥχει καὶ διὰ τῶν Ζ, Η σημείων ἐπιζευγνυμένων
τῶν ΖΛ, ΛΗ καὶ ὀρθῶν ὁμοίως γινομένων τῶν πρὸς τοῖς
Ζ, Η γωνιῶν; καὶ ἔσται ἡ πυραμὶς σφαίρᾳ περιειλημμένη τῇ
δοθείσῇ. ἡ γὰρ ΚΛ τῆς σφαίρας διάμετρος ἴση ἐστὶ τῇ
τῆς δοθείσης σφαίρας διαμετρῳ τῇ ΑΒ, ἐπειδήπερ τῇ μὲν
ΑΓ ἴση κεῖται ἡ ΚΘ, τῇ δὲ ΓΒ ἡ ΘΛ.}

{\greekfont Λέγω δή, ὅτι ἡ τῆς σφαίρας διάμετρος ἡμιολία
ἐστὶ δυνάμει τῆς πλευρᾶς τῆς πυραμίδος.}

{\greekfont Ἐπεὶ γὰρ διπλῆ ἐστιν ἡ ΑΓ τῆς ΓΒ, τριπλῆ ἄρα
ἐστὶν ἡ ΑΒ τῆς ΒΓ; ἀναστρέψαντι ἡμιολία ἄρα ἐστὶν
ἡ ΒΑ τῆς
ΑΓ. ὡς δὲ ἡ ΒΑ πρὸς τὴν ΑΓ, οὕτως
τὸ ἀπὸ τῆς ΒΑ πρὸς τὸ ἀπὸ τῆς ΑΔ [ἐπειδήπερ ἐπιζευγνμένης
τῆς ΔΒ ἐστιν ὡς ἡ ΒΑ πρὸς τὴν ΑΔ, οὕτως ἡ ΔΑ πρὸς
τὴν ΑΓ διὰ τὴν ὁμοιότητα τῶν ΔΑΒ, ΔΑΓ τριγώνων, καὶ
εἶναι ὡς τὴν πρώτην πρὸς τὴν τρίτην, οὕτως
τὸ ἀπὸ τῆς τρώτης πρὸς τὸ ἀπὸ τῆς δευτέρας].
ἡμιόλιον ἄρα καὶ τὸ ἀπὸ τῆς ΒΑ τοῦ ἀπὸ τῆς ΑΔ. καί
ἐστιν ἡ μὲν ΒΑ ἡ τῆς δοθείσης
σφαίρας διάμετρος, ἡ δὲ ΑΔ ἴση τῇ πλευρᾷ
τῆς πυραμίδος.}

{\greekfont Ἡ ἄρα τῆς σφαίρας διάμετρος ἡμιολία ἐστὶ τῆς
πλευρᾶς τῆς πυραμίδος; ὅπερ ἔδει δεῖχαι.}}

\ParallelRText{
\begin{center}
{\large Proposition 13}
\end{center}

To construct a (regular) pyramid ({\em i.e.}, a tetrahedron), and to
enclose (it) in a given sphere, and to show that the square on the
diameter of the sphere is one and a half times the (square) on the side 
of the pyramid. 

Let the diameter $AB$ of the given sphere be laid out, and let it be
cut at point $C$ such that $AC$ is double $CB$ [Prop.~6.10]. And let the semi-circle 
$ADB$  be drawn on $AB$. And let $CD$ be drawn from point
$C$ at right-angles to $AB$. And let $DA$ be joined. 
And let the circle $EFG$ be laid down having a radius equal to $DC$,
and let the equilateral triangle $EFG$ be inscribed in circle
$EFG$ [Prop.~4.2]. And let the center of the circle, point $H$,
be found [Prop.~3.1]. And let $EH$, $HF$, and $HG$ be
joined. And let $HK$ be set up, at point $H$, at right-angles to the
plane of circle $EFG$ [Prop.~11.12]. And let $HK$, equal to the straight-line $AC$, be cut off from $HK$. And let $KE$, $KF$, and $KG$ be joined.
And since $KH$ is at right-angles to the plane of circle $EFG$, it will
thus also make right-angles with all of the straight-lines joining it
(which are) also in the plane of circle $EFG$ [Def.~11.3]. And $HE$, $HF$,
and $HG$ each join it. Thus, $HK$ is at right-angles to each of $HE$,
$HF$, and $HG$. And since  $AC$ is equal to $HK$, and $CD$ to
$HE$, and they contain right-angles, the base $DA$ is thus equal to the base 
$KE$ [Prop.~1.4]. So, for the same (reasons),  $KF$ and
$KG$ is each equal to $DA$. Thus, the three (straight-lines)
$KE$, $KF$, and $KG$ are equal to one another. And since 
$AC$ is double $CB$, $AB$ (is) thus triple $BC$. And as 
$AB$ (is) to $BC$, so the (square) on $AD$ (is) to the (square)
on $DC$, as will be shown later [see lemma]. Thus, the (square)
on $AD$ (is) three times the (square) on $DC$. And the (square)
on $FE$ is also three times the (square) on $EH$ [Prop.~13.12], 
and $DC$ is equal to $EH$. Thus, $DA$ (is) also equal to $EF$.
But, $DA$ was shown (to be) equal to each of $KE$, $KF$, and
$KG$. Thus, $EF$, $FG$, and $GE$ are equal to 
$KE$, $KF$, and $KG$, respectively. Thus, the four triangles $EFG$, $KEF$,
$KFG$, and $KEG$ are equilateral. Thus, a pyramid, whose base is triangle
$EFG$, and apex the point $K$,   has been
constructed from four equilateral triangles.

\epsfysize=3.5in
\centerline{\epsffile{Book13/fig13e.eps}}

So, it is also necessary to enclose it in the given sphere, and to show that
the square on the diameter of the sphere is one and a half times the (square)
on the side of the pyramid. 

For let the straight-line $HL$ be produced in a straight-line with $KH$, and let
$HL$ be made equal to $CB$. And since as $AC$ (is) to $CD$, 
so $CD$ (is) to $CB$ [Prop.~6.8~corr.], and $AC$ (is) equal to $KH$, and $CD$ to
$HE$, and $CB$ to $HL$, thus as $KH$ is to $HE$, so $EH$ (is) to $HL$.
Thus,  the (rectangle contained) by $KH$ and $HL$ is equal to the
(square) on $EH$ [Prop.~6.17]. And each of the angles $KHE$ and
$EHL$ is a right-angle. Thus, the semi-circle drawn on $KL$ 
will also pass through $E$ [inasmuch as if we join $EL$ then the angle
$LEK$ becomes a right-angle, on account of triangle $ELK$
becoming equiangular to each of the triangles $ELH$ and $EHK$ [Props.~6.8, 3.31]\,]. So, if
$KL$ remains (fixed), and the semi-circle is carried around, and again
established at the same (position) from which it began to be moved, 
it will also pass through points $F$ and $G$, (because) if  $FL$ and
$LG$ are joined, the angles at $F$ and $G$ will similarly become  right-angles. And the pyramid will be enclosed by the given sphere.
For the diameter, $KL$,  of the sphere is equal to the diameter, $AB$,
of the given sphere---inasmuch as $KH$ was made equal to $AC$, 
and $HL$ to $CB$. 

So, I say that the square on the diameter of the sphere is
one and a half times the (square) on the side of the pyramid.

For since $AC$ is double $CB$, $AB$ is thus triple $BC$. Thus,
via conversion, $BA$ is one and a half times $AC$. And as
$BA$ (is) to $AC$, so the (square) on $BA$ (is) to the (square) on $AD$
[inasmuch as if $DB$ is joined then as $BA$ is to $AD$, so
$DA$ (is) to $AC$, on account of the similarity of triangles
$DAB$ and $DAC$. And as the first is to the third (of four
proportional magnitudes), so the (square) on the first (is) to the
(square) on the second.] Thus, the (square) on $BA$
(is) also one and a half times the (square) on $AD$. 
And $BA$ is the diameter of the given sphere, and $AD$
(is) equal to the side of the pyramid.

Thus, the square on the diameter of the sphere is one and a half
times the (square) on the side of the pyramid.$^\dag$ (Which is) the very thing
it was required to show.}
\end{Parallel}
{\footnotesize\noindent$^\dag$ If the radius of the sphere is unity then the side of the pyramid ({\em i.e.}, tetrahedron) is $\sqrt{8/3}$.}

\begin{Parallel}{}{}
\ParallelLText{
\epsfysize=2.5in
\centerline{\epsffile{Book13/fig13ag.eps}}

\begin{center}\vspace*{-7pt}
{\large {\greekfont Λῆμμα}.}
\end{center}\vspace*{-7pt}

{\greekfont Δεικτέον, ὅτι ἐστὶν ὡς ἡ ΑΒ πρὸς τὴν ΒΓ, οὕτως
τὸ ἀπὸ τῆς ΑΔ πρὸς τὸ ἀπὸ τῆς ΔΓ.}

{\greekfont Ἐκκείσθω γὰρ ἡ τοῦ ἡμικυκλίου καταγραφή, καὶ ἐπεζεύθθω
ἡ ΔΒ, καὶ ἀναγεγράφθω ἀπὸ τῆς ΑΓ τετράγωνον τὸ ΕΓ,
καὶ συμπεπληρώσθω τὸ ΖΒ παραλληλόγραμμον. ἐπεὶ οὖν
διὰ τὸ ἰσογώνιον εἶναι τὸ ΔΑΒ τρίγωνον τῷ ΔΑΓ
τριγώνῳ ἐστὶν ὡς ἡ ΒΑ πρὸς τὴν ΑΔ, οὕτως ἡ ΔΑ πρὸς
τὴν ΑΓ, τὸ ἄρα ὑπὸ τῶν ΒΑ, ΑΓ ἴσον ἐστὶ τῷ ἀπὸ
τῆς ΑΔ. καὶ ἐπεί ἐστιν ὡς ἡ ΑΒ πρὸς τὴν ΒΓ, οὕτως
τὸ ΕΒ πρὸς τὸ ΒΖ, καί ἐστι τὸ μὲν ΕΒ τὸ ὑπὸ τῶν ΒΑ,
ΑΓ; ἴση γὰρ ἡ ΕΑ τῇ ΑΓ; τὸ δὲ ΒΖ τὸ ὑπὸ τῶν
ΑΓ, ΓΒ, ὡς ἄρα ἡ ΑΒ πρὸς τὴν ΒΓ, οὕτως τὸ ὑπὸ τῶν
ΒΑ, ΑΓ πρὸv
τὸ ὑπὸ τῶν ΑΓ, ΓΒ. καί
ἐστι τὸ μὲν  ὑπὸ τῶν ΒΑ, ΑΓ ἴσον τῷ ἀπὸ τῆς ΑΔ, τὸ
δὲ ὑπὸ τῶν ΑΓΒ ἴσον τῷ ἀπὸ τῆς ΔΓ; ἡ γὰρ ΔΓ
κάθετος τῶν τῆς βάσεως τμ ημάτων τῶν ΑΓ, ΓΒ μέση
ἀνάλογόν ἐστι διὰ τὸ ὀρθὴν εἶναι τὴν ὑπὸ ΑΔΒ.
ὡς ἄρα ἡ ΑΒ πρὸς τὴν ΒΓ, οὕτως τὸ ἀπὸ τῆς
ΑΔ πρὸς τὸ ἀπὸ τῆς ΔΓ; ὅπερ ἔδει δεῖχαι.}}

\ParallelRText{
\epsfysize=2.5in
\centerline{\epsffile{Book13/fig13ae.eps}}

\begin{center}
{\large Lemma}
\end{center}

It must be shown that as $AB$ is to $BC$, so the (square) on
$AD$ (is) to the (square) on $DC$.

For, let the figure of the semi-circle be set out, and let $DB$ have
been joined. And let the square $EC$ be described on
$AC$. And let the parallelogram $FB$ be completed. 
Therefore, since, on account of triangle $DAB$ being
equiangular to  triangle $DAC$ [Props.~6.8, 6.4], (proportionally) as $BA$
is to $AD$, so $DA$ (is) to $AC$,  the (rectangle contained)
by $BA$ and $AC$  is thus equal to the (square) on $AD$ [Prop.~6.17]. 
And since as $AB$ is to $BC$, so $EB$ (is) to $BF$ [Prop.~6.1].
And $EB$ is  the (rectangle contained) by $BA$ and $AC$---for $EA$ (is) equal to $AC$. And $BF$  the (rectangle
contained) by $AC$ and $CB$. Thus, as $AB$ (is) to $BC$,
so the (rectangle contained) by $BA$ and $AC$ (is) to the
(rectangle contained) by $AC$ and $CB$. And the (rectangle contained)
by $BA$ and $AC$ is equal to the (square) on $AD$, and the
(rectangle contained) by $ACB$ (is) equal to the (square) on 
$DC$. For the perpendicular $DC$ is the mean proportional
to the pieces  of the base, $AC$ and $CB$, on account of
$ADB$ being a right-angle [Prop.~6.8~corr.]. Thus, as
$AB$ (is) to $BC$, so the (square) on $AD$ (is) to the (square)
on $DC$. (Which is) the very thing it was required to show.}
\end{Parallel}

%%%%
%13.14
%%%%
\pdfbookmark[1]{Proposition 13.14}{pdf13.14}
\begin{Parallel}{}{}
\ParallelLText{
\begin{center}
{\large \ggn{14}.}
\end{center}\vspace*{-7pt}

{\greekfont Ὀκτάεδρον συστήσασθαι καὶ σφαίρᾳ περιλαβεῖν, ᾗ καὶ τὰ
πρότερα, καὶ δεῖχαι, ὅτι ἡ τῆς σφαίρας διάμετρος δυνάμει
διπλασία ἐστὶ τῆς πλευρᾶς τοῦ ὀκταέδρου.}\\

\epsfysize=3.25in
\centerline{\epsffile{Book13/fig14g.eps}}

{\greekfont Ἐκκείσθω ἡ τῆς δοθείσης σφαίρας διάμετρος ἡ ΑΒ, καὶ
τετμήσθω δίθα κατὰ τὸ Γ, καὶ γεγράφθω ἐπὶ τῆς ΑΒ
ἡμικύκλιον τὸ ΑΔΒ, καὶ ἤθθω ἀπὸ τοῦ
Γ τῇ ΑΒ πρὸς ὀρθὰς ἡ ΓΔ,  καὶ ἐπεζεύθθω ἡ ΔΒ,
καὶ ἐκκείσθω τετράγωνον τὸ ΕΖΗΘ ἴσην ἔθον ἑκάστην
τῶν πλευρῶν τῇ ΔΒ, καὶ ἐπεζεύθθωσαν αἱ ΘΖ, ΕΗ, καὶ
ἀνεστάτω ἀπὸ τοῦ Κ σημείου τῷ τοῦ ΕΖΗΘ τετραγώνου
ἐπιπέδῳ πρὸς ὀρθὰς εὐθεῖα ἡ ΚΛ καὶ διήθθω
ἐπὶ τὰ ἕτερα μέρη τοῦ ἐπιπέδου ὡς ἡ ΚΜ, καὶ
ἀφῃρήσθω ἀφ ἑκατέρας τῶν ΚΛ, ΚΜ μιᾷ τῶν ΕΚ, ΖΚ, ΗΚ, ΘΚ
ἴση ἑκατέρα τῶν ΚΛ, ΚΜ, καὶ ἐπεζεύθθωσαν αἱ ΛΕ, ΛΖ, ΛΗ, ΛΘ,
ΜΕ, ΜΖ, ΜΗ, ΜΘ.}

{\greekfont Καὶ ἐπεὶ ἴση ἐστὶν ἡ ΚΕ τῇ ΚΘ, καί
ἐστιν ὀρθὴ ἡ ὑπὸ ΕΚΘ γωνία, τὸ ἄρα ἀπὸ τῆς ΘΕ
διπλάσιόν ἐστι τοῦ ἀπὸ τῆς ΕΚ. πάλιν, ἐπεὶ ἴση ἐστὶν
ἡ ΛΚ τῇ ΚΕ, καί ἐστιν ὀρθὴ ἡ ὑπὸ ΛΚΕ γωνία,
τὸ ἄρα ἀπὸ τῆς ΕΛ διπλάσιόν ἐστι τοῦ ἀπὸ ΕΚ.
ἐδείθθη δὲ καὶ τὸ ἀπὸ τῆς ΘΕ διπλάσιον τοῦ ἀπὸ τῆς
ΕΚ; τὸ ἄρα ἀπὸ τῆς ΛΕ ἴσον ἐστὶ τῷ ἀπὸ τῆς
ΕΘ; ἴση ἄρα ἐστὶν ἡ ΛΕ τῇ ΕΘ. διὰ τὰ α ὐτὰ δὴ καὶ ἡ
ΛΘ τῇ ΘΕ ἐστιν ἴση; ἰσόπλευρον ἄρα ἐστὶ τὸ ΛΕΘ τρίγωνον.
ὁμοίως δὴ δείχομεν, ὅτι καὶ ἕκαστον τῶν λοιπῶν
τριγώνων, ὧν βάσεις μέν εἰσιν αἱ τοῦ ΕΖΗΘ τετραγώνου
πλευραί, κορυφαὶ δὲ τὰ Λ, Μ σημεῖα, ἰσόπλευρόν ἐστιν; ὀκτάεδρον
ἄρα συνέσταται ὑπὸ ὀκτὼ τριγώνων ἰσοπλεύρων περιεχόμενον.}

{\greekfont Δεῖ δὴ α ὐτὸ καὶ σφαίρᾳ περιλαβεῖν τῇ δοθείσῃ καὶ
δεῖχαι, ὅτι ἡ τῆς σφαίρας διάμετρος δυνάμει διπλασίων
ἐστὶ τῆς τοῦ ὀκταέδρου πλευρᾶς.}

{\greekfont Ἐπεὶ γὰρ αἱ τρεῖς αἱ ΛΚ, ΚΜ, ΚΕ ἴσαι
ἀλλήλαις εἰσίν, τὸ ἄρα ἐπὶ τῆς ΛΜ γραφόμενον
ἡμικύκλιον ἥχει καὶ διὰ τοῦ Ε. καὶ διὰ τὰ α ὐτά,
ἐὰν μενούσης τῆς ΛΜ περιενεθθὲν τὸ ἡμικύκλιον
εἰς τὸ α ὐτὸ ἀποκατασταθῇ, ὅθεν ἤρχατο φέρεσθαι,
ἥχει καὶ διὰ τῶν Ζ, Η, Θ σημείων, καὶ ἔσται σφαίρᾳ
περιειλημμένον τὸ ὀκτάεδρον. λέγω δή, ὅτι καὶ τῇ δοθείσῃ.
ἐπεὶ γὰρ ἴση ἐστὶν ἡ ΛΚ τῇ ΚΜ, κοινὴ δὲ ἡ ΚΕ, καὶ
γωνίας ὀρθὰς περιέθουσιν, βάσις ἄρα ἡ ΛΕ βάσει τῇ ΕΜ
ἐστιν ἴση. καὶ ἐπεὶ ὀρθή ἐστιν ἡ ὑπὸ ΛΕΜ γωνία;
ἐν ἡμικυκλίῳ γάρ; τὸ ἄρα ἀπὸ τῆς ΛΜ διπλάσιόν
ἐστι τοῦ ἀπὸ τῆς ΛΕ. πάλιν, ἐπεὶ ἴση ἐστὶν ἡ ΑΓ τῇ
ΓΒ, διπλασία ἐστὶν ἡ ΑΒ τῆς ΒΓ. ὡς δὲ ἡ ΑΒ πρὸς τὴν
ΒΓ, οὕτως τὸ ἀπὸ τῆς ΑΒ πρὸς τὸ ἀπὸ τῆς ΒΔ; διπλάσιον
ἄρα ἐστὶ τὸ ἀπὸ τῆς ΑΒ τοῦ ἀπὸ τῆς ΒΔ. ἐδείθθη
δὲ καὶ τὸ ἀπὸ τῆς ΛΜ διπλάσιον τοῦ ἀπὸ τῆς ΛΕ.
καί ἐστιν ἴσον τὸ ἀπὸ τῆς ΔΒ τῷ ἀπὸ τῆς ΛΕ;
ἴση γὰρ κεῖται ἡ ΕΘ τῇ ΔΒ. ἴσον ἄρα καὶ τὸ ἀπὸ τῆς
ΑΒ τῷ ἀπὸ τῆς ΛΜ; ἴση ἄρα ἡ ΑΒ τῇ ΛΜ. καί ἐστιν
ἡ ΑΒ ἡ τῆς δοθείσης σφαίρας διάμετρος; ἡ ΛΜ ἄρα ἴση
ἐστὶ τῇ τῆς δοθείσης σφαίρας διαμέτρῳ.}

{\greekfont Περιείληπται ἄρα τὸ ὀκτάεδρον τῇ δοθείσῃ σφαίρᾳ.
καὶ συναποδέδεικται, ὅτι ἡ τῆς σφαίρας διάμετρος δυνάμει
διπλασίων ἐστὶ τῆς τοῦ ὀκταέδρου πλευρᾶς; ὅπερ
ἔδει δεῖχαι.}}

\ParallelRText{
\begin{center}
{\large Proposition 14}
\end{center}

To construct an octahedron, and to enclose (it) in a (given) sphere, like in the preceding (proposition), and to show that the square
on the diameter of the sphere is double the (square) on the
side of the octahedron.

\epsfysize=3.25in
\centerline{\epsffile{Book13/fig14e.eps}}

Let the diameter $AB$ of the given sphere be laid out, and let it have
been cut in half at $C$. And let the semi-circle $ADB$ be
drawn on $AB$. And let $CD$ be drawn from $C$ at right-angles to
$AB$. And let $DB$ be joined. And let the square $EFGH$,
having each of its sides equal to $DB$, be laid out. And let $HF$
and $EG$ be joined. And let the straight-line $KL$
be set up, at point $K$, at right-angles to the plane of square 
$EFGH$ [Prop.~11.12]. And let it be drawn across  on the other side of
the plane, like $KM$. And let $KL$ and $KM$, equal to
one of $EK$, $FK$, $GK$, and $HK$,  be cut off from 
$KL$ and $KM$, respectively. And let $LE$, $LF$, $LG$, $LH$, $ME$, $MF$,
$MG$, and $MH$ be joined.  

And since  $KE$ is equal to $KH$, and angle $EKH$ is a right-angle, 
the (square) on the  $HE$ is thus double the (square) on $EK$ [Prop.~1.47].
Again, since $LK$ is equal to $KE$, and angle $LKE$
is a right-angle, the (square) on $EL$ is thus double the (square) on
$EK$ [Prop.~1.47]. And the (square) on $HE$ was also shown (to be)
double the (square) on $EK$.  Thus, the (square) on $LE$
is equal to the (square) on $EH$. Thus, $LE$ is equal to $EH$. So,
for the same (reasons), $LH$ is also equal to $HE$. Triangle
$LEH$ is thus equilateral. So, similarly, we can show that each of
the remaining triangles, whose bases are the sides of the
square $EFGH$, and apexes the points $L$ and $M$, are equilateral.
Thus, an octahedron contained by eight
equilateral triangles has been constructed.

So, it is also necessary to enclose it by the given sphere, and to show that the
square on the diameter of the sphere is double the (square) on the
side of the octahedron.

For since the three (straight-lines) $LK$, $KM$, and $KE$ are equal
to one another, the semi-circle drawn on $LM$ will thus also pass
through $E$. And, for the same (reasons), if $LM$ remains (fixed), and the semi-circle is carried around, and again established at the same
(position) from which it began to be moved, then it will also pass
through points $F$, $G$, and $H$, and the octahedron will be
enclosed by a sphere. So, I say that (it is) also (enclosed)
by the given (sphere). For since $LK$ is equal to $KM$, and $KE$
(is) common, and they contain right-angles, the base
$LE$ is thus equal to the base $EM$ [Prop.~1.4]. And since angle $LEM$
is a right-angle---for (it is) in a semi-circle [Prop.~3.31]---the
(square) on $LM$ is thus double the (square) on $LE$ [Prop.~1.47]. 
Again, since $AC$ is equal to $CB$, $AB$ is double $BC$. 
And as $AB$ (is) to $BC$, so the (square) on $AB$ (is) to the
(square) on $BD$ [Prop.~6.8, Def.~5.9]. Thus, the (square) on 
$AB$ is double the (square) on $BD$. And the (square) on $LM$
was also shown (to be) double the (square) on $LE$. And the (square)
on $DB$ is equal to the (square) on $LE$. For $EH$ was made equal
to $DB$. Thus, the (square) on $AB$ (is) also equal to the (square) on 
$LM$. Thus, $AB$ (is) equal to $LM$. And $AB$ is the diameter of the
given sphere. Thus, $LM$ is equal to the diameter of the given sphere.

Thus, the octahedron has been enclosed by the given sphere,  and it has
been simultaneously  proved that the square on the diameter of the sphere
is double the (square) on the side of the octahedron.$^\dag$ (Which is) the
very thing it was required to show.}
\end{Parallel}
{\footnotesize\noindent$^\dag$ If the radius of the sphere is unity then the
side of octahedron is $\sqrt{2}$.}

%%%%
%13.15
%%%%
\pdfbookmark[1]{Proposition 13.15}{pdf13.15}
\begin{Parallel}{}{}
\ParallelLText{
\begin{center}
{\large \ggn{15}.}
\end{center}\vspace*{-7pt}

{\greekfont Κύβον συστήσασθαι καὶ σφαίρᾳ περιλαβεῖν, ᾗ καὶ τὴν
πυραμίδα, καὶ δεῖχαι, ὅτι ἡ τῆς σφαίρας διάμετρος δυνάμει
τριπλασίων ἐστὶ τῆς τοῦ κύβου πλευρᾶς.}

{\greekfont Ἐκκείσθω ἡ τῆς δοθείσης σφαίρας διάμετρος ἡ ΑΒ καὶ
τετμήσθω κατὰ τὸ Γ ὥστε διπλῆν εἶναι τὴν ΑΓ τῆς
ΓΒ, καὶ γεγράφθω ἐπὶ τῆς ΑΒ ἡμικύκλιον τὸ ΑΔΒ, καὶ
ἀπὸ τοῦ Γ τῇ ΑΒ πρὸς ὀρθὰς ἤθθω ἡ ΓΔ, καὶ
ἐπεζεύθθω ἡ ΔΒ, καὶ ἐκκείσθω τετράγωνον τὸ ΕΖΗΘ
ἴσην ἔθον τὴν πλευρὰν τῇ ΔΒ, καὶ ἀπὸ τῶν Ε, Ζ, Η, Θ
τῷ τοῦ ΕΖΗΘ τετραγώνου ἐπιπέδῳ πρὸς ὀρθὰς ἤθθωσαν
αἱ ΕΚ, ΖΛ, ΗΜ, ΘΝ, καὶ ἀφῃρήσθω ἀπὸ ἑκάστης
τῶν ΕΚ, ΖΛ, ΗΜ, ΘΝ μιᾷ τῶν ΕΖ, ΖΗ, ΗΘ, ΘΕ ἴση
ἑκάστη τῶν ΕΚ, ΖΛ, ΗΜ, ΘΝ, καὶ ἐπεζεύθθωσαν αἱ ΚΛ,
ΛΜ, ΜΝ, ΝΚ; κύβος ἄρα συνέσταται ὁ ΖΝ ὑπὸ ἓχ
τετραγώνων ἴσων περιεθόμενος.}

{\greekfont Δεῖ δὴ α ὐτὸν καὶ σφαίρᾳ
περιλαβεῖν τῇ δοθείσῃ καὶ δεῖχαι, ὅτι ἡ τῆς σφαίρας
διάμετρος δυνάμει τριπλασία ἐστὶ τῆς πλευρᾶς τοῦ κύβου.}\\

\epsfysize=3.in
\centerline{\epsffile{Book13/fig15g.eps}}

{\greekfont Ἐπεζεύθθωσαν γὰρ αἱ ΚΗ, ΕΗ. καὶ ἐπεὶ ὀρθή
ἐστιν ἡ ὑπὸ ΚΕΗ γωνία διὰ τὸ καὶ τὴν ΚΕ ὀρθὴν
εἶναι πρὸς τὸ ΕΗ ἐπίπεδον δηλαδὴ καὶ πρὸς τὴν ΕΗ
εὐθεῖαν, τὸ ἄρα ἐπὶ τῆς ΚΗ γραφόμενον ἡμικύκλιον
ἥχει καὶ διὰ τοῦ Ε σημείου. πάλιν, ἐπεὶ ἡ ΗΖ ὀρθή
ἐστι πρὸς ἑκατέραν τῶν ΖΛ, ΖΕ, καὶ πρὸς τὸ ΖΚ ἄρα
ἐπίπεδον ὀρθή ἐστιν ἡ ΗΖ; ὥστε καὶ ἐὰν ἐπιζεύχωμεν
τὴν ΖΚ, ἡ ΗΖ ὀρθὴ ἔσται καὶ πρὸς τὴν ΖΚ;
καὶ δὶα τοῦτο πάλιν τὸ ἐπὶ τῆς ΗΚ γραφόμενον ἡμικύκλιον
ἥχει καὶ διὰ τοῦ Ζ. ὁμοίως καὶ δὶα τῶν λοιπῶν
τοῦ κύβου σημείων ἥχει. ἐὰν δὴ μενούσης τῆς ΚΗ
περιενεθθὲν τὸ ἡμικύκλιον εἰς τὸ α ὐτὸ ἀποκατασταθῇ, 
ὅθεν ἤρχατο φέρεσθαι, ἔσται σφαίρᾳ περιειλημμένος
ὁ κύβος. λέγω δή, ὅτι καὶ τῇ δοθείσῃ. ἐπεὶ γὰρ ἴση ἐστὶν
ἡ ΗΖ τῇ ΖΕ, καί ἐστιν ὀρθὴ ἡ πρὸς τῷ Ζ γωνία, τὸ ἄρα
ἀπὸ τῆς ΕΗ διπλάσιόν ἐστι τοῦ ἀπὸ τῆς ΕΖ. ἴση δὲ ἡ ΕΖ
τῇ ΕΚ; τὸ  ἄρα ἀπὸ τῆς ΕΗ διπλάσιόν ἐστι τοῦ ἀπὸ
τῆς ΕΚ; ὥστε τὰ ἀπὸ τῶν ΗΕ, ΕΚ, τουτέστι τὸ ἀπὸ τῆς
ΗΚ, τριπλάσιόν ἐστι τοῦ ἀπὸ τῆς ΕΚ. καὶ ἐπεὶ τριπλασίων
ἐστὶν ἡ ΑΒ τῆς ΒΓ, ὡς δὲ ἡ ΑΒ πρὸς τὴν ΒΓ, οὕτως
τὸ ἀπὸ τῆς ΑΒ πρὸς τὸ ἀπὸ τῆς ΒΔ, τριπλάσιον ἄρα τὸ
ἀπὸ τῆς ΑΒ τοῦ ἀπὸ τῆς ΒΔ. ἐδείθθη δὲ καὶ τὸ ἀπὸ
τῆς ΗΚ τοῦ ἀπὸ τῆς ΚΕ τριπλάσιον. καὶ κεῖται ἴση
ἡ ΚΕ τῇ ΔΒ; ἴση ἄρα καὶ ἡ ΚΗ τῇ ΑΒ. καί
ἐστιν ἡ ΑΒ τῆς δοθείσης σφαίρας διάμετρος; καὶ ἡ ΚΗ
ἄρα ἴση ἐστὶ τῇ τῆς δοθείσης σφαίρας διαμέτρῳ.}

{\greekfont Τῇ δοθείσῃ ἄρα σφαίρα περιείληπται ὁ κύβος; καὶ συναποδέδεικται,
ὅτι ἡ τῆς σφαίρας διάμετρος δυνάμει τριπλασίων ἐστὶ
τῆς τοῦ κύβου πλευρᾶς; ὅπερ ἔδει δεῖχαι.}}

\ParallelRText{
\begin{center}
{\large Proposition 15}
\end{center}

To construct a cube, and to enclose (it) in a sphere, like in the (case of the)
pyramid, and to show that the square on the diameter of the sphere
is three times the (square) on the side of the cube.

Let the diameter $AB$ of the given sphere be laid out, and let it be
cut at $C$ such that $AC$ is double $CB$. And let the semi-circle
$ADB$ be drawn on $AB$. And let $CD$ be drawn from $C$
at right-angles to $AB$. And let $DB$ be joined. And let the
square $EFGH$, having (its) side equal to $DB$, be laid out. And
let $EK$, $FL$, $GM$, and $HN$ be drawn from (points)
$E$, $F$, $G$, and $H$, (respectively), at right-angles to the plane of
square $EFGH$. And let $EK$, $FL$, $GM$, and $HN$, equal
to one of $EF$, $FG$, $GH$, and $HE$, be cut off from
$EK$, $FL$, $GM$, and $HN$, respectively. And let $KL$,
$LM$, $MN$, and $NK$ be joined. Thus, a cube contained
by six equal squares has been constructed. 

So, it is also necessary
to enclose it by the given sphere, and to show that the square on
the diameter of the sphere is three times the (square) on the side of the cube.

\epsfysize=3.in
\centerline{\epsffile{Book13/fig15e.eps}}

For let $KG$ and $EG$ be joined. And since
angle $KEG$ is a right-angle---on account of $KE$ also being at right-angles
to the plane $EG$, and manifestly also to the straight-line $EG$ [Def.~11.3]---the semi-circle drawn on $KG$ will thus also pass through
point $E$. Again, since $GF$ is at right-angles to each of $FL$
and $FE$, $GF$ is thus also at right-angles to the plane $FK$ [Prop.~11.4]. 
Hence, if we also join $FK$ then $GF$ will also be at right-angles to
$FK$. And, again, on account of this, the semi-circle drawn on
$GK$ will also pass through point $F$. Similarly, it will also pass through
the remaining (angular) points of the cube. So, if $KG$ remains (fixed), 
and the semi-circle is carried around, and again established at the
same (position) from which it began to be moved, then the cube will
be enclosed by a sphere. So, I say that (it is) also (enclosed)
by the given (sphere). For since $GF$ is equal to $FE$, and the
angle at $F$ is a right-angle, the (square) on $EG$ is thus double
the (square) on $EF$ [Prop.~1.47]. And $EF$ (is) equal to
$EK$. Thus, the (square) on $EG$ is double the (square) on $EK$.
Hence, the (sum of the squares) on $GE$ and $EK$---that is to say, the
(square) on $GK$ [Prop.~1.47]---is three times the  (square) on $EK$. And since
$AB$ is three times $BC$, and as $AB$ (is) to $BC$, so the
(square) on $AB$ (is) to the (square) on $BD$ [Prop.~6.8, Def.~5.9], 
the (square) on $AB$ (is) thus three times the (square) on $BD$. 
And the (square) on $GK$ was also shown (to be) three times
the (square) on $KE$. And $KE$ was made equal to $DB$. 
Thus, $KG$ (is) also equal to $AB$. And $AB$ is the
radius of the given sphere. Thus, $KG$ is also equal to the diameter of the
given sphere.

Thus, the cube has been enclosed by the given sphere. And it
has simultaneously been shown that the square on the diameter of
the sphere is three times the (square) on the side of the cube.$^\dag$ (Which is)
the very thing it was required to show.}
\end{Parallel}
{\footnotesize\noindent$^\dag$ If the radius of the sphere is unity 
then the side of the cube is $\sqrt{4/3}$.}

%%%%
%13.16
%%%%
\pdfbookmark[1]{Proposition 13.16}{pdf13.16}
\begin{Parallel}{}{}
\ParallelLText{
\begin{center}
{\large \ggn{16}.}
\end{center}\vspace*{-7pt}

{\greekfont Εἰκοσάεδρον συστήσασθαι καὶ σφαίρᾳ περιλαβεῖν,
ᾗ καὶ τὰ προειρημένα σθήματα, καὶ δεῖχαι, ὅτι
ἡ τοῦ εἰκοσαέδρου πλευρὰ ἄλογός ἐστιν ἡ καλουμένη
ἐλάττων.}

\epsfysize=1.5in
\centerline{\epsffile{Book13/fig16ag.eps}}

{\greekfont Ἐκκείσθω ἡ τῆς δοθείσης σφαίρας διάμετρος
ἡ ΑΒ καὶ τετμήσθω κατὰ τὸ Γ ὥστε τετραπλῆν εἶναι
τὴν ΑΓ τῆς ΓΒ, καὶ γεγράφθω ἐπὶ τῆς ΑΒ ἡμικύκλιον
τὸ ΑΔΒ, καὶ ἤθθω ἀπὸ τοῦ Γ τῇ ΑΒ πρὸς ορθὰς
γωνίας εὐθεῖα γραμμὴ ἡ ΓΔ, καί ἐπεζεύθθω ἡ ΔΒ,
καὶ ἐκκείσθω κύκλος ὁ ΕΖΗΘΚ, οὗ ἡ ἐν τοῦ κέντρου
ἴση ἔστω τῇ ΔΒ, καὶ ἐγγεγράφθω εἰς τὸν
ΕΖΗΘΚ κύκλον πεντάγωνον ἰσόπλευρόν τε καὶ ἰσογώνιον
τὸ ΕΖΗΘΚ, καὶ τετμήσθωσαν αἱ ΕΖ, ΖΗ, ΗΘ, ΘΚ, ΚΕ περιφέρειαι
δίθα κατὰ τὸ Λ, Μ, Ν, Χ, Ο σημεῖα, καὶ ἐπεζεύθθωσαν αἱ
ΛΜ, ΜΝ, ΝΧ, ΧΟ, ΟΛ, ΕΟ. ἴσόπλευρον ἄρα ἐστὶ καὶ τὸ
ΛΜΝΧΟ πεντάγωνον, καὶ δεκαγώνου ἡ ΕΟ εὐθεῖα.
καὶ ἀνεστάτωσαν ἄπὸ τῶν Ε, Ζ, Η, Θ, Κ σημείων
τῷ τοῦ κύκλου ἐπιπέδῳ πρὸς ὀρθὰς γωνίας
εὐθεῖαι αἱ ΕΠ, ΖΡ, ΗΣ, ΘΤ, ΚΥ ἴσαι οὖσαι τῇ ἐκ
τοῦ κέντρου τοῦ ΕΖΗΘΚ κύκλου, καὶ ἐπεζεύθθωσαν
αἱ ΠΡ, ΡΣ, ΣΤ, ΤΥ, ΥΠ, ΠΛ, ΛΡ, ΡΜ, ΜΣ, ΣΝ, ΝΤ, ΤΧ, ΧΥ,
ΥΟ, ΟΠ.}

{\greekfont Καὶ ἐπεὶ ἑκατέρα τῶν ΕΠ, ΚΥ τῷ α ὐτῷ ἐπιπέδῳ
πρὸς ὀρθάς ἐστιν, παράλληλος ἄρα ἐστὶν ἡ ΕΠ
τῇ ΚΥ. ἔστι δὲ α ὐτῇ καὶ ἴση; αἱ δὲ τὰς ἴσας
τε καὶ παραλλήλους ἐπιζευγνύουσαι ἐπὶ τὰ α ὐτὰ μέρη
εὐθεῖαι ἴσαι τε καὶ παράλληλοί εἰσιν. ἡ ΠΥ ἄρα τῇ
ΕΚ ἴση τε καὶ παράλληλός ἐστιν. πενταγώνου δὲ
ἰσοπλεύρου ἡ ΕΚ; πενταγώνου ἄρα ἰσοπλεύρου
καὶ ἡ ΠΥ τοῦ εἰς τὸν ΕΖΗΘΚ κύκλον ἐγγραφομένου.
διὰ τὰ α ὐτὰ δὴ καὶ ἑκάστη τῶν ΠΡ, ΡΣ, ΣΤ, ΤΥ πενταγώνου
ἐστὶν ἰσοπλεύρου τοῦ εἰς τὸν ΕΖΗΘΚ κύκλον ἐγγραφομένου;
ἰσόπλευρον ἄρα τὸ ΠΡΣΤΥ πεντάγωνον. καὶ ἐπεὶ ἑχαγώνου
μέν ἐστιν ἡ ΠΕ, δεκαγώνου δὲ ἡ ΕΟ, καί ἐστιν ὀρθὴ ἡ ὑπὸ
ΠΕΟ, πενταγώνου ἄρα ἐστὶν ἡ ΠΟ; ἡ γὰρ τοῦ πενταγώνου
πλευρὰ δύναται τήν τε τοῦ ἑχαγώνου καὶ τὴν τοῦ
δεκαγώνου τῶν εἰς τὸν α ὐτὸν κύκλον ἐγγραφομένων.
διὰ τὰ α ὐτὰ δὴ καὶ ἡ ΟΥ πενταγώνου ἐστὶ πλευρά.
ἔστι δὲ καὶ ἡ ΠΥ πενταγώνου; ἰσόπλευρον ἄρα ἐστὶ
τὸ ΠΟΥ τρίγωνον. διὰ τὰ α ὐτὰ δὴ καὶ ἕκαστον τῶν ΠΛΡ, ΡΜΣ, ΣΝΤ,
ΤΧΥ ἰσόπλευρόν ἐστιν. καὶ ἐπεὶ πενταγώνου ἐδείθθη ἑκατέρα
τῶν ΠΛ, ΠΟ, ἔστι δὲ καὶ ἡ ΛΟ πενταγώνου, ἰσόπλευρον ἄρα
ἐστὶ τὸ ΠΛΟ τρίγωνον. διὰ τὰ α ὐτὰ δὴ καὶ ἕκαστον τῶν
ΛΡΜ, ΜΣΝ, ΝΤΧ, ΧΥΟ τριγώνων ἰσόπλευρόν ἐστιν.}~\\~\\~\\~\\

\epsfysize=3.5in
\centerline{\epsffile{Book13/fig16g.eps}}

{\greekfont Εἰλήφθω
τὸ κέντρον τοῦ ΕΖΗΘΚ κύκλου τὸ Φ σημεῖον; καὶ ἀπὸ
τοῦ Φ τῷ τοῦ κύκλου ἐπιπέδῳ πρὸς ὀρθὰς
ἀνεστάτω ἡ ΦΩ, καὶ ἐκβεβλήσθω ἐπὶ τὰ
ἕτερα μέρη ὡς ἡ ΦΨ, καὶ ἀφῃρήσθω ἑχαγώνου μὲν
ἡ ΦΘ, δεκαγώνου δὲ ἑκατέρα τῶν ΦΨ, ΘΩ, καὶ ἐπεζεύθθωσαν αἱ ΠΩ, ΠΘ, ΥΩ, ΕΦ, ΛΦ, ΛΨ,
ΨΜ.}

{\greekfont Καὶ ἐπεὶ ἑκατέρα τῶν ΦΘ, ΠΕ τῷ τοῦ κύκλου
ἐπιπέδῳ πρὸς ὀρθάς ἐστιν, παράλληλος ἄρα ἐστὶν
ἡ ΦΘ τῇ ΠΕ. εἰσὶ δὲ καὶ ἴσαι; καὶ αἱ ΕΦ, ΠΘ
ἄρα ἴσαι τε καὶ παράλληλοί εἰσιν. ἑχαγώνου δὲ ἡ ΕΦ; 
ἑχαγώνου ἄρα καὶ ἡ ΠΘ. καὶ ἐπεὶ ἑχαγώνου
μέν ἐστιν ἡ ΠΘ, δεκαγώνου δὲ ἡ ΘΩ, καὶ ὀρθή ἐστιν
ἡ ὑπὸ ΠΘΩ γωνία, πενταγώνου ἄρα ἐστὶν ἡ ΠΩ.
διὰ τὰ α ὐτὰ δὴ καὶ ἡ ΥΩ πενταγώνου ἐστίν, ἐπειδήπερ,
ἐὰν ἐπιζεύχωμεν τὰς ΦΚ, ΘΥ, ἴσαι καὶ ἀπεναντίον
ἔσονται, καί ἐστιν ἡ ΦΚ ἐκ τοῦ κέντρου οὖσα
ἑχαγώνου.  ἑχαγώνου ἄρα καὶ ἡ ΘΥ. δεκαγώνου
δὲ ἡ ΘΩ, καὶ ὀρθὴ ἡ ὑπὸ ΥΘΩ; πενταγώνου
ἄρα ἡ ΥΩ. ἔστι δὲ καὶ ἡ ΠΥ πενταγώνου; ἰσόπλευρον ἄρα
ἐστὶ τὸ ΠΥΩ τρίγωνον. διὰ τὰ α ὐτὰ δὴ καὶ ἕκαστον
τῶν λοιπῶν τριγώνων, ὧν βάσεις μέν εἰσιν αἱ ΠΡ, ΡΣ,
ΣΤ, ΤΥ εὐθεῖαι, κορυφὴ δὲ τὸ Ω σημεῖον,  ἰσόπλευρόν
ἐστιν. πάλιν, ἐπεὶ ἑχαγώνου μὲν ἡ ΦΛ, δεκαγώνου
δὲ ἡ ΦΨ, καὶ ὀρθή ἐστιν ἡ ὑπὸ ΛΦΨ γωνία, πενταγώνου
ἄρα ἐστὶν ἡ ΛΨ. διὰ τὰ α ὐτὰ δὴ ἐὰν ἐπιζεύχωμεν τὴν
ΜΦ οὖσαν ἑχαγώνου, συνάγεται καὶ ἡ ΜΨ πενταγώνου.
ἔστι δὲ καὶ ἡ ΛΜ πενταγώνου; ἰσόπλευρον ἄρα ἐστὶ τὸ
ΛΜΨ τρίγωνον. ὁμοίως δὴ δειθθήσεται, ὅτι καὶ ἕκαστον
τῶν λοιπῶν τριγώνων, ὧν βάσεις μέν εἰσιν αἱ ΜΝ,
ΝΧ, ΧΟ, ΟΛ, κορυφὴ δὲ τὸ Ψ σημεὶον, ἰσόπλευρόν
ἐστιν. συνέσταται ἄρα εἰκοσάεδρον ὑπὸ εἴκοσι τριγώνων
ἰσοπλεύρων περιεχόμενον.}

{\greekfont Δεῖ δὴ α ὐτὸ καὶ σφαίρᾳ περιλαβεῖν τῇ δοθείσῃ
καὶ δεῖχαι, ὅτι ἡ τοῦ εἰκοσαέδρου πλευρὰ ἄλογός
ἐστιν ἡ καλουμένη ἐλάσσων.}

{\greekfont Ἐπεὶ γὰρ ἑχαγώνου ἐστὶν ἡ ΦΘ, δεκαγώνου δὲ ἡ ΘΩ,
ἡ ΦΩ ἄρα ἄκρον καὶ μέσον λόγον τέτμ ηται κατὰ τὸ Θ,
καὶ τὸ μεῖζον α ὐτῆς τμῆμά ἐστιν ἡ ΦΘ; ἔστιν
ἄρα ὡς ἡ ΩΦ πρὸς τὴν ΦΘ, οὕτως ἡ ΦΘ πρὸς
τὴν ΘΩ. ἴση δὲ ἡ μὲν ΦΘ τῇ ΦΕ, ἡ δὲ ΘΩ
τῇ ΦΨ; ἔστιν ἄρα ὡς ἡ ΩΦ πρὸς τὴν ΦΕ, οὕτως
ἡ ΕΦ πρὸς τὴν ΦΨ. καί εἰσιν ὀρθαὶ αἱ ὑπὸ
ΩΦΕ, ΕΦΨ γωνίαι; ἐὰν ἄρα ἐπιζεύχωμεν τὴν ΕΩ εὐθεὶαν,
ὀρθὴ ἔσται ἡ ὑπὸ ΨΕΩ γωνία διὰ τὴν ὁμοιότητα
τῶν ΨΕΩ, ΦΕΩ τριγώνων. διὰ τὰ α ὐτὰ δὴ ἐπεί
ἐστιν ὡς ἡ ΩΦ πρὸς τὴν ΦΘ, οὕτως ἡ ΦΘ πρὸς
τὴν ΘΩ, ἴση δὲ ἡ μὲν ΩΦ τῇ ΨΘ, ἡ δὲ ΦΘ τῇ ΘΠ, ἔστιν ἄρα ὡς ἡ ΨΘ πρὸς τὴν ΘΠ, οὕτως ἡ ΠΘ πρὸς τὴν ΘΩ. καὶ διὰ τοῦτο
πάλιν ἐὰν ἐπιζεύχωμεν τὴν ΠΨ, ὀρθὴ ἔσται ἡ πρὸς τῷ
Π γωνία; τὸ ἄρα ἐπὶ τῆς ΨΩ γραφόμενον ἡμικύκλιον
ἥχει καὶ
δὶα τοῦ Π. καὶ ἐὰν μενούσης τῆς ΨΩ περιενεθθὲν
τὸ ἡμικύκλιον εἰς τὸ α ὐτὸ πάλιν ἀποκατασταθῇ,
ὅθεν ἤρχατο φέρεσθαι, ἥχει καὶ διὰ τοῦ Π καὶ 	τῶν
λοιπῶν σημείων τοῦ εἰκοσαέδρου, καὶ ἔσται σφαίρᾳ
περιειλημμένον τὸ εἰκοσάεδρον. λέγω δή, ὅτι καὶ τῇ
δοθείσῃ. τετμήσθω γὰρ ἡ ΦΘ δίθα κατὰ τὸ α. καὶ ἐπεὶ
εὐθεῖα γραμμὴ ἡ ΦΩ ἄκρον καὶ μέσον λόγον τέτμ ηται
κατὰ τὸ Θ, καὶ τὸ ἔλασσον α ὐτῆς τμῆμά ἐστιν ἡ ΩΘ,  ἡ ἄρα
ΩΘ προσλαβοῦσα 
τὴν ἡμίσειαν τοῦ μείζονος τμήματος τὴν Θα πενταπλάσιον
δύναται τοῦ ἀπὸ τῆς
ἡμισείας τοῦ μείζονος τμήματος; πενταπλάσιον ἄρα
ἐστὶ τὸ ἀπὸ τῆς Ωα τοῦ ἀπὸ τῆς αΘ. καί ἐστι τῆς
μὲν Ωα διπλῆ ἡ ΩΨ, τὴς δὲ αΘ διπλῆ ἡ ΦΘ;
πενταπλάσιον ἄρα ἐστὶ τὸ ἀπὸ τῆς ΩΨ τοῦ  ἀπὸ
τῆς ΘΦ. καὶ ἐπεὶ τετραπλῆ ἐστιν ἡ ΑΓ τῆς ΓΒ, πενταπλῆ
ἄρα ἐστὶν ἡ ΑΒ τῆς ΒΓ. ὡς δὲ ἡ ΑΒ πρὸς τὴν ΒΓ,
οὕτως τὸ ἀπὸ τῆς ΑΒ πρὸς τὸ ἀπὸ τῆς ΒΔ;
πενταπλάσιον ἄρα ἐστὶ τὸ ἀπὸ τῆς ΑΒ τοῦ ἀπὸ
τῆς ΒΔ. ἐδείθθη δὲ καὶ τὸ ἀπὸ τῆς ΩΨ πενταπλάσιον
τοῦ ἀπὸ τῆς ΦΘ. καί ἐστιν ἴση ἡ ΔΒ τῇ ΦΘ;
ἑκατέρα γὰρ α ὐτῶν ἴση ἐστὶ τῇ ἐκ τοῦ κέντρου τοῦ
ΕΖΗΘΚ κύκλου; ἴση ἄρα καὶ ἡ ΑΒ τῇ ΨΩ. καί
ἐστιν ἡ ΑΒ ἡ τῆς δοθείσης σφαίρας διάμετρος; καὶ
ἡ ΨΩ ἄρα ἴση ἐστὶ τῇ τῆς δοθείσης σφαίρας
διαμέτρῳ; τῇ ἄρα δοθείσῃ σφαίρᾳ περιείληπται τὸ
εἰκοσάεδρον.}

{\greekfont Λέγω δή, ὅτι ἡ τοῦ εἰκοσαέδρου πλευρὰ ἄλογός
ἐστιν ἡ καλουμένη ἐλάττων. ἐπεὶ γὰρ ῥητή ἐστιν ἡ τῆς
σφαίρας διάμετρος, καί ἐστι δυνάμει πενταπλασίων
τῆς ἐκ τοῦ κέντρου τοῦ ΕΖΗΘΚ κύκλου,
ῥητὴ ἄρα ἐστὶ καὶ ἡ ἑκ τοῦ κέντρου τοῦ ΕΖΗΘΚ κύκλου;
ὥστε καὶ ἡ διάμετρος α ὐτοῦ ῥητή ἐστιν. ἐὰν δὲ εἰς
κύκλον ῥητὴν ἔθοντα τὴν διάμετρον πεντάγωνον ἰσόπλευρον ἐγγραφῃ,
ἡ τοῦ πενταγώνου πλευρὰ ἄλογός ἐστιν ἡ καλουμένη ἐλάττων.
ἡ δὲ τοῦ ΕΖΗΘΚ
πενταγώνου πλευρὰ ἡ τοῦ εἰκοσαέδρου ἐστίν. ἡ ἄρα τοῦ
είκοσαέδρου πλευρὰ ἄλογός ἐστιν ἡ καλουμένη
ἐλάττων.}}

\ParallelRText{
\begin{center}
{\large Proposition 16}
\end{center}

To construct an icosahedron, and to enclose (it) in a sphere, like the
aforementioned figures, and to show that the side of the
icosahedron is that irrational (straight-line) called minor.

\epsfysize=1.5in
\centerline{\epsffile{Book13/fig16ae.eps}}

Let the diameter $AB$ of the given sphere be laid out, and let it be
cut at $C$ such that $AC$ is four times $CB$ [Prop.~6.10]. And
let the semi-circle $ADB$ be drawn on $AB$. And let the straight-line $CD$ be drawn from $C$ at right-angles to $AB$. And let
$DB$ be joined. And let the circle $EFGHK$ be set down,
and let its radius be equal to $DB$.  And let the equilateral and
equiangular pentagon $EFGHK$ be inscribed in circle
$EFGHK$ [Prop.~4.11]. And let the circumferences $EF$, $FG$,
$GH$, $HK$, and $KE$ be cut in half at points $L$, $M$,
$N$, $O$, and $P$ (respectively). And let $LM$, $MN$, $NO$,
$OP$, $PL$, and $EP$ be joined. Thus, pentagon $LMNOP$
is also equilateral, and $EP$ (is) the side of the decagon (inscribed in the circle).
And let the straight-lines $EQ$, $FR$, $GS$, $HT$, and $KU$, which
are equal to the radius of circle $EFGHK$,  have
been set up at right-angles to the plane of the circle, at points $E$, $F$, $G$, $H$, and $K$ (respectively). And let $QR$, $RS$, $ST$, $TU$,
$UQ$, $QL$, $LR$, $RM$, $MS$, $SN$, $NT$, $TO$,
$OU$, $UP$, and $PQ$ be joined.

And since $EQ$ and $KU$ are each at right-angles to the same plane, $EQ$ is thus parallel to $KU$ [Prop.~11.6]. And it is also equal to it. And straight-lines joining equal and parallel (straight-lines) on the same side are
(themselves) equal and parallel [Prop.~1.33]. Thus, $QU$
is equal and parallel to $EK$. And $EK$ (is the side) of an equilateral
pentagon (inscribed in circle $EFGHK$). Thus, $QU$ (is) also  the
side of an equilateral pentagon inscribed in circle $EFGHK$.  So, for the
same (reasons), $QR$, $RS$, $ST$, and $TU$ are also  the sides of
an equilateral pentagon inscribed in circle $EFGHK$. Pentagon
$QRSTU$ (is) thus equilateral. And side $QE$
is (the side) of a hexagon (inscribed in circle $EFGHK$),
and $EP$ (the side) of a decagon, and (angle) $QEP$ is a right-angle, 
thus $QP$ is  (the side) of a pentagon (inscribed in the same circle).
For the square on the side of a pentagon is (equal to the sum of)
the (squares) on (the sides of) a hexagon and a decagon inscribed in the
same circle  [Prop.~13.10]. So, for the same (reasons), $PU$
is also the side of a pentagon. And $QU$ is also (the
side) of a pentagon. Thus, triangle $QPU$ is equilateral. So, for the
same (reasons), (triangles) $QLR$, $RMS$, $SNT$, and $TOU$
are each also equilateral. And since $QL$ and $QP$ were each shown
(to be the sides) of a pentagon,  and $LP$ is  also (the side) of a pentagon, triangle $QLP$ is thus equilateral. So,
for the same (reasons), triangles $LRM$, $MSN$, $NTO$, and
$OUP$ are each also equilateral. 

\epsfysize=3.5in
\centerline{\epsffile{Book13/fig16e.eps}}

Let the center, point $V$, of circle $EFGHK$ be found [Prop.~3.1]. 
And let $VZ$ be set up, at (point) $V$,  at right-angles to the plane of the
circle. And let it be produced on the other side (of the circle), 
like $VX$. And let $VW$ be cut off 
(from $XZ$ so as to be equal to the side) of a hexagon, and each of $VX$ and 
$WZ$ (so as to be equal to the side) of a decagon. And let
$QZ$,  $QW$, $UZ$, $EV$, $LV$, $LX$, and
$XM$ be joined.

And since $VW$ and $QE$ are each at right-angles to the plane of the
circle, $VW$ is thus parallel to $QE$ [Prop.~11.6].  And they are also
equal. $EV$ and $QW$ are thus equal and parallel (to one another) [Prop.~1.33]. And $EV$ (is the side) of a hexagon. Thus, $QW$
(is) also (the side) of a hexagon. And since $QW$ is (the side) of a hexagon, 
and $WZ$ (the side) of a decagon, and angle $QWZ$ is a right-angle
[Def.~11.3, Prop.~1.29], 
$QZ$ is thus (the side) of a pentagon [Prop.~13.10]. So, for the same
(reasons), $UZ$ is also (the side) of a pentagon---inasmuch as, if we join
$VK$ and $WU$ then they will be equal and opposite. 
And $VK$, being (equal)
to the radius (of the circle), is (the side) of a hexagon [Prop.~4.15~corr.].
Thus, $WU$ (is) also the side of a hexagon.
 And $WZ$
(is the side) of a decagon, and (angle) $UWZ$ (is) a right-angle. 
Thus, $UZ$ (is the side) of a pentagon [Prop.~13.10]. And $QU$ is
also (the side) of a pentagon. Triangle $QUZ$ is thus equilateral. So,
for the same (reasons), each of the remaining triangles, whose bases
are the straight-lines $QR$, $RS$, $ST$, and $TU$, and apexes the point
$Z$, are also equilateral. Again, since $VL$ (is the side) of a hexagon, and $VX$ (the side) of a decagon, and angle $LVX$ is a right-angle, 
$LX$ is thus (the side) of a pentagon [Prop.~13.10]. So, for the same
(reasons), if we join $MV$, which is (the side) of a hexagon, $MX$
is also inferred (to be the side) of a pentagon. And $LM$ is also (the side)
of a pentagon. Thus, triangle $LMX$ is equilateral. So, similarly,
it can be shown that each of the remaining triangles, whose bases
are the (straight-lines) $MN$, $NO$, $OP$, and $PL$, and apexes the
point $X$,  are also equilateral. Thus, an icosahedron contained by
twenty equilateral triangles has been constructed.

So, it is also necessary to enclose it in the given sphere, and to show
that the side of the icosahedron is that irrational (straight-line) called minor.

For, since $VW$ is (the side) of a hexagon, and $WZ$ (the side) of
a decagon, $VZ$ has thus been cut in extreme and mean
ratio at $W$, and $VW$ is its greater piece  [Prop.~13.9]. Thus, as
$ZV$ is to $VW$, so $VW$ (is) to $WZ$. And $VW$ (is) equal
to $VE$, and $WZ$ to $VX$.  Thus, as $ZV$ is to $VE$,
so $EV$ (is) to $VX$. And angles $ZVE$ and
$EVX$ are right-angles. Thus, if we join straight-line $EZ$ then angle
$XEZ$ will be a right-angle, on account of the similarity of triangles
$XEZ$ and $VEZ$.  [Prop.~6.8]. So, for the same (reasons), since
as $ZV$ is to $VW$, so $VW$ (is) to $WZ$,  and $ZV$ (is) equal to
$XW$, and $VW$ to $WQ$, thus as $XW$ is to $WQ$, so $QW$
(is) to $WZ$. And, again, on account of this,  if we join $QX$ then the
angle at $Q$ will be a right-angle [Prop.~6.8]. Thus, the semi-circle
drawn on $XZ$ will also pass through $Q$ [Prop.~3.31]. And if $XZ$ remains
fixed, and the semi-circle is carried around, and again established at
the same (position) from which it began to be moved, then it will
also pass through (point) $Q$, and (through) the remaining (angular) points
of the icosahedron. And the icosahedron will be enclosed by a sphere.
 So, I say that (it is) also (enclosed) by the given (sphere). 
 For let $VW$ be cut in half at $a$. And since the straight-line $VZ$ has been cut in extreme and mean ratio at $W$, and  $ZW$ is its lesser piece,
 then the square on $ZW$ added to half of the greater piece, $Wa$, is
 five times the (square) on half of the greater piece [Prop.~13.3]. 
 Thus, the (square) on $Za$ is five times the (square) on $aW$. And
 $ZX$ is double $Za$, and $VW$ double $aW$. Thus,
 the (square) on $ZX$ is five times the (square) on $WV$. And since
 $AC$ is four times $CB$, $AB$ is thus five times $BC$. And as
 $AB$ (is) to $BC$, so the (square) on $AB$ (is) to the
 (square) on $BD$ [Prop.~6.8, Def.~5.9]. Thus, the (square)
 on $AB$ is five times the (square) on $BD$. And the (square) on 
 $ZX$ was also shown (to be) five times the (square) on $VW$. 
 And $DB$ is equal to $VW$. For each of them is equal to the
 radius of circle $EFGHK$.  Thus, $AB$ (is) also equal to
 $XZ$.  And $AB$ is the diameter of the given sphere. Thus,
 $XZ$ is equal to the diameter of the given sphere. Thus, the icosahedron
 has been enclosed by the given sphere.
 
 So, I say that the side of the icosahedron is that irrational (straight-line)
 called minor. For since the diameter of the sphere is rational, and the
 square on it is five times the (square) on the radius of circle
 $EFGHK$, the radius of circle $EFGHK$ is thus also rational.
 Hence, its diameter is also rational. And if an equilateral
 pentagon is inscribed in a circle having a rational diameter then
 the side of the pentagon is that irrational (straight-line) called minor [Prop.~13.11]. And the side of pentagon $EFGHK$ is (the side)
 of the icosahedron. Thus, the side of the icosahedron is that irrational
 (straight-line) called minor.}
\end{Parallel}

\begin{Parallel}{}{}
\ParallelLText{
\begin{center}
{\large {\greekfont Πόρισμα}.}
\end{center}\vspace*{-7pt}

{\greekfont Ἐκ δὴ τούτου φανερόν, ὅτι ἡ τῆς σφαίρας διάμετρος δυνάμει
πενταπλασίων ἐστὶ τῆς ἐκ τοῦ κέντρου τοῦ κύκλου,
ἀφ οὗ τὸ εἰκοσάεδρον ἀναγέγραπται, καὶ ὅτι
ἡ τῆς σφαίρας διάμετρος σύγκειται ἔκ τε τῆς τοῦ
ἑχαγώνου καὶ δύο τῶν τοῦ δεκαγώνου τῶν εἰς
τὸν α ὐτὸν κύκλον ἐγγραφομένων. ὅπερ ἔδει
δεῖχαι.}}

\ParallelRText{
\begin{center}
{\large Corollary}
\end{center}

So,  (it is) clear, from this, that the square on the diameter of the sphere
is five times the square on the radius of the circle from which the icosahedron has been described,
and that the the diameter of the sphere is the sum of (the side) of the hexagon, 
and two of (the sides) of the decagon, inscribed in the same circle.$^\dag$ }
\end{Parallel}
{\footnotesize\noindent$^\dag$ If the radius of the sphere is unity then the
radius of the circle is $2/\sqrt{5}$, and the sides of the hexagon, decagon,
and pentagon/icosahedron are $2/\sqrt{5}$, $1-1/\sqrt{5}$, and $(1/\sqrt{5})\,\sqrt{10-2\,\sqrt{5}}$, respectively.}

%%%%
%13.17
%%%%
\pdfbookmark[1]{Proposition 13.17}{pdf13.17}
\begin{Parallel}{}{}
\ParallelLText{
\begin{center}
{\large \ggn{17}.}
\end{center}\vspace*{-7pt}

{\greekfont Δωδεκάεδρον συστήσασθαι καὶ σφαίρᾳ περιλαβεῖν, ᾗ καὶ τὰ
προειρημένα σθήματα, καὶ δεῖχαι, ὅτι ἡ τοῦ δωδεκαέδρου
πλευρὰ ἄλογός ἐστιν ἡ καλουμένη ἀποτομή.}

{\greekfont Ἐκκείσθωσαν τοῦ προειρημένου κύβου δύο ἐπίπεδα πρὸς
ὀρθὰς ἀλλήλοις τὰ ΑΒΓΔ, ΓΒΕΖ, καὶ τετμήσθω ἑκάστη τῶν
ΑΒ, ΒΓ, ΓΔ, ΔΑ, ΕΖ, ΕΒ, ΖΓ πλευρῶν δίθα κατὰ τὰ Η, Θ, Κ, Λ, Μ, Ν, Χ,
καὶ ἐπεζεύθθωσαν αἱ ΗΚ, ΘΛ, ΜΘ, ΝΧ, καὶ τετηήσθω ἑκάστη τῶν
ΝΟ, ΟΧ, ΘΠ ἄκρον καὶ μέσον λόγον κατὰ τὰ Ρ, Σ, Τ σημεῖα,
καὶ ἔστω α ὐτῶν μείζονα τμήματα τὰ ΡΟ, ΟΣ, ΤΠ, καὶ ἀνεστάτωσαν
ἀπὸ τῶν Ρ, Σ, Τ σημείων τοῖς τοῦ κύβου ἐπιπέδοις
πρὸς ὀρθὰς ἐπὶ τὰ ἐκτὸς μέρη τοῦ κύβου αἱ ΡΥ, ΣΦ, ΤΘ,
καὶ κείσθωσαν ἴσαι ταῖς ΡΟ, ΟΣ, ΤΠ, καὶ ἐπεζεύθθωσαν
αἱ ΥΒ, ΒΘ, ΘΓ, ΓΦ, ΦΥ.}\\~\\~\\

\epsfysize=3in
\centerline{\epsffile{Book13/fig17g.eps}}

{\greekfont Λέγω, ὅτι τὸ ΥΒΘΓΦ πεντάγωνον
ἰσόπλευρόν τε καὶ ἐν ἑνὶ ἐπιπέδῳ καὶ ἔτι ἰσογώνιόν ἐστιν.
ἐπεζεύθθωσαν γὰρ αἱ ΡΒ, ΣΒ, ΦΒ. καὶ ἐπεὶ εὐθεῖα ἡ ΝΟ ἄκρον
καὶ μέσον λόγον τέτμ ηται κατὰ τὸ Ρ, καὶ τὸ μεῖζον τμῆμά
ἐστιν ἡ ΡΟ, τὰ ἄρα ἀπὸ τῶν ΟΝ, ΝΡ τριπλάσιά ἐστι τοῦ 
ἀπὸ τῆς ΡΟ. ἴση δὲ ἡ μὲν ΟΝ τῇ ΝΒ, ἡ δὲ ΟΡ τῇ ΡΥ; τὰ
ἄρα ἀπὸ τῶν ΒΝ, ΝΡ τριπλάσιά ἐστι τοῦ ἀπὸ τῆς ΡΥ.
τοῖς δὲ ἀπὸ τῶν ΒΝ, ΝΡ τὸ ἀπὸ τῆς ΒΡ ἐστιν
ἴσον; τὸ ἄρα ἀπὸ τῆς ΒΡ τριπλάσιόν ἐστι τοῦ ἀπὸ τῆς ΡΥ;
ὥστε τὰ ἀπὸ τῶν ΒΡ, ΡΥ τετραπλάσιά ἐστι τοῦ ἀπὸ τῆς ΡΥ.
τοῖς δὲ ἀπὸ τῶν ΒΡ, ΡΥ ἴσον ἐστι τὸ ἀπὸ τῆς ΒΥ; τὸ
ἄρα ἄπὸ τῆς ΒΥ τετραπλάσιόν ἐστι τοῦ ἀπὸ τῆς ΥΡ;  διπλῆ ἄρα
ἐστὶν ἡ ΒΥ τῆς ΡΥ. ἔστι δὲ καὶ ἡ ΦΥ τῆς ΥΡ διπλῆ, 
ἐπειδήπερ καὶ ἡ ΣΡ τῆς ΟΡ, τουτέστι τῆς ΡΥ, ἐστι διπλῆ;
ἴση ἄρα ἡ ΒΥ τῇ ΥΦ. ὁμοίως δὴ δειθθήσεται, ὅτι
καὶ ἑκάστη  τῶν ΒΘ, ΘΓ, ΓΦ ἑκατέρᾳ τῶν ΒΥ, ΥΦ
ἐστιν ἴση. ἰσόπλευρον ἄρα ἐστὶ τὸ ΒΥΦΓΘ πεντάγωνον. λέγω
δή, ὅτι καὶ ἐν ἑνί  ἐστιν ἐπιπέδῳ. ἤθθω γὰρ ἀπὸ τοῦ Ο
ἑκατέρᾳ τῶν ΡΥ, ΣΦ παράλληλος ἐπὶ τὰ ἐκτὸς τοῦ
κύβου μέρη ἡ ΟΨ, καὶ ἐπεζεύθθωσαν αἱ ΨΘ, ΘΘ;
λέγω,  ὅτι ἡ ΨΘΘ εὐθεῖά ἐστιν. ἐπεὶ γὰρ ἡ ΘΠ ἄκρον καὶ
μέσον λόγον τέτμ ηται κατὰ τὸ Τ, καὶ τὸ μεῖζον α ὐτῆς
τμῆμά ἐστιν ἡ ΠΤ, ἔστιν ἄρα ὡς ἡ ΘΠ πρὸς τὴν 
ΠΤ, οὕτως ἡ ΠΤ πρὸς τὴν ΤΘ.  ἴση δὲ ἡ μὲν ΘΠ
τῇ ΘΟ, ἡ δὲ ΠΤ ἑκατέρᾳ τῶν ΤΘ, ΟΨ; ἔστιν ἄρα ὡς
ἡ ΘΟ πρὸς τὴν ΟΨ, οὕτως ἡ ΘΤ πρὸς τὴν ΤΘ. καί ἐστι
παράλληλος ἡ μὲν ΘΟ τῇ ΤΘ; ἑκατέρα γὰρ α ὐτῶν
τῷ ΒΔ ἐπιπέδῳ πρὸς ὀρθάς ἐστιν; ἡ δὲ ΤΘ
τῇ ΟΨ; ἑκατέρα γὰρ α ὐτῶν τῷ ΒΖ ἐπιπέδῳ
πρὸς ὀρθάς ἐστιν.  ἐὰν δὲ δύο τρίγωνα συντεθῇ κατὰ μίαν
γωνίαν, ὡς τὰ ΨΟΘ, ΘΤΘ, τὰς δύο πλευρὰς ταῖς δυνὶν ἀνάλογον
ἔθοντα, ὥστε τὰς ὁμολόγους α ὐτῶν πλευρὰς καὶ παραλλήλους
εἶναι, αἱ λοιπαὶ εὐθεῖαι ἐπ εὐθείας ἔσονται; ἐπ εὐθείας
ἄρα ἐστὶν ἡ ΨΘ τῇ ΘΘ. πᾶσα δὲ εὐθεῖα ἐν ἑνί
ἐστιν ἐπιπέδῳ; ἐν ἑνὶ ἄρα ἐπιπέδῳ ἐστὶ τὸ ΥΒΘΓΦ πεντάγωνον.}

{\greekfont Λέγω δή, ὅτι καὶ ἰσογώνιόν ἐστιν.}

{\greekfont Ἐπεὶ γὰρ εὐθεῖα γραμμὴ ἡ ΝΟ ἄκρον καὶ μέσον λόγον
τέτμ ηται κατὰ τὸ Ρ, καὶ τὸ μεῖζον τμῆμά
ἐστιν ἡ ΟΡ [ἔστιν ἄρα ὡς συναμφότερος ἡ ΝΟ, ΟΡ πρὸς
τὴν ΟΝ, οὕτως ἡ ΝΟ πρὸς τὴν ΟΡ], ἴση δὲ ἡ ΟΡ
τῇ ΟΣ [ἔστιν ἄρα ὡς ἡ ΣΝ πρὸς τὴν ΝΟ, οὕτως
ἡ ΝΟ πρὸς τὴν ΟΣ], ἡ ΝΣ ἄρα ἄκρον καὶ μέσον λόγον
τέτμ ηται κατὰ τὸ Ο, καὶ τὸ μεῖζον τμῆμά
ἐστιν ἡ ΝΟ; τὰ ἄρα ἀπὸ τῶν ΝΣ, ΣΟ τριπλάσιά ἐστι τοῦ ἀπὸ
τῆς ΝΟ. ἴση δὲ ἡ μὲν ΝΟ τῇ ΝΒ, ἡ δὲ ΟΣ τῇ ΣΦ; τὰ
ἄρα ἀπὸ τῶν ΝΣ, ΣΦ τετράγωνα τριπλάσιά ἐστι τοῦ ἀπὸ τῆς
ΝΒ; ὥστε τὰ ἀπὸ τῶν ΦΣ, ΣΝ, ΝΒ τετραπλάσιά ἐστι  τοῦ ἀπὸ
τῆς ΝΒ. τοῖς δὲ ἀπὸ τῶν ΣΝ, ΝΒ ἴσον ἐστὶ τὸ ἀπὸ
τῆς ΣΒ; τὰ ἄρα ἀπὸ τῶν ΒΣ, ΣΦ, τουτέστι τὸ ἀπὸ τῆς
ΒΦ [ὀρθὴ γὰρ ἡ ὑπὸ ΦΣΒ γωνία], τετραπλάσιόν
ἐστι τοῦ ἀπὸ τῆς ΝΒ; διπλῆ ἄρα ἐστὶν ἡ ΦΒ τῆς ΒΝ.
ἔστι δὲ καὶ ἡ ΒΓ τῆς ΒΝ διπλῆ; ἴση  ἄρα ἐστὶν ἡ ΒΦ τῇ ΒΓ. καὶ
ἐπεὶ δύο αἱ ΒΥ, ΥΦ δυσὶ ταῖς ΒΘ, ΘΓ ἴσαι εἰσίν, καὶ
βάσις ἡ ΒΦ βάσει τῇ ΒΓ ἴση, γωνία ἄρα ἡ ὑπὸ ΒΥΦ
γωνίᾳ τῇ ὑπὸ ΒΘΓ ἐστιν ἴση. ὁμοίως δὴ δείχομεν,
ὅτι καὶ ἡ ὑπὸ ΥΦΓ γωνία ἴση ἐστὶ τῇ ὑπὸ ΒΘΓ;
αἱ ἄρα ὑπὸ ΒΘΓ, ΒΥΦ, ΥΦΓ τρεῖς γωνίαι ἴσαι
ἀλλήλαις εἰσίν. ἐὰν δὲ πενταγώνου ἰσοπλεύρου αἱ 
τρεῖς γωνίαι ἴσαι ἀλλήλαις ὦσιν, ἰσογώνιον ἔσται
τὸ πεντάγωνον; ἰσογώνιον ἄρα ἐστὶ τὸ ΒΥΦΓΘ πεντάγωνον.
ἐδείθθη δὲ καὶ ἰσόπλευρον;
τὸ ἄρα ΒΥΦΓΘ πεντάγωνον ἰσόπλευρόν ἐστι καὶ ἰσογώνιον,
καί ἐστιν ἐπὶ μιᾶς τοῦ κύβου πλευρᾶς τῆς ΒΓ. ἐὰν
ἄρα ἐφ ἑκάστης τῶν τοῦ κύβου δώδεκα
πλευρῶν τὰ α ὐτὰ κατασκευάσωμεν, συσταθήσεταί τι σχῆμα
στερεὸν ὑπὸ δώδεκα πενταγώνων ἰσοπλεύρων
τε καὶ ἰσογωνίων περιεχόμενον, ὃ καλεῖται δωδεκάεδρον.}

{\greekfont Δεῖ δὴ α ὐτὸ καὶ σφαίρᾳ περιλαβεῖν τῇ δοθείσῃ καὶ δεῖχαι,
ὅτι ἡ τοῦ δωδεκαέδρου πλευρὰ ἄλογός ἐστιν ἡ καλουμένη 
ἀποτομή.}

{\greekfont Ἐκβεβλήσθω γὰρ ἡ ΨΟ, καὶ ἔστω ἡ ΨΩ; συμβάλλει ἄρα
ἡ ΟΩ τῇ τοῦ κύβου διαμέτρῳ, καὶ δίθα
τέμνουσιν ἀλλήλας; τοῦτο γὰρ δέδεικται ἐν τῷ
παρατελεύτῳ θεωρήματι  τοῦ ἑνδεκάτου βιβλίου. τεμνέτωσαν
κατὰ τὸ Ω; τὸ Ω ἄρα κέντρον ἐστὶ τῆς σφαίρας τῆς περιλαμβανούσης
τὸν κύβον,
καὶ ἡ ΩΟ ἡμίσεια τῆς πλευρᾶς τοῦ κύβου.
 ἐπεζεύθθω δὴ ἡ ΥΩ. καὶ ἐπεὶ εὐθεῖα γραμμὴ ἡ ΝΣ
ἄκρον καὶ μέσον λόγον τέτμ ηται κατὰ τὸ Ο, καὶ
τὸ μεῖζον α ὐτῆς τμῆμά ἐστιν ἡ ΝΟ, τὰ ἄρα ἀπὸ τῶν ΝΣ, ΣΟ
τριπλάσιά ἐστι τοῦ ἀπὸ τῆς ΝΟ. ἴση δὲ ἡ μὲν ΝΣ τῇ ΨΩ,
ἐπειδήπερ καὶ ἡ μὲν ΝΟ τῇ ΟΩ ἐστιν ἴση, ἡ δὲ ΨΟ τῇ
ΟΣ. ἀλλὰ μὴν
καὶ ἡ ΟΣ τῇ ΨΥ, ἐπεὶ καὶ τῇ ΡΟ; τὰ ἄρα ἀπὸ τῶν
 ΩΨ, ΨΥ τριπλάσιά ἐστι τοῦ ἀπὸ τῆς ΝΟ. τοῖς
δὲ ἀπὸ τῶν ΩΨ, ΨΥ ἴσον ἐστὶ τὸ ἀπὸ τῆς ΥΩ; τὸ
ἄρα ἀπὸ τῆς ΥΩ τριπλάσιόν ἐστι τοῦ ἀπὸ τῆς ΝΟ.
ἔστι δὲ καὶ ἡ ἐκ τοῦ κέντρου τῆς σφαίρας τῆς περιλαμβανούσης
τὸν κύβον δυνάμει τριπλασίων τῆς ἡμισείας
τῆς τοῦ κύβου πλευρᾶς; προδέδεικται γὰρ κύβον συστήσασθαι
καὶ σφαίρᾳ περιλαβεῖν καὶ δεῖχαι, ὅτι ἡ τῆς σφαίρας
διάμετρος δυνάμει τριπλασίων ἐστὶ τῆς πλευρᾶς τοῦ κύβου.
εἰ δὲ ὅλη τῆς ὅλης, καὶ [ἡ] ἡμίσεια τῆς ἡμισείας; καί ἐστιν
ἡ ΝΟ ἡμίσεια τῆς
τοῦ κύβου πλευρᾶς; ἡ ἄρα
ΥΩ ἴση ἐστὶ τῇ ἐκ τοῦ κέντρου τῆς σφαίρας τῆς περιλαμβανούσης
τὸν κύβον. καί ἐστι τὸ Ω κέντρον τῆς σφαίρας τῆς περιλαμβανούσης
τὸν κύβον; τὸ Υ ἄρα σημεῖον πρὸς τῇ ἐπιφανείᾳ ἐστι τῆς
σφαίρας. ὁμοίως δὴ δείχομεν, ὅτι καὶ 
ἑκάστη τῶν λοιπῶν γωνιῶν τοῦ δωδεκαέδρου πρὸς τῇ
ἐπιφανείᾳ ἐστὶ τῆς
σφαίρας; περιείληπται ἄρα τὸ δωδεκαέδρον τῇ δοθείσῃ σφαίρᾳ.}

{\greekfont Λέγω δή, ὅτι ἡ τοῦ δωδεκαέδρου πλευρὰ ἄλογός
ἐστιν ἡ καλουμένη ἀποτομή.}

{\greekfont Ἐπεὶ γὰρ τῆς ΝΟ ἄκρον καὶ μέσον λόγον τετμ ημένης τὸ
μεῖζον τμῆμά ἐστιν ὁ ΡΟ, τῆς δὲ ΟΧ ἄκρον καὶ μέσον
λόγον τετμ ημένης τὸ μεῖζον τμῆμά ἐστιν ἡ ΟΣ, ὅλης ἄρα
τῆς ΝΧ ἄκρον καὶ μέσον λόγον
τεμνομένης τὸ μεῖζον τμῆμά ἐστιν ἡ ΡΣ. [οἷον ἐπεί
ἐστιν ὡς ἡ ΝΟ πρὸς τὴν ΟΡ, ἡ ΟΡ πρὸς τὴν ΡΝ, καὶ
τὰ διπλάσια; τὰ γὰρ μέρη τοῖς ἰσάκις πολλαπλασίοις
τὸν α ὐτὸν ἔχει λόγον; ὡς ἄρα ἡ ΝΧ πρὸς τὴν ΡΣ, οὕτως
ἡ ΡΣ πρὸς συναμφότερον τὴν ΝΡ, ΣΧ. μείζων δὲ ἡ ΝΧ
τῆς ΡΣ; μείζων ἄρα καὶ ἡ ΡΣ συναμφοτέρου τῆς ΝΡ, ΣΧ; ἡ ΝΧ
ἄρα ἄκρον  καὶ μέσον λόγον τέτμ ηται, καὶ
  τὸ μεῖζον  α ὐτῆς τμῆμά
ἐστιν ἡ ΡΣ.] ἴση δὲ ἡ ΡΣ τῇ ΥΦ; τῆς ἄρα ΝΧ ἄκρον
καὶ μέσον λόγον τεμνομένης τὸ μεῖζον τμῆμά ἐστιν ἡ
ΥΦ. καὶ ἐπεὶ ῥητή ἐστιν τῆς σφαίρας διάμετρος
καί ἐστι δυνάμει τριπλασίων τῆς τοῦ κύβου πλευρᾶς, ῥητὴ
ἄρα ἐστὶν ἡ ΝΧ πλευρὰ οὖσα τοῦ κύβου. ἐὰν δὲ
ῥητὴ γραμμὴ ἄκρον καὶ μέσον λόγον τμ ηθῇ, ἑκάτερον
τῶν τμ ημάτων ἄλογός ἐστιν ἀποτομή.}

{\greekfont Ἡ ΥΦ ἄρα πλευρὰ οὖσα τοῦ δωδεκαέδρου ἄλογός
ἐστιν ἀποτομή.}}

\ParallelRText{
\begin{center}
{\large Proposition 17}
\end{center}

To construct a dodecahedron, and to enclose (it) in a sphere, like the
aforementioned figures, and to show that the side of the dodecahedron 
is that irrational (straight-line) called an apotome.

Let two planes of the aforementioned cube [Prop. 13.15], $ABCD$ and $CBEF$,  (which are) at right-angles to one another,
be laid out. And let the sides $AB$, $BC$, $CD$, $DA$, $EF$,
$EB$, and $FC$  have each been cut in half at points $G$, $H$, $K$, $L$,
$M$, $N$, and $O$ (respectively). And let $GK$, $HL$, $MH$, and
$NO$ be joined. And let $NP$, $PO$, and $HQ$ have each been
cut in extreme and mean ratio at points $R$, $S$, and $T$ (respectively).
And let their greater pieces be $RP$, $PS$, and $TQ$ (respectively).
And let $RU$, $SV$, and $TW$ be set up on the exterior side of the cube,  at points
$R$, $S$, and $T$ (respectively), at right-angles to the planes of the cube.
 And let them be made equal to $RP$,
$PS$, and $TQ$. And let
$UB$, $BW$, $WC$, $CV$, and $VU$ be joined.

\epsfysize=3in
\centerline{\epsffile{Book13/fig17e.eps}}

I say that the pentagon $UBWCV$ is equilateral, and in one plane, and,
further, equiangular. For let $RB$, $SB$, and $VB$ be joined. 
And since the straight-line $NP$ has been cut in extreme and mean
ratio at $R$, and $RP$ is the greater piece,  the (sum of the squares)
on $PN$ and $NR$ is thus three times the (square) on $RP$ [Prop.~13.4].
And $PN$ (is) equal to $NB$, and $PR$ to $RU$. 
Thus, the (sum of the squares) on $BN$ and $NR$ is three times the
(square) on $RU$. And the (square) on $BR$ is equal to the (sum of the squares) on $BN$ and $NR$ [Prop.~1.47].  
Thus, the (square) on $BR$ is three times the (square) on $RU$. 
Hence, the (sum of the
squares) on $BR$ and $RU$ is four times the (square) on $RU$. And
the (square)
on $BU$ is equal to the (sum of the squares) on $BR$ and $RU$ [Prop.~1.47]. Thus, the (square) on $BU$ is four times the
(square) on $UR$.  Thus, $BU$ is double $RU$.  And
$VU$ is also double $UR$, inasmuch as $SR$ is also double $PR$---that is to say, $RU$. Thus, $BU$ (is) equal to $UV$. So, similarly, it can be shown
that  each of $BW$, $WC$, $CV$ is equal to each of $BU$ and $UV$. 
Thus, pentagon $BUVCW$ is equilateral. So, I say that it is also in one plane. For let $PX$ be drawn from $P$, parallel to
each of $RU$ and $SV$, on the exterior side of the cube. And let
$XH$ and $HW$ be joined. I say that $XHW$ is a straight-line.
For since $HQ$ has been cut in extreme and mean ratio at $T$,
and  $QT$  is its greater piece, thus as $HQ$ is to $QT$, so $QT$
(is) to $TH$. And $HQ$ (is) equal to $HP$, and $QT$ to each of
$TW$ and $PX$. Thus, as $HP$ is to $PX$, so $WT$ (is) to
$TH$. And $HP$ is parallel to $TW$. For of each of them is at right-angles
to the plane $BD$ [Prop.~11.6]. And $TH$ (is parallel) to
$PX$. For each of them is at right-angles to the plane $BF$ [Prop.~11.6].
And if two triangles, like
$XPH$ and $HTW$, having two sides proportional to two sides,
are placed together at a single angle  such that their corresponding sides are also parallel then the remaining sides will be straight-on (to one another) [Prop.~6.32]. Thus, $XH$ is straight-on to
$HW$. And every straight-line is in one plane [Prop.~11.1]. Thus,
pentagon $UBWCV$ is in one plane.

So, I say that it is also equiangular.

For since the straight-line $NP$ has been cut in extreme and mean ratio at
$R$, and $PR$ is the greater piece [thus as the sum of
$NP$ and $PR$ is to $PN$, so $NP$ (is) to $PR$], and $PR$
(is) equal to $PS$ [thus as $SN$ is to $NP$, so $NP$ (is) to $PS$], 
$NS$ has thus also been cut in extreme and mean ratio at $P$, and $NP$
is the greater piece [Prop.~13.5]. Thus, the (sum of the squares) on
$NS$ and $SP$ is three times the (square) on $NP$ [Prop.~13.4].
And $NP$ (is) equal to $NB$, and $PS$ to $SV$. Thus, the (sum of the)
squares on  $NS$ and $SV$ is three times the (square) on $NB$.
Hence, the (sum of the squares) on $VS$, $SN$, and $NB$
is four times the (square) on $NB$. And the (square) on $SB$
is equal to the (sum of the squares) on
$SN$ and $NB$  [Prop.~1.47].
Thus, the (sum of the squares) on $BS$ and $SV$---that is to say, 
the (square) on $BV$ [for angle $VSB$ (is) a right-angle]---is
four times the (square) on $NB$ [Def.~11.3, Prop.~1.47].
Thus, $VB$ is double $BN$. And $BC$ (is) also double $BN$.
Thus, $BV$ is equal to $BC$. And since the two (straight-lines)
$BU$ and $UV$ are equal to the two (straight-lines)
$BW$ and $WC$ (respectively), and the base $BV$ (is) equal to
the base $BC$, angle $BUV$ is thus equal to angle $BWC$ [Prop.~1.8].
So, similarly, we can show that angle $UVC$ is equal to angle
$BWC$.  Thus, the three angles
$BWC$, $BUV$, and $UVC$ are equal to one another.
And if three angles of an equilateral pentagon are equal to one another
then the pentagon is equiangular [Prop.~13.7].  Thus, pentagon
$BUVCW$ is equiangular. And it was also shown (to be)
equilateral. Thus, pentagon $BUVCW$ is equilateral and equiangular,
and it is on one of the sides, $BC$, of the cube. Thus, if we make the
same construction on each of the twelve sides of the cube then some solid
figure contained by twelve equilateral and equiangular pentagons 
will be constructed,  which is called a dodecahedron.

So, it is necessary to enclose it in the given sphere, and to show that the
side of the dodecahedron is that irrational (straight-line) called an apotome.

For let $XP$ be produced, and let (the produced straight-line) be $XZ$. Thus, $PZ$ meets
the diameter of the cube, and they cut one another in half.
For, this has been proved in the penultimate theorem of the eleventh book [Prop.~11.38]. Let them cut (one another) at $Z$. Thus, $Z$ is the
center of the sphere enclosing the cube, and $ZP$ (is) half the side of
the cube. So, let $UZ$ be joined. And since the straight-line 
$NS$ has been cut in extreme and mean ratio at $P$, and its greater piece
is $NP$, the (sum of the squares) on $NS$ and $SP$
is thus three times the (square) on $NP$ [Prop.~13.4]. And $NS$ (is)
equal to $XZ$, inasmuch as $NP$ is also equal to $PZ$, and
$XP$ to $PS$. But, indeed, $PS$ (is) also (equal) to $XU$,
since (it is) also (equal) to $RP$. Thus, the (sum of the squares) on 
$ZX$ and $XU$ is three times the (square) on $NP$.
And  the (square)
on $UZ$ is equal to the (sum of the squares) on $ZX$ and $XU$ [Prop.~1.47].  Thus, the (square) on $UZ$ is three times the
(square) on $NP$. 
 And the square on the
radius of the sphere enclosing the cube is also three times the (square) on half
the side of the cube. For it has previously been demonstrated (how to) construct 
the cube, and to enclose (it) in a sphere, and to show that the square on
the diameter of the sphere is three times the (square) on the
side of the cube [Prop.~13.15]. And if the (square on the) whole (is three
times) the (square on the) whole, then the (square on the) half (is) also (three times) the (square on the) half. 
And $NP$ is half of the side of the cube. Thus, $UZ$ is equal to
the radius of the sphere enclosing the cube. And $Z$ is the center of the
sphere enclosing the cube.
Thus, point $U$
is on the surface of the sphere. So, similarly, we can show that each of
the remaining angles of the dodecahedron is also  on the surface of the sphere. 
Thus, the dodecahedron has been enclosed by the given sphere.

So, I say that the side of the dodecahedron is that irrational straight-line
called an apotome.

For since  $RP$ is the greater piece of $NP$, which has been cut in extreme and mean ratio, and $PS$ is the greater piece of $PO$, which
has been cut in extreme and mean ratio,  $RS$ is thus the greater piece of the
whole of $NO$, which has been cut in extreme and mean ratio.
[Thus, since as $NP$ is to $PR$, (so) $PR$ (is) to $RN$, and  (the same is also true) of the doubles.
For parts have the same ratio as similar multiples (taken in corresponding
order) [Prop.~5.15]. Thus, as $NO$ (is) to $RS$, so $RS$
(is) to the sum of $NR$ and $SO$. And $NO$
(is) greater than $RS$. Thus, $RS$ (is) also
greater than the sum of $NR$ and $SO$ [Prop.~5.14]. Thus, $NO$
has been cut in extreme and mean ratio, and  $RS$ is its greater piece.]
And $RS$ (is) equal to $UV$. Thus, $UV$ is the greater piece of
$NO$, which has been cut in extreme and mean ratio. 
And since the diameter of the sphere is rational, and the square on it
is three times the (square) on the side of the cube,  $NO$, which is
the side of the cube, is thus rational. And if a rational (straight)-line
is cut in extreme and mean ratio then each of the
pieces is the irrational (straight-line called) an apotome.

Thus, $UV$, which is the side of the dodecahedron, is the
irrational (straight-line called) an apotome [Prop. 13.6].
}
\end{Parallel}

\begin{Parallel}{}{}
\ParallelLText{
\begin{center}
{\large {\greekfont Πόρισμα}.}
\end{center}\vspace*{-7pt}

{\greekfont Ἐκ δὴ τούτου φανερόν, ὅτι τῆς τοῦ κύβου πλευρᾶς
ἄκρον καὶ μέσον λόγον τεμνομένης τὸ μεῖζον
τμῆμά ἐστιν ἡ τοῦ δωδεκαέδρου πλευρά. ὅπερ
ἔδει δεῖχαι.}}

\ParallelRText{
\begin{center}
{\large Corollary}
\end{center}

So, (it is) clear, from this, that the side of the dodecahedron
is the greater piece of the side of the cube, when it is cut in extreme
and mean ratio.$^\dag$ (Which is) the very thing it was required to show.}
\end{Parallel}
{\footnotesize\noindent$^\dag$ If the radius of the circumscribed sphere is
unity then the side of the cube is $\sqrt{4/3}$, and the side of the
dodecahedron is $(1/3)\,(\sqrt{15}-\sqrt{3})$.}

%%%%
%13.18
%%%%
\pdfbookmark[1]{Proposition 13.18}{pdf13.18}
\begin{Parallel}{}{}
\ParallelLText{
\begin{center}
{\large \ggn{18}.}
\end{center}\vspace*{-7pt}

{\greekfont Τὰς πλευρὰς τῶν πέντε σθημάτων ἐκθέσθαι καὶ
συγκρῖν-αι πρὸς ἀλλήλας.}

\epsfysize=2.25in
\centerline{\epsffile{Book13/fig18g.eps}}

{\greekfont Ἐκκείσθω ἡ τῆς δοθείσης σφαίρας διάμετρος ἡ ΑΒ, καὶ
τετμήσθω κατὰ τὸ Γ ὥστε ἴσην εἶναι τὴν ΑΓ τῇ ΓΒ, κατὰ
δὲ τὸ Δ ὥστε διπλασίονα εἶναι τὴν ΑΔ τῆς ΔΒ, καὶ
γεγράφθω ἐπὶ τῆς ΑΒ ἡμικύκλιον τὸ ΑΕΒ, καὶ ἀπὸ τῶν
Γ, Δ τῇ ΑΒ πρὸς ὀρθὰς ἤθθωσαν αἱ ΓΕ, ΔΖ, καὶ ἐπεζεύθθωσαν
αἱ ΑΖ, ΖΒ, ΕΒ. καὶ ἐπεὶ διπλῆ ἐστιν ἡ ΑΔ τῆς ΔΒ,
τριπλῆ ἄρα ἐστὶν ἡ ΑΒ τῆς ΒΔ. ἀναστρέψαντι ἡμιολία
ἄρα ἐστὶν ἡ ΒΑ τῆς ΑΔ. ὡς δὲ ἡ ΒΑ πρὸς τὴν ΑΔ,
οὕτως τὸ ἀπὸ τῆς ΒΑ πρὸς τὸ ἀπὸ τῆς ΑΖ;
ἰσογώνιον γάρ ἐστι τὸ ΑΖΒ τρίγωνον τῷ ΑΖΔ
τριγώνῳ; ἡμιόλιον ἄρα ἐστὶ τὸ ἀπὸ τῆς ΒΑ τοῦ ἀπὸ
τῆς ΑΖ. ἔστι δὲ καὶ ἡ τῆς σφαίρας διάμετρος δυνάμει
ἡμιολία τῆς πλευρᾶς τῆς πυραμίδος. καί ἐστιν ἡ ΑΒ
ἡ τῆς σφαίρας διάμετρος; ἡ ΑΖ ἄρα ἴση ἐστὶ τῇ
πλευρᾷ τῆς πυραμίδος.}

{\greekfont Πάλιν, ἐπεὶ διπλασίων ἐστὶν ἡ ΑΔ τῆς ΔΒ, τριπλῆ
ἄρα ἐστὶν ἡ ΑΒ τῆς ΒΔ. ὡς δὲ ἡ ΑΒ πρὸς τὴν ΒΔ,
οὕτως τὸ ἀπὸ τῆς ΑΒ πρὸς τὸ ἀπὸ τῆς ΒΖ; τριπλάσιον
ἄρα ἐστὶ τὸ ἀπὸ τῆς ΑΒ τοῦ ἀπὸ τῆς ΒΖ. ἔστι
δὲ καὶ ἡ τῆς σφαίρας διάμετρος δυνάμει τριπλασίων
τῆς τοῦ κύβου πλευρᾶς. καί ἐστιν ἡ ΑΒ ἡ τῆς σφαίρας
διάμετρος; ἡ ΒΖ ἄρα τοῦ κύβου ἐστὶ πλευρά.}

{\greekfont Καὶ ἐπεὶ ἴση ἐστὶν ἡ ΑΓ τῇ ΓΒ, διπλῆ ἄρα
ἐστὶν ἡ ΑΒ τῆς ΒΓ. ὡς δὲ ἡ ΑΒ πρὸς τὴν ΒΓ, οὕτως
τὸ ἀπὸ τῆς ΑΒ πρὸς τὸ ἀπὸ τῆς ΒΕ; διπλάσιον
ἄρα ἐστὶ τὸ ἀπὸ τῆς ΑΒ τοῦ ἀπὸ τῆς ΒΕ. ἔστι
δὲ καὶ ἡ τῆς σφαίρας διάμετρος δυνάμει διπλασίων
τῆς τοῦ ὀκταέδρου πλευρᾶς. καὶ ἐστιν ἡ ΑΒ ἡ
τῆς δοθείσης σφαίρας διάμετρος; ἡ ΒΕ ἄρα τοῦ ὀκταέδρου
ἐστὶ πλευρά.}

{\greekfont Ἤθθω δὴ ἀπὸ τοῦ Α σημείου τῇ ΑΒ εὐθείᾳ
πρὸς ὀρθὰς ἡ ΑΗ, καὶ κείσθω ἡ ΑΗ ἴση τῇ ΑΒ,
καὶ ἐπεζεύθθω ἡ ΗΓ, καὶ ἀπὸ τοῦ Θ ἐπὶ
τὴν ΑΒ κάθετος ἤθθω ἡ ΘΚ. καὶ ἐπεὶ διπλῆ ἐστιν
ἡ ΗΑ τῆς ΑΓ; ἴση γὰρ ἡ ΗΑ τῇ ΑΒ; ὡς δὲ ἡ ΗΑ
πρὸς τὴν ΑΓ, οὕτως ἡ ΘΚ πρὸς τὴν ΚΓ, διπλῆ ἄρα
καὶ ἡ ΘΚ τῆς ΚΓ. τετραπλάσιον ἄρα ἐστὶ τὸ ἀπὸ
τῆς ΘΚ τοῦ ἀπὸ τῆς ΚΓ; τὰ ἄρα ἀπὸ τῶν ΘΚ, ΚΓ,
ὅπερ ἐστὶ τὸ ἀπὸ τῆς ΘΓ, πενταπλάσιόν  ἐστι τοῦ
ἀπὸ τῆς ΚΓ. ἴση δὲ ἡ ΘΓ τῇ ΓΒ; πενταπλάσιον
ἄρα ἐστὶ τὸ ἀπὸ τῆς ΒΓ τοῦ ἀπὸ τῆς ΓΚ. καὶ
ἐπεὶ διπλῆ ἐστιν ἡ ΑΒ τῆς ΓΒ, ὧν ἡ ΑΔ τῆς
ΔΒ ἐστι διπλῆ,  λοιπὴ ἄρα ἡ ΒΔ λοιπῆς τῆς ΔΓ ἐστι
διπλῆ. τριπλῆ ἄρα ἡ ΒΓ τῆς ΓΔ; ἐνναπλάσιον ἄρα τὸ
ἀπὸ τῆς ΒΓ τοῦ ἀπὸ τῆς ΓΔ. πενταπλάσιον δὲ τὸ ἀπὸ
τῆς ΒΓ τοῦ ἀπὸ τῆς ΓΚ; μεῖζον ἄρα τὸ ἀπὸ τῆς
ΓΚ τοῦ ἀπὸ τῆς ΓΔ. μείζων ἄρα ἐστὶν ἡ ΓΚ τῆς
ΓΔ. κείσθω τῇ ΓΚ ἴση ἡ ΓΛ, καὶ ἀπὸ τοῦ Λ τῇ
ΑΒ πρὸς ὀρθὰς ἤθθω ἡ ΛΜ, καὶ ἐπεζεύθθω ἡ ΜΒ.
καὶ ἐπεὶ πενταπλάσιόν ἐστι τὸ ἀπὸ  τῆς ΒΓ τοῦ ἀπὸ τῆς
ΓΚ, καί ἐστι τῆς μὲν ΒΓ διπλῆ ἡ ΑΒ, τῆς δὲ ΓΚ διπλῆ
ἡ ΚΛ, πενταπλάσιον ἄρα ἐστὶ τὸ ἀπὸ τῆς ΑΒ τοῦ ἀπὸ τῆς
ΚΛ.
ἔστι δὲ καὶ ἡ τῆς σφαίρας διάμετρος δυνάμει πενταπλασίων
τῆς ἐκ τοῦ κέντρου τοῦ κύκλου, ἀφ οὗ τὸ εἰκοσάεδρον
ἀναγέγραπται. καί ἐστιν ἡ ΑΒ ἡ τῆς σφαίρας διάμετρος;
ἡ ΚΛ ἄρα ἐκ τοῦ κέντρου ἐστὶ τοῦ κύκλου, ἀφ οὗ
τὸ εἰκοσάεδρον ἀναγέγραπται; ἡ ΚΛ ἄρα ἑχαγώνου ἐστὶ
πλευρὰ τοῦ εἰρημένου κύκλου. καὶ ἐπεὶ ἡ τῆς σφαίρας
διάμετρος σύγκειται ἔκ τε τῆς τοῦ ἑχαγώνου καὶ δύο τῶν
τοῦ δεκαγώνου τῶν εἰς τὸν εἰρημένον κύκλον ἐγγραφομένων,
καί ἐστιν ἡ μὲν ΑΒ ἡ τῆς σφαίρας διάμετρος, ἡ δὲ ΚΛ
ἑχαγώνου πλευρά, καὶ ἴση ἡ ΑΚ τῇ ΛΒ, ἑκατέρα ἄρα τῶν
ΑΚ, ΛΒ δεκαγώνου ἐστὶ πλευρὰ τοῦ ἐγγραφομένου εἰς
τὸν κύκλον, ἀφ οὗ τὸ εἰκοσάεδρον ἀναγέγραπται. καὶ
ἐπεὶ δεκαγώνου μὲν ἡ ΛΒ, ἑχαγώνου δὲ ἡ ΜΛ;
ἴση γάρ ἐστι τῇ ΚΛ, ἐπεὶ καὶ τῇ ΘΚ; ἴσον γὰρ ἀπέθουσιν
ἀπὸ τοῦ κέντρου; καί ἐστιν ἑκατέρα τῶν ΘΚ, ΚΛ διπλασίων
τῆς ΚΓ; πενταγώνου ἄρα ἐστὶν ἡ ΜΒ. ἡ δὲ τοῦ πενταγώνου
ἐστὶν ἡ τοῦ εἰκοσαέδρου; εἰκοσαέδρου ἄρα ἐστὶν ἡ ΜΒ.}

{\greekfont Καὶ ἐπεὶ ἡ ΖΒ κύβου ἐστὶ πλευρά, τετμήσθω ἄκρον καὶ
μέσον λόγον κατὰ τὸ Ν, καὶ ἔστω μεῖζον τμῆμα τὸ
ΝΒ; ἡ ΝΒ ἄρα δωδεκαέδρου ἐστὶ πλευρά.}

{\greekfont Καὶ ἐπεὶ ἡ τῆς σφαίρας διάμετρος ἐδείθθη τῆς μὲν
ΑΖ πλευρᾶς τῆς πυραμίδος δυνάμει ἡμιολία, τῆς δὲ
τοῦ ὀκταέδρου τῆς ΒΕ δυνάμει διπλασίων, τῆς δὲ τοῦ
κύβου τῆς ΖΒ δυνάμει τριπλασίων, οἵων ἄρα ἡ τῆς
σφαίρας διάμετρος δυνάμει ἕχ, τοιούτων ἡ μὲν τῆς
πυραμίδος τεσσάρων, ἡ δὲ τοῦ ὀκταέδρου τριῶν, ἡ δὲ
τοῦ κύβου δύο.  ἡ μὲν ἄρα τῆς πυραμίδος πλευρὰ
τῆς μὲν τοῦ ὀκταέδρου πλευρᾶς δυνάμει ἐστὶν ἐπίτριτος,
τῆς δὲ τοῦ κύβου δυνάμει διπλῆ, ἡ δὲ τοῦ ὀκταέδρου τῆς
τοῦ κύβου δυνάμει ἡμιολία. αἱ μὲν οὖν εἰρημέναι
τῶν τριῶν σθημάτων πλευραί, λέγω δὴ πυραμίδος καὶ ὀκταέδρου
καὶ κύβου, πρὸς ἀλλήλας εἰσὶν ἐν λόγοις ῥητοῖς. αἱ δὲ λοιπαὶ
δύο, λέγω δὴ ἥ τε τοῦ εἰκοσαέδρου καὶ ἡ τοῦ δωδεκαέδρου,
οὔτε πρὸς ἀλλήλας οὔτε πρὸς τὰς προειρημένας εἰσὶν ἐν λόγοις
ῥητοῖς; ἄλογοι γάρ εἰσιν, ἡ μὲν ἐλάττων, ἡ δὲ ἀποτομή.}

{\greekfont Ὅτι μείζων ἐστὶν ἡ τοῦ εἰκοσαέδρου πλευρὰ ἡ ΜΒ
τῆς τοῦ δωδεκαέδρου τῆς ΝΒ, δείχομεν οὕτως.}

{\greekfont Ἐπεὶ γὰρ ἰσογώνιόν ἐστι τὸ ΖΔΒ τρίγωνον τῷ
ΖΑΒ τριγώνῳ, ἀνάλογόν ἐστιν ὡς ἡ ΔΒ πρὸς
τὴν ΒΖ, οὕτως ἡ ΒΖ πρὸς τὴν ΒΑ. καὶ ἐπεὶ τρεῖς
εὐθεῖαι ἀνάλογόν εἰσιν, ἔστιν ὡς ἡ πρώτη πρὸς
τὴν τρίτην, οὕτως τὸ ἀπὸ τῆς πρώτης πρὸς τὸ ἀπὸ
τῆς δευτέρας;  ἔστιν ἄρα ὡς ἡ ΔΒ πρὸς τὴν ΒΑ, οὕτως
τὸ ἀπὸ τῆς ΔΒ πρὸς τὸ ἀπὸ τῆς ΒΖ; ἀνάπαλιν ἄρα ὡς ἡ ΑΒ
πρὸς τὴν ΒΔ, οὕτως τὸ ἀπὸ τῆς ΖΒ πρὸς τὸ ἀπὸ
τῆς ΒΔ. τριπλῆ δὲ ἡ ΑΒ τῆς ΒΔ; τριπλάσιον
ἄρα τὸ ἀπὸ τῆς ΖΒ τοῦ ἀπὸ τῆς ΒΔ. ἔστι
δὲ καὶ τὸ ἀπὸ τῆς ΑΔ τοῦ ἀπὸ τῆς ΔΒ τετραπλάσιον;
διπλῆ γὰρ ἡ ΑΔ τῆς ΔΒ; μεῖζον ἄρα τὸ ἀπὸ τῆς
ΑΔ τοῦ ἀπὸ τῆς ΖΒ; μείζων ἄρα ἡ ΑΔ τῆς ΖΒ;
πολλῷ ἄρα ἡ ΑΛ τῆς ΖΒ μείζων ἐστίν. καὶ τῆς
μὲν ΑΛ ἄκρον καὶ μέσον λόγον τεμνομένης τὸ μεῖζον
τμῆμά ἐστιν ἡ ΚΛ, ἐπειδήπερ ἡ μὲν ΛΚ ἑχαγώνου ἐστίν,
ἡ δὲ ΚΑ δεκαγώνου; τῆς δὲ ΖΒ ἄκρον καὶ μέσον
λόγον τεμνομένης τὸ μεῖζον τμῆμά ἐστιν ἡ ΝΒ; μείζων
ἄρα ἡ ΚΛ τῆς ΝΒ. ἴση δὲ ἡ ΚΛ τῇ ΛΜ; μείζων
ἄρα ἡ ΛΜ τῆς ΝΒ [τῆς δὲ ΛΜ μείζων ἐστὶν ἡ ΜΒ].
πολλῷ ἄρα ἡ ΜΒ πλευρὰ οὖσα τοῦ εἰκοσαέδρου μείζων
ἐστὶ τῆς ΝΒ πλευρᾶς οὔσης τοῦ δωδεκαέδρου; ὅπερ
ἔδει δεῖχαι.}}

\ParallelRText{
\begin{center}
{\large Proposition 18}
\end{center}

To set out the sides of the five (aforementioned) figures, and to compare (them)
with one another.$^\dag$

\epsfysize=2.25in
\centerline{\epsffile{Book13/fig18e.eps}}

Let the diameter, $AB$, of the given sphere be laid out. And let it
be cut at $C$, such that $AC$ is equal to $CB$, and at $D$, such
that $AD$ is double $DB$. And let the semi-circle $AEB$ be
drawn on $AB$. And let $CE$ and $DF$ be drawn from $C$
and $D$ (respectively), at right-angles to $AB$. And let $AF$, $FB$,
and $EB$ be joined. And since $AD$ is double $DB$, 
$AB$ is thus triple $BD$. Thus, via conversion, $BA$ is one and
a half times $AD$. And as $BA$ (is) to $AD$, so the (square)
on $BA$ (is) to the (square) on $AF$ [Def.~5.9].  For
triangle $AFB$ is equiangular to triangle $AFD$ [Prop.~6.8]. 
Thus, the (square) on $BA$ is one and a half times the (square)
on $AF$. And the square on the diameter of the sphere
is also one and a half times the (square) on the side of the pyramid
[Prop.~13.13].  And $AB$ is the diameter of the sphere. Thus,
$AF$ is equal to the side of the pyramid.

Again, since $AD$ is double $DB$,  $AB$ is thus triple $BD$. And
as $AB$ (is) to $BD$, so the (square) on $AB$ (is) to the (square)
on $BF$ [Prop.~6.8, Def.~5.9]. Thus, the (square) on $AB$
is three times the (square) on $BF$.  And the square on the diameter of
the sphere is also three times the (square) on the side of the cube [Prop.~13.15]. And $AB$ is the diameter of the sphere. Thus, $BF$ is the
side of the cube.

And since $AC$ is equal to $CB$, $AB$ is thus double $BC$.  And
as $AB$ (is) to $BC$, so the (square) on $AB$ (is) to the (square)
on $BE$ [Prop.~6.8, Def.~5.9]. Thus, the (square) on $AB$
is double the (square) on $BE$. And the square on the diameter
of the sphere is also double the (square) on the side of the octagon [Prop.~13.14]. And $AB$ is the diameter of the given sphere. Thus,
$BE$ is the side of the octagon.

So let $AG$ be drawn from point $A$ at right-angles to the straight-line
$AB$. And let $AG$ be made equal to $AB$. And let $GC$
be joined. And let $HK$ be drawn from $H$, perpendicular to $AB$. And since $GA$ is double $AC$. 
For $GA$ (is) equal to $AB$.  And as $GA$ (is) to $AC$, so
$HK$ (is) to $KC$ [Prop.~6.4]. $HK$ (is) thus also double
$KC$. Thus, the (square) on $HK$ is four times the (square)
on $KC$. Thus, the (sum of the squares) on $HK$ and $KC$,
which is the (square) on $HC$ [Prop.~1.47],  is five times the
(square) on $KC$.  And $HC$ (is) equal to $CB$. Thus,
the (square) on $BC$ (is) five times the (square) on  $CK$. 
And since $AB$ is double $CB$, of which $AD$ is double $DB$, 
the remainder $BD$ is thus double the remainder $DC$.
$BC$ (is) thus triple $CD$. The (square) on $BC$ (is) thus nine times 
the (square) on $CD$.  And the (square) on $BC$ (is) five times the (square)
on $CK$.  Thus, the (square) on $CK$ (is) greater than the (square) on $CD$.
$CK$ is thus greater than $CD$. Let $CL$ be made equal to $CK$.
And let $LM$ be drawn from $L$ at right-angles to $AB$. 
And let $MB$ be joined. And since the (square) on  $BC$
is five times the (square) on $CK$, and $AB$ is double $BC$, and $KL$
double $CK$, the (square) on $AB$ is thus five times the (square)
on $KL$. And the square on the diameter of the sphere  is also five times
the (square) on the radius of the circle from which the icosahedron has been
described [Prop.~13.16~corr.]. And $AB$ is the diameter of the sphere.
Thus, $KL$ is the radius of the circle from which the icosahedron has
been described. Thus, $KL$ is (the side) of the hexagon (inscribed) in the aforementioned
circle [Prop.~4.15~corr.]. And since the diameter of the sphere is composed of (the side) of the hexagon, and two of (the sides) of the decagon, inscribed
in the aforementioned circle, and $AB$ is the diameter of the sphere, 
and $KL$ the side of the hexagon, and $AK$ (is) equal to $LB$, thus
$AK$ and $LB$ are each sides of the decagon inscribed in the  circle
from which the icosahedron has been described. And since
$LB$ is (the side) of the decagon. And $ML$ (is the side) of the hexagon---for (it is) equal to $KL$, since (it is) also (equal) to $HK$, for they are equally
far from the center. And $HK$ and $KL$ are each double $KC$. 
$MB$ is thus (the side) of the pentagon (inscribed in the circle) [Props.~13.10, 1.47]. 
And (the side) of the pentagon is (the side) of the icosahedron [Prop.~13.16]. 
Thus, $MB$ is (the side) of the icosahedron. 

And since $FB$ is the side of the cube, let it be cut in extreme and
mean ratio at $N$, and let $NB$ be the greater piece. Thus, $NB$
is the side of the dodecahedron [Prop.~13.17~corr.].

And since the (square) on the diameter of the sphere was shown (to be)
one and a half times the square on the side, $AF$,  of the pyramid,
and twice the square on (the side), $BE$,  of the octagon, and
three times the square on (the side), $FB$, of the cube, 
thus, of whatever (parts) the (square) on the diameter of the
sphere (makes) six, of such (parts) the (square) on (the side) of the pyramid (makes)
four, and (the square) on (the side) of the octagon three, and (the
square) on (the side) of the cube  two. Thus, the (square) on the
side of the pyramid is one and a third times the square on the side of the
octagon, and double the square on (the side) of the cube. 
And the (square) on (the side) of the octahedron is one and a half
times the square on (the side) of the cube. Therefore, the aforementioned
sides of the three figures---I mean, of the pyramid, and of the octahedron, and 
of the cube---are
in rational ratios to one another. And (the sides of) the remaining two (figures)---I mean, of
the icosahedron, and of the dodecahedron---are neither in rational ratios to
one another, nor to the (sides) of the aforementioned (three figures). For they are
irrational (straight-lines): (namely), a minor [Prop.~13.16], and 
an apotome [Prop.~13.17].

(And), we can show that the side, $MB$, of the icosahedron is greater
that the (side), $NB$, or the dodecahedron, as follows.

For, since triangle $FDB$ is equiangular to triangle $FAB$ [Prop.~6.8],
proportionally, as $DB$ is to $BF$, so $BF$ (is) to $BA$ [Prop.~6.4]. 
And since three straight-lines are (continually) proportional, as the
first (is) to the third, so the (square) on the first (is) to the (square) on the
second [Def.~5.9, Prop.~6.20~corr.]. Thus, as $DB$ is to $BA$, so the (square) on $DB$
(is) to the (square) on $BF$. Thus, inversely, as $AB$ (is) to $BD$, so the
(square) on $FB$ (is) to the (square) on $BD$. And $AB$ (is) triple
$BD$. Thus, the (square) on $FB$ (is) three times the (square)
on $BD$. And the (square) on $AD$ is also four times the (square)
on $DB$. For $AD$ (is) double $DB$. Thus, the (square) on $AD$
(is) greater than the (square) on $FB$.  Thus, $AD$ (is) greater than
$FB$. 
Thus, $AL$ is much greater than $FB$. And $KL$ is the greater piece of $AL$, which is cut in extreme and
mean ratio---inasmuch as $LK$ is (the side) of the hexagon, and $KA$
(the side) of the decagon [Prop.~13.9]. And $NB$
is the greater piece of $FB$, which is cut in extreme and mean ratio. 
Thus, $KL$ (is) greater than $NB$.  And $KL$ (is) equal to $LM$. 
Thus, $LM$ (is) greater than $NB$ [and $MB$ is greater than $LM$].
Thus, $MB$, which is (the side) of the icosahedron, is much greater
than $NB$, which is (the side) of the dodecahedron.  (Which is)
the very thing it was required to show.}
\end{Parallel}
{\footnotesize\noindent$\dag$ If the radius of the given sphere is unity 
then the sides of the pyramid ({\em i.e.}, tetrahedron), octahedron, cube,
icosahedron, and dodecahedron, respectively, satisfy the following
inequality: $\sqrt{8/3} > \sqrt{2}>\sqrt{4/3}>(1/\sqrt{5})\,\sqrt{10 -2\,\sqrt{5}}
> (1/3)\,(\sqrt{15}-\sqrt{3})$.}

\begin{Parallel}{}{}
\ParallelLText{~\\

{\greekfont Λέγω δή, ὅτι παρὰ τὰ εἰρημένα πέντε σθήματα
οὐ συσταθήσεται ἕτερον σχῆμα περιεχόμενον ὑπὸ
ἰσοπλεύρων τε καὶ ἰσογωνίων ἴσων ἀλλήλοις.}

{\greekfont Ὑπὸ μὲν γὰρ δύο τριγώνων ἢ ὅλως ἐπιπέδων
στερεὰ γωνία οὐ συνίσταται. ὑπὸ δὲ τριῶν τριγώνων
ἡ τῆς πυραμίδος, ὑπὸ δὲ τεσσάρων ἡ τοῦ ὀκταέδρου,
ὑπὸ δὲ πέντε ἡ τοῦ εἰκοσαέδρου; ὑπὸ δὲ ἓχ τριγώνων
ἰσοπλεύρων τε καὶ ἰσογωνίων πρὸς ἑνὶ σημείῳ
συνισταμένων οὐκ ἔσται στερεὰ γωνία; οὔσης γὰρ τῆς
τοῦ ἰσοπλεύρου τριγώνου γωνίας διμοίρου ὀρθῆς
ἔσονται αἱ ἓχ τέσσαρσιν ὀρθαῖς ἴσαι; ὅπερ ἀδύνατον;
ἅπασα γὰρ στερεὰ γωνία ὑπὸ ἐλασσόνων ἢ τεσσάρων
ὀρθῶν περέθεται. διὰ τὰ α ὐτὰ δὴ οὐδὲ ὑπὸ
πλειόνων ἢ ἓχ γωνιῶν ἐπιπέδων στερεὰ γωνία
συνίσταται. ὑπὸ δὲ τετραγώνων τριῶν ἡ τοῦ κύβου
γωνία περιέθεται; ὑπὸ δὲ τεσσάρων ἀδύνατον; ἔσονται
γὰρ πάλιν τέσσαρες ὀρθαί. ὑπὸ δὲ πενταγώνων ἰσοπλεύρων
καὶ ἰσογωνίων, ὑπὸ μὲν τριῶν ἡ τοῦ δωδεκαέδρου;
ὑπὸ δὲ τεσσάρων ἀδύνατον; οὔσης γὰρ τῆς τοῦ
πενταγώνου ἰσοπλεύρου γωνίας ὀρθῆς καὶ πέμπτου,
ἔσονται αἱ τέσσαρες γωνίαι τεσσάρων ὀρθῶν μείζους;
ὅπερ ἀδύνατον. οὐδὲ μὴν ὑπὸ πολυγώνων ἑτέρων
σθημάτων περισθεθήσεται στερεὰ γωνία διὰ τὸ α ὐτὸ
ἄτοπον.}

{\greekfont Οὐκ ἄρα παρὰ τὰ εἰρημένα πέντε σθήματα ἕτερον
σχῆμα στερεὸν συσταθήσεται ὑπὸ ἰσοπλεύρων
τε καὶ ἰσογωνίων περιεχόμενον; ὅπερ
ἔδει δεῖχαι.}}

\ParallelRText{~\\

So, I say that, beside the five aforementioned  figures,  no other (solid) figure
can be constructed (which is) contained by equilateral and equiangular (planes), equal to one another.

For a solid angle cannot be constructed from two triangles, or indeed (two) planes (of any sort) [Def.~11.11]. And (the solid angle) of the
pyramid (is constructed) from three (equiangular) triangles, and (that) of the octahedron
from four (triangles), and (that) of the icosahedron from (five) triangles. 
And a solid angle cannot be (made) from six equilateral and
equiangular triangles set up together at one point. For, since the angles of a
equilateral triangle are (each)  two-thirds of a right-angle, the (sum of the)
six  (plane) angles (containing the solid angle) will be four right-angles. The very thing (is) impossible.
For every solid angle is contained by (plane angles whose sum is) less than four right-angles [Prop.~11.21].  So, for the same (reasons), a solid angle cannot be constructed from more than six  plane angles (equal to two-thirds of a right-angle) either. And the (solid) angle of a cube
is contained by three squares. And (a solid angle contained) by four
(squares is) impossible. For, again, the (sum of the plane angles
containing the solid angle) will be four right-angles. And (the solid angle)
of a dodecahedron (is contained) by three equilateral and equiangular
pentagons. And (a solid angle contained) by four (equiangular
pentagons is) impossible. For, the angle of an equilateral
pentagon being one and one-fifth of right-angle, four (such) angles will be
greater (in sum) than four right-angles. The very thing (is) impossible. And, on account of the same absurdity, a solid angle cannot be constructed from any other
(equiangular) polygonal figures either.

Thus, beside the five aforementioned figures, no other solid figure can be constructed (which is) contained by equilateral and equiangular (planes).
(Which is) the very thing it was required to show.}
\end{Parallel}

\begin{Parallel}{}{}
\ParallelLText{

\epsfysize=2in
\centerline{\epsffile{Book13/fig18ag.eps}}

\begin{center}\vspace*{-7pt}
{\large {\greekfont Λῆμμα}.}
\end{center}\vspace*{-7pt}

{\greekfont Ὅτι δὲ ἡ τοῦ ἰσοπλεύρου καὶ ἰσογωνίου πενταγώνου
γωνία ὀρθή ἐστι καὶ πέμπτου, οὕτω δεικτέον.}

{\greekfont Ἔστω γὰρ πεντάγωνον ἰσόπλευρον καὶ ἰσογώνιον
τὸ ΑΒΓΔΕ, καὶ περιγεγράφθω περὶ α ὐτὸ κύκλος ὁ ΑΒΓΔΕ,
καὶ εἰλήφθω α ὐτοῦ τὸ κέντρον τὸ Ζ, καὶ ἐπεζεύθθωσαν
αἱ ΖΑ, ΖΒ, ΖΓ, ΖΔ, ΖΕ. δίθα ἄρα τέμνουσι τὰς πρὸς τοῖς
Α, Β, Γ, Δ, Ε τοῦ πενταγώνου γωνίας. καὶ ἐπεὶ αἱ πρὸς
τῷ Ζ πέντε γωνίαι τέσσαρσιν ὀρθαῖς ἴσαι εἰσὶ καί
εἰσιν ἴσαι, μία ἄρα α ὐτῶν, ὡς ἡ ὑπὸ ΑΖΒ, μιᾶς
ὀρθῆς ἐστι παρὰ πέμπτον; λοιπαὶ ἄρα αἱ ὑπὸ ΖΑΒ, ΑΒΖ
μιᾶς εἰσιν ὀρθῆς καὶ πέμπτου. ἴση δὲ ἡ ὑπὸ ΖΑΒ τῇ
ὑπὸ ΖΒΓ; καὶ ὅλη ἄρα ἡ ὑπὸ ΑΒΓ τοῦ πενταγώνου
γωνία μιᾶς
ἐστιν  ὀρθῆς καὶ πέμπτου; ὅπερ ἔδει δεῖχαι.}}

\ParallelRText{

\epsfysize=2in
\centerline{\epsffile{Book13/fig18ae.eps}}

\begin{center}\vspace*{-7pt}
{\large Lemma}
\end{center}\vspace*{-7pt}

It can be shown that the angle of an equilateral and
equiangular pentagon is one and one-fifth of a right-angle,
as follows.

For let $ABCDE$ be an equilateral and equiangular pentagon, and
let the circle $ABCDE$ be circumscribed about it [Prop.~4.14]. 
And let its center, $F$, be found [Prop.~3.1]. 
And let $FA$, $FB$, $FC$, $FD$, and $FE$ be joined. Thus,
they cut the angles of the pentagon in half at (points)
 $A$, $B$, $C$, $D$, and $E$  [Prop.~1.4].
 And since the five
angles at $F$ are equal (in sum) to four right-angles, and
are also equal (to one another),  (any) one of them, like $AFB$, is thus
one less a fifth of a right-angle. Thus, the (sum of the) remaining
(angles in triangle $ABF$), $FAB$ and $ABF$,  is one plus a fifth  of a right-angle
[Prop.~1.32]. And $FAB$ (is) equal to $FBC$. Thus, the
whole angle, $ABC$, of the pentagon is also one and one-fifth
of a right-angle. (Which is) the very thing it was required to show.}
\end{Parallel}
